\documentclass[msc]{cls_bst/cls_file_ufam_v3}

\usepackage{amsmath}
\usepackage{amssymb}
\usepackage{amsfonts}

\usepackage[dvipsnames]{xcolor}
%% Paleta neutra (nude/taupe) do documento. Usar estas cores nas figuras de
%% todos os capitulos para manter consistencia visual.
\definecolor{palPaper}{HTML}{F4EFE8}
\definecolor{palFill}{HTML}{EAE1D2}
\definecolor{palStrong}{HTML}{D9C4A0}
\definecolor{palEdge}{HTML}{A4927A}
\definecolor{palEdgeDark}{HTML}{6F6051}
\usepackage{tikz}
\usetikzlibrary{positioning,shapes.geometric,arrows.meta,calc,fit,backgrounds}
\usepackage{listings}
\lstset{inputencoding=utf8,extendedchars=true,literate=
  {á}{{\'a}}1 {à}{{`a}}1 {â}{{\^a}}1 {ã}{{\~a}}1 {é}{{\'e}}1 {ê}{{\^e}}1
  {í}{{\'i}}1 {ó}{{\'o}}1 {ô}{{\^o}}1 {õ}{{\~o}}1 {ú}{{\'u}}1 {ç}{{\c c}}1
  {Á}{{\'A}}1 {À}{{`A}}1 {Â}{{\^A}}1 {Ã}{{\~A}}1 {É}{{\'E}}1 {Ê}{{\^E}}1
  {Í}{{\'I}}1 {Ó}{{\'O}}1 {Ô}{{\^O}}1 {Õ}{{\~O}}1 {Ú}{{\'U}}1 {Ç}{{\c C}}1
  {°}{{\textdegree}}1}

\usepackage{algorithm}
\usepackage{algpseudocode}

\usepackage{tcolorbox}

\usepackage{enumitem}

\usepackage{pdfpages}

\usepackage[english,brazil]{babel}
\usepackage[utf8]{inputenc}
\usepackage[T1]{fontenc}
\usepackage{mathptmx}
\usepackage{indentfirst}
\usepackage{fancyhdr}
\usepackage{url}
\usepackage{cite}
\usepackage{float}
\usepackage{multirow}
\usepackage{booktabs}

\fancypagestyle{fancy}{%
  \fancyhf{}%
  \renewcommand{\headrulewidth}{0.8pt}%
  \renewcommand{\footrulewidth}{0pt}%
  \fancyhead[L]{\footnotesize\scshape\nouppercase{\leftmark}}%
  \fancyhead[R]{\footnotesize\scshape\nouppercase{\rightmark}}%
  \fancyfoot[C]{\small\thepage}%
}
\fancypagestyle{plain}{%
  \fancyhf{}%
  \renewcommand{\headrulewidth}{0pt}%
  \fancyfoot[C]{\small\thepage}%
}
\renewcommand{\chaptermark}[1]{\markboth{Capítulo \thechapter\ -- #1}{}}
\renewcommand{\sectionmark}[1]{\markright{\thesection\ \ #1}}

\newif\ifbanca
\bancafalse

%% ┌─────────────────────────────────────────────────────────────────────┐
%% │ ANTES DE VIRAR PARA \bancatrue, RESOLVER:                           │
%% │                                                                     │
%% │ A única pendência vermelha do documento está no Apêndice            │
%% │ \ref{apend:contraexemplos}: os contraexemplos NÃO são cópias        │
%% │ literais dos logs. Ocultar o \pendente sem corrigir o conteúdo      │
%% │ esconde o problema do leitor, não do examinador. — LC               │
%% │                                                                     │
%% │ Ação: substituir cada contraexemplo pelo trecho literal do log do   │
%% │ servidor iartes. Valor canônico do candidato 1 é                    │
%% │ _next_buffer_idx = 74 (o 75 vem da reprodução em macOS e não entra  │
%% │ na dissertação).                                                    │
%% └─────────────────────────────────────────────────────────────────────┘

\newcommand{\revisao}[1]{#1}
\newcommand{\corr}[1]{#1}

\newcommand{\artigo}[1]{#1}
\newcommand{\novo}[1]{#1}

\ifbanca
  \newcommand{\pendente}[1]{#1}
  \newcommand{\falta}[1]{#1}
\else
  \newcommand{\pendente}[1]{\textcolor{orange!85!black}{#1}}
  \newcommand{\falta}[1]{\textcolor{red!80!black}{#1}}
\fi

%% ─── Compactação de espaços verticais ────────────────────────────────────
%% Não altera margens nem entrelinha (norma ABNT/UFAM preservada); apenas
%% aproxima floats, listagens e listas do texto. Reduz páginas sem mexer no
%% corpo do texto. Para afrouxar, aumentar os valores abaixo.
\setlength{\textfloatsep}{10pt plus 2pt minus 2pt}  % float no topo/base da página x texto
\setlength{\intextsep}{8pt plus 2pt minus 2pt}       % floats [h] no meio do texto
\setlength{\floatsep}{8pt plus 2pt minus 2pt}        % entre dois floats
\setlength{\abovecaptionskip}{4pt}
\setlength{\belowcaptionskip}{2pt}
\lstset{aboveskip=6pt, belowskip=6pt}                % folga ao redor de código
\setlist{topsep=3pt, itemsep=2pt, parsep=0pt}        % listas mais justas (enumitem)
%% ──────────────────────────────────────────────────────────────────────────

%% ─── Topo de capítulo mais justo ─────────────────────────────────────────
%% A classe v3 herda o \@makechapterhead do book padrão, que reserva ~50pt de
%% espaço em branco acima do "Capítulo N". Reduz-se esse espaço (e o vão abaixo
%% do título) sem mexer em fonte, numeração ou estilo. Se o v3 customizar o
%% cabeçalho e os alvos não existirem, os \patchcmd simplesmente não agem (ramo
%% de falha vazio), sem erro. Para afrouxar, aumentar os "10\p@" / "18\p@".
\usepackage{etoolbox}
\makeatletter
\patchcmd{\@makechapterhead}{\vspace*{50\p@}}{\vspace*{12\p@}}{}{}
\patchcmd{\@makechapterhead}{\vskip 40\p@}{\vskip 20\p@}{}{}
\patchcmd{\@makeschapterhead}{\vspace*{50\p@}}{\vspace*{12\p@}}{}{}
\patchcmd{\@makeschapterhead}{\vskip 40\p@}{\vskip 20\p@}{}{}
\makeatother
%% ──────────────────────────────────────────────────────────────────────────

\setlength{\emergencystretch}{3em}

\setcounter{secnumdepth}{3}
\setcounter{tocdepth}{2}
\setlength{\parindent}{1.25cm}
\hyphenpenalty=1000
\tolerance=2000

\newif\ifperguntas
\ifbanca \perguntasfalse \else \perguntastrue \fi

\newtcolorbox{caixapergunta}[1]{%
  colback=orange!6!white, colframe=orange!75!black,
  boxrule=0.8pt, arc=2pt, boxsep=1pt,
  left=6pt, right=6pt, top=4pt, bottom=4pt,
  fonttitle=\bfseries\footnotesize,
  title={Pergunta ao orientador --- #1}}

\newcommand{\pergunta}[2]{%
  \ifperguntas
    \begin{caixapergunta}{#1}\footnotesize #2\end{caixapergunta}
  \fi}

\tcbuselibrary{breakable}
\newif\iflc
\ifbanca \lcfalse \else \lctrue \fi

\newtcolorbox{caixalc}{%
  breakable,
  colback=red!3!white, colframe=red!65!black,
  boxrule=0.6pt, arc=1pt, boxsep=1pt,
  left=6pt, right=6pt, top=4pt, bottom=4pt}

\newcommand{\lc}[1]{%
  \iflc
    \begin{caixalc}%
      \footnotesize\color{red!72!black}%
      \textbf{LC:}\ #1
    \end{caixalc}%
  \fi}

\AtBeginDocument{%
  \ifbanca
    \typeout{^^J*** VERSAO PARA A BANCA: marcacoes ocultas. ***^^J}%
  \else
    \typeout{^^J*** VERSAO DE TRABALHO: verde/vermelho/laranja/LC VISIVEIS. ***^^J%
             *** Para a banca, trocar \string\bancafalse por \string\bancatrue. ***^^J}%
  \fi}

\begin{document}

\title{Verificação Formal de Módulos Críticos do PX4 Autopilot usando ESBMC}

\foreigntitle{Formal Verification of Critical PX4 Autopilot Modules using ESBMC}

\author{Marcela Alves}{Rodrigues}

\advisor{Prof.}{Lucas Carvalho}{Cordeiro}{D.Sc.}

\socoadvisor{Prof.}{José Reginaldo}{Hughes de Carvalho}{D.Sc.}

\examiner{Prof.}{Raimundo da Silva Barreto}{D.Sc.}

\department{MEE}

\date{12}{2026}

\maketitle

\renewcommand\makefrontpage{%
  \thispagestyle{empty}
  \vspace*{\fill}
  \begin{center}Ficha Catalográfica\end{center}
  \begin{center}
  \fbox{\begin{minipage}{110mm}\vspace*{25mm}
  \vspace*{25mm}\end{minipage}}
  \end{center}
  \vspace*{\fill}
  \clearpage
  \thispagestyle{empty}
  \begin{center}FOLHA DE APROVAÇÃO\end{center}
  \vspace*{10mm}
  \begin{center}
  {\large\textbf{Marcela Alves Rodrigues}}\\[8mm]
  {\large\textbf{Verificação Formal de Módulos Críticos do PX4 Autopilot usando ESBMC}}\\[12mm]
  \end{center}
  %% LC: \textbf{[BLOQUEANTE]} Membros da banca ainda comentados. Preencher
  %% LC: antes de gerar a versão de depósito — a folha de aprovação sem os
  %% LC: examinadores impede o registro no PPGEE.
  \begin{center}
  BANCA EXAMINADORA\\[8mm]
  Prof. Dr. Lucas Carvalho Cordeiro --- Presidente (Orientador)\\
  PPGEE --- Universidade Federal do Amazonas (UFAM)\\[6mm]
  Prof. Dr. Raimundo da Silva Barreto --- Membro\\
  PPGI --- Universidade Federal do Amazonas (UFAM)\\[6mm]
  Prof. Dr. Yuri Bessa --- Membro\\
  PPGEE --- Universidade Federal do Amazonas (UFAM)\\[10mm]
  \end{center}
  \vspace*{\fill}
  \clearpage
  \thispagestyle{empty}
  \vspace*{\fill}
  \mbox{}
  \vspace*{\fill}
  \clearpage}

\frontmatter

\chapter*{Agradecimentos}
\addcontentsline{toc}{chapter}{Agradecimentos}

\begin{flushright}
\begin{minipage}{0.62\textwidth}
\raggedleft
\textit{``Consagre ao Senhor tudo o que você faz,\\
e os seus planos serão bem-sucedidos.''}\\[1mm]
Provérbios 16:3
\end{minipage}
\end{flushright}

\vspace{1mm}

Sempre acreditei que o amor transforma vidas; ao longo desta trajetória, o
mestrado me mostrou que o estudo também transforma.

Agradeço ao projeto SAGRES pela oportunidade de integrar uma cooperação tão
significativa, e estendo meu agradecimento ao Prof. Dr. Lucas Carvalho Cordeiro
e ao Prof. Dr. José Reginaldo Hughes de Carvalho pela orientação ao longo desta
pesquisa. Agradeço também à Universidade de Manchester pelo acolhimento durante
a visita técnica, experiência que ampliou minha formação e enriqueceu este
trabalho.

Agradeço, por fim, à minha família, que é o meu alicerce. À minha mãe, Maria, e
à minha avó, Maria, pelo amor e pelo apoio de sempre; e à minha tia Gely, pelo
carinho e pelo incentivo em cada etapa. De modo especial, agradeço ao meu tio
Raul, que me inspirou a seguir por esta área e cuja influência foi decisiva para
que eu chegasse até aqui.

O presente trabalho foi realizado com apoio da Coordenação de Aperfeiçoamento de
Pessoal de Nível Superior -- Brasil (CAPES) -- Código de Financiamento 001.

\begin{abstract}

Veículos Aéreos Não Tripulados dependem de software embarcado confiável para operar com segurança. O PX4 Autopilot é uma plataforma de código aberto para controle de voo cuja árvore \texttt{src/}, no commit analisado nesta pesquisa, reúne aproximadamente 710 mil linhas físicas de arquivos C/C++. Nesse contexto, este trabalho desenvolve e avalia uma abordagem de verificação formal aplicada a módulos críticos do PX4, organizada em duas campanhas experimentais. A primeira utiliza o ESBMC (\textit{Efficient SMT-Based Context-Bounded Model Checker}) para verificar propriedades de segurança em módulos selecionados, com \textit{harnesses} e \textit{stubs} para representar suas dependências. Foram analisados 12 módulos do PX4 Autopilot, no ramo \textit{main}, commit \texttt{8be2590c15}, utilizando o ESBMC v8.1.0 e o solver Boolector 3.2.4. O tempo total da campanha foi de aproximadamente 159 segundos, ou 94 segundos quando desconsiderado o módulo GPS uBlox, considerando a mediana de três execuções. Dez módulos não apresentaram violações das propriedades instrumentadas até o limite de desenrolamento adotado. Os outros dois, GPS uBlox e IMU BMI088, produziram três contraexemplos. Após a triagem, um foi classificado como violação alcançável no modelo, embora uma prova indutiva tenha demonstrado que o estado violador não ocorre no driver de produção. Os outros dois foram classificados como falso positivo e especificação disputada. Nenhum dos três achados foi classificado como defeito confirmado.

A segunda campanha aplica o mesmo método a dez módulos complementares --- oito módulos adicionais do PX4, o driver RoboClaw e o driver \texttt{batmon} ---, com o objetivo de avaliar a generalidade da abordagem além da amostra principal. A investigação do RoboClaw identificou um defeito na rotina de terminação de voo, em que o driver envia ao controlador do motor valores de byte superiores a 127, fora da faixa aceita pelo protocolo de comunicação; a ocorrência foi reproduzida por meio de \textit{replay} no código original, caracterizando um defeito reproduzido na implementação de produção, e foi também observada em bancada, numa Pixhawk~6X~RT com o controlador de motores emulado. No driver \texttt{batmon}, foi identificada uma escrita fora dos limites de um arranjo de tensões de célula. Os desfechos dos oito módulos adicionais foram heterogêneos, incluindo duplicatas de defeitos já corrigidos, um falso positivo e casos de baixa severidade. Em conjunto, as duas campanhas evidenciam a importância de complementar os resultados do verificador com a inspeção da implementação e a reprodução dos achados no código original.

\noindent
\textbf{Palavras-chave}: Verificação Formal; PX4 Autopilot; ESBMC; Bounded Model Checking; VANT; Sistemas Embarcados.
\end{abstract}

\begin{foreignabstract}

Unmanned Aerial Vehicles rely on reliable embedded software for safe operation. PX4 Autopilot is an open-source flight control platform whose \texttt{src/} tree, at the commit analysed in this study, contains approximately 710,000 physical lines in C/C++ files. This work develops and evaluates a formal verification approach applied to critical PX4 modules through two experimental campaigns. The first uses ESBMC (\textit{Efficient SMT-Based Context-Bounded Model Checker}) to verify safety properties in selected modules, with \textit{harnesses} and \textit{stubs} representing their dependencies. Twelve PX4 Autopilot modules were analysed on the \textit{main} branch, commit \texttt{8be2590c15}, using ESBMC v8.1.0 and the Boolector 3.2.4 solver. The campaign took approximately 159 seconds in total, or 94 seconds when the GPS uBlox module was excluded, based on the median of three runs. Ten modules showed no violations of the instrumented properties up to the adopted unwinding bound. The other two modules, GPS uBlox and IMU BMI088, produced three counterexamples. After triage, one was classified as a model-reachable violation, although an inductive proof demonstrated that the violating state does not occur in the production driver. The other two were classified as a false positive and a disputed specification. None of these three findings was classified as a confirmed defect.

The second campaign applies the same method to ten complementary modules --- eight additional PX4 modules, the RoboClaw driver and the \texttt{batmon} driver --- to assess the generality of the approach beyond the main sample. The investigation of RoboClaw identified a defect in the flight termination routine, in which the driver sends byte values greater than 127 to the motor controller, outside the range accepted by the communication protocol; the occurrence was reproduced through \textit{replay} on the original source code, characterising a defect reproduced in the production implementation; it was also observed on a bench, on a Pixhawk~6X~RT with an emulated motor controller. In the \texttt{batmon} driver, an out-of-bounds write to a cell-voltage array was identified. The outcomes of the eight additional modules were heterogeneous, including duplicates of already-fixed defects, a false positive and low-severity cases. Taken together, the two campaigns highlight the importance of complementing verifier results with implementation inspection and reproduction of findings in the original source code.

\noindent
\textbf{Keywords}: Formal Verification; PX4 Autopilot; ESBMC; Bounded Model Checking; UAV; Embedded Systems.
\end{foreignabstract}

\chapter*{Declaração de Uso de Inteligência Artificial}
\addcontentsline{toc}{chapter}{Declaração de Uso de Inteligência Artificial}
\vspace*{-0.9cm}

{\small

Durante a escrita desta dissertação, foi utilizada a ferramenta de Inteligência Artificial Generativa Claude (Anthropic), em atendimento à Portaria CNPq nº 2.664, de 6 de março de 2026, que institui a Política de Integridade na Atividade Científica do CNPq\corr{, e em consonância com as recomendações do ICMJE~\cite{ICMJE2026} e da WAME~\cite{WAME2023} sobre a declaração do uso de IA em trabalhos científicos}. O Claude foi empregado entre novembro de 2025 e a conclusão deste trabalho, em versões disponibilizadas pelo fornecedor ao longo do período, como apoio na escrita, na revisão bibliográfica e na parte experimental. Em todas as frentes, a autoria, as decisões técnicas e o conteúdo final permanecem de responsabilidade da autora.

Na escrita, a ferramenta auxiliou na revisão gramatical, na clareza do texto em português e em ajustes pontuais de formatação em LaTeX. Na revisão bibliográfica, o Google Scholar e gerenciadores de referências com recursos de IA foram usados para localizar trabalhos relacionados; toda referência incluída no arquivo \textit{.bib} foi conferida pela autora, após leitura do resumo original, antes de ser citada.

Na parte experimental, os \textit{harnesses} de verificação foram escritos pela autora a partir do código-fonte do PX4, com o Claude como apoio na sintaxe de \textit{assertions} do ESBMC, no uso de valores não determinísticos e na composição dos comandos de execução. A definição das propriedades e a versão final de cada \textit{harness} são da autora, que conduziu todas as execuções do ESBMC e do CBMC. A execução das verificações, a interpretação dos contraexemplos e a triagem dos achados foram feitas exclusivamente pela autora; o Claude não apontou achados nem indicou trechos suspeitos. Na revisão final, a ferramenta também auxiliou a conferir os números contra os logs de execução e a montar em \LaTeX{} as figuras dos contraexemplos. Todo o material gerado com apoio da ferramenta foi revisado pela autora, que permanece responsável pelo conteúdo final desta dissertação.

\par}

\tableofcontents

\cleardoublepage
\addcontentsline{toc}{chapter}{\listfigurename}
\listoffigures

\cleardoublepage
\addcontentsline{toc}{chapter}{\listtablename}
\listoftables

\chapter*{Lista de Siglas}
\addcontentsline{toc}{chapter}{Lista de Siglas}

\begin{description}
\item[AST]     Abstract Syntax Tree
\item[BDD]     Binary Decision Diagram
\item[BMC]     Bounded Model Checking
\item[BSD]     Berkeley Software Distribution
\item[CAN]     Controller Area Network
\item[CAPES]   Coordenação de Aperfeiçoamento de Pessoal de Nível Superior
\item[CBMC]    C Bounded Model Checker
\item[CI/CD]   Continuous Integration / Continuous Delivery
\item[CRC]     Cyclic Redundancy Check
\item[CTL]     Computation Tree Logic
\item[CVE]     Common Vulnerabilities and Exposures
\item[CWE]     Common Weakness Enumeration
\item[DDS]     Data Distribution Service
\item[DOI]     Digital Object Identifier
\item[EKF2]    Extended Kalman Filter 2
\item[ESBMC]   Efficient SMT-Based Context-Bounded Model Checker
\item[FIFO]    First In, First Out
\item[GPS]     Global Positioning System
\item[HAL]     Hardware Abstraction Layer
\item[HIL]     Hardware-in-the-Loop
\item[I2C]     Inter-Integrated Circuit
\item[IA]      Inteligência Artificial
\item[IEEE]    Institute of Electrical and Electronics Engineers
\item[IMU]     Inertial Measurement Unit
\item[LLVM]    Low Level Virtual Machine
\item[LoC]     Lines of Code
\item[LPE]     Local Position Estimator
\item[LTL]     Linear Temporal Logic
\item[MAVLink] Micro Air Vehicle Link
\item[MMU]     Memory Management Unit
\item[PID]     Proporcional-Integral-Derivativo
\item[POSIX]   Portable Operating System Interface
\item[PPGEE]   Programa de Pós-Graduação em Engenharia Elétrica
\item[PR]      \textit{Pull Request}
\item[PWM]     Pulse Width Modulation
\item[PX4]     PX4 Autopilot
\item[QE]      Questão Experimental
\item[RC]      Radio Control
\item[RTPS]    Real-Time Publish-Subscribe
\item[SAT]     Satisfiability
\item[SBESC]   Simpósio Brasileiro de Engenharia de Sistemas Computacionais
\item[SITL]    Software-in-the-Loop
\item[SMT]     Satisfiability Modulo Theories
\item[SPI]     Serial Peripheral Interface
\item[SV-COMP] International Competition on Software Verification
\item[UART]    Universal Asynchronous Receiver/Transmitter
\item[UAV]     Unmanned Aerial Vehicle
\item[UAVCAN]  Uncomplicated Application-level Vehicular Computing and Networking
\item[UBX]     u-blox Binary Protocol
\item[UNSAT]   Unsatisfiable
\item[uORB]    micro Object Request Broker
\item[VAM]     Violação Alcançável no Modelo
\item[VANT]    Veículo Aéreo Não Tripulado
\item[VCC]     Verification Condition
\item[VTOL]    Vertical Take-Off and Landing
\end{description}

\mainmatter
\pagestyle{fancy}
\chapter{Introdução}
\label{cap:introducao}

Os Veículos Aéreos Não Tripulados (VANTs) são utilizados em aplicações como monitoramento ambiental, agricultura de precisão, inspeção de infraestrutura crítica e operações de busca e salvamento~\cite{Wanner2024,Siddiqi2022}. A expansão desses usos foi favorecida pelos avanços tecnológicos e pela redução dos custos dos equipamentos. Independentemente da finalidade da missão, a operação desses veículos depende de sistemas de controle embarcados, cujo funcionamento exige confiabilidade e segurança~\cite{Mohsan2023}. Falhas nesses sistemas podem ocasionar a perda do equipamento e, dependendo do contexto operacional, oferecer riscos às pessoas e ao meio ambiente~\cite{Wanner2024}.

Esse requisito de confiabilidade deixa de ser abstrato em missões realizadas em regiões remotas. No monitoramento ambiental da Amazônia, por exemplo, uma falha em pleno voo pode significar a perda de um equipamento de difícil recuperação e a interrupção de uma coleta de dados que nem sempre pode ser repetida. É nesse cenário que se insere esta pesquisa, vinculada ao projeto SAGRES, que emprega VANTs nesse tipo de operação. Examinar o comportamento do firmware que controla o voo, antes de confiá-lo a uma missão dessas, é parte do que torna essa operação segura.

Entre as plataformas utilizadas no desenvolvimento de sistemas de controle de voo está o PX4 Autopilot, uma solução de código aberto para VANTs, implementada em C e C++~\cite{PX4_GitHub}. Em um projeto de grande porte como esse, testes e revisão de código são fundamentais para avaliar a qualidade do software. Entretanto, essas atividades não asseguram que todos os comportamentos possíveis sejam examinados, pois os testes exercitam cenários definidos durante sua elaboração~\cite{Godefroid2005}. Assim, problemas associados a combinações específicas de entradas, condições extremas ou limites de operação podem permanecer sem identificação.

Entre esses problemas estão os \textit{overflows} aritméticos, os acessos inválidos à memória e as violações de limites estabelecidos para sensores. Para investigar situações dessa natureza, esta pesquisa aplica verificação formal por meio de \textit{Bounded Model Checking} (BMC). A técnica analisa sistematicamente as execuções representadas no modelo até um limite de desenrolamento previamente definido. Quando uma propriedade é violada dentro desse limite, o verificador pode produzir um contraexemplo que representa uma sequência de execução associada à violação~\cite{Biere1999}. Dessa forma, o BMC permite examinar combinações de entradas e condições que podem não ter sido contempladas nos testes tradicionais~\cite{Beyer2022}.

O alcance dessa análise depende do modelo construído, das propriedades especificadas e do limite de desenrolamento adotado. Portanto, a ausência de contraexemplos indica que nenhuma violação foi encontrada nas condições analisadas, mas não demonstra que o programa esteja livre de defeitos. Essa ressalva é especialmente importante em sistemas embarcados, nos quais dependências de \textit{hardware} e componentes externos precisam ser representados de maneira controlada para permitir a análise dos módulos isoladamente.

Nesta pesquisa, o BMC é aplicado por meio do ESBMC (\textit{Efficient SMT-Based Context-Bounded Model Checker})~\cite{Cordeiro2012}, ferramenta de verificação formal com suporte às linguagens C e C++, utilizadas na implementação do PX4. A escolha considera o suporte a construções modernas de C++, a integração com diferentes solvers SMT e a experiência do grupo de pesquisa com a ferramenta~\cite{Menezes2024,Song2024}. Entre os solvers disponíveis estão Z3~\cite{deMoura2008}, Boolector~\cite{Brummayer2009}, cvc5~\cite{Barbosa2022} e Bitwuzla~\cite{Niemetz2023}.

Outras ferramentas, como CBMC e LLBMC, também foram consideradas. O LLBMC opera sobre a representação intermediária do LLVM, exigindo uma cadeia de compilação compatível com essa representação~\cite{Merz2012}. O ESBMC, por sua vez, utiliza um \textit{front-end} baseado na AST do Clang para processar construções de C++~\cite{Song2024}, aspecto relevante para a base do PX4, que utiliza C++17. O CBMC não foi descartado: nesta pesquisa, ele é empregado na validação cruzada de achados selecionados. Assim, o ESBMC foi adotado como ferramenta principal, enquanto o CBMC foi utilizado como ferramenta independente para confrontar resultados.

Como a análise produz contraexemplos relativos ao modelo, sua interpretação exige uma etapa de triagem. Um contraexemplo demonstra que a propriedade codificada foi violada nas condições representadas pelo modelo, mas não comprova, isoladamente, que a mesma situação ocorra no código de produção. Por isso, esta dissertação adota uma classificação que distingue o resultado da verificação da avaliação de sua correspondência com a implementação real.

\begin{tcolorbox}[colback=gray!4!white, colframe=gray!55!black,
boxrule=0.6pt, arc=2pt, left=8pt, right=8pt, top=6pt, bottom=6pt,
title={\textbf{Classificação dos achados}}]
\footnotesize
\begin{description}[leftmargin=3.0cm, style=nextline, font=\normalfont\bfseries]

\item[Contraexemplo] Traço de execução produzido pelo verificador quando uma propriedade é violada no modelo. Não comprova, isoladamente, que a situação ocorra no código de produção.

\item[Candidato] Contraexemplo submetido à triagem, mas ainda não classificado como violação alcançável no modelo, falso positivo ou especificação disputada.

\item[Violação alcançável no modelo (VAM)] Candidato cuja violação ocorre no modelo verificado, sem que o \textit{harness} tenha omitido guardas presentes no código de produção. A possibilidade de o estado violador ocorrer no driver real deve ser examinada separadamente, por meio da conferência com a implementação.

\item[Falso positivo] Candidato cuja violação decorre da omissão, no \textit{harness}, de uma guarda existente no código de produção. A conferência com a implementação demonstra que a execução apresentada não representa o comportamento do trecho real nas mesmas condições.

\item[Especificação disputada] Candidato cuja violação ocorre em relação à propriedade codificada, mas cuja interpretação, baseada em documentação técnica do fabricante, não é reconhecida pelos mantenedores como requisito do código.

\item[Defeito reproduzido no código de produção] Problema cuja ocorrência foi reproduzida na implementação original, por meio de uma execução ou \textit{replay} que preserva as condições relevantes do comportamento observado.

\item[Defeito confirmado] Achado reconhecido como defeito pelos mantenedores do projeto. Essa classificação é distinta da reprodução técnica de um problema no código de produção.

\end{description}
\end{tcolorbox}

\noindent
\footnotesize
\textit{Nota:} \textit{harness} é o programa que instancia o trecho de código sob verificação, fornece entradas não determinísticas e propriedades a serem avaliadas e substitui dependências externas por \textit{stubs}.
\normalsize

Essa distinção é importante para interpretar os resultados desta pesquisa. No GPS uBlox (Seção~\ref{sec:caso-gps}), por exemplo, a violação foi classificada como alcançável no modelo porque o estado violador era admitido pelas condições do \textit{harness}. Entretanto, a análise do código de produção e a prova indutiva apresentada nesta dissertação (Seção~\ref{subsec:gps-inducao}) demonstraram que esse estado não ocorre no driver real. O mantenedor do PX4 enquadrou a guarda proposta como endurecimento defensivo. Já o candidato relacionado ao contador de FIFO do BMI088 foi classificado como falso positivo, pois o \textit{harness} não reproduzia uma guarda presente no driver de produção (Seção~\ref{subsec:imu-fifo}). Esses casos mostram por que a análise dos contraexemplos precisa considerar tanto o modelo quanto a implementação real.

A aplicação da técnica ao PX4 considera um conjunto delimitado de módulos, correspondente ao ramo \textit{main} no commit \texttt{8be2590c15}. O objetivo não é verificar integralmente o firmware, mas avaliar componentes selecionados conforme sua relevância operacional, sua representatividade na arquitetura e a possibilidade de isolá-los para análise. Para isso, a metodologia contempla a seleção dos módulos, a criação de \textit{stubs} para abstrair dependências de \textit{hardware}, a especificação de propriedades de segurança e a análise dos contraexemplos. Os achados também podem ser conferidos com o código de produção, reportados à comunidade responsável pelo projeto e submetidos à validação cruzada com outra ferramenta.

Foram selecionados 12 módulos, abrangendo sete classes funcionais do sistema de controle de voo --- de sensores seriais à estimação de estado ---, cuja composição completa é apresentada no Capítulo~\ref{cap:metodologia} (Tabela~\ref{tab:modulos}). Essa distribuição busca contemplar componentes com responsabilidades distintas, evitando concentrar a análise em uma única parte do firmware.

Além da criticidade operacional e da representatividade arquitetural, a possibilidade de isolar cada módulo por meio de \textit{stubs} foi considerada na seleção. Esse critério facilita a execução da verificação, mas também pode introduzir viés na composição da amostra, uma limitação discutida na Seção~\ref{sec:ameacas-validade}. A seleção buscou, portanto, equilibrar a profundidade da análise com a representação de diferentes componentes do PX4.

Como o estudo de caso integra o projeto SAGRES, voltado a esse tipo de missão em região remota, o enquadramento regulatório dessas operações também merece registro.

O enquadramento das operações em linha de visada visual (VLOS) e além da linha de visada visual (BVLOS) no espaço aéreo brasileiro é estabelecido pela ICA~100-40 do DECEA~\cite{DECEA_ICA10040}. Como os requisitos variam conforme o tipo de operação, a realização de uma missão em região remota não significa, por si só, que ela seja BVLOS. Para esta pesquisa, o aspecto relevante é a dificuldade de acesso e recuperação do equipamento em determinadas localidades, independentemente do regime de operação.

A partir desse contexto, a pesquisa busca responder à seguinte questão: é viável aplicar \textit{Bounded Model Checking} a um firmware de controle de voo de porte industrial, com tempo de verificação por módulo compatível com o uso prático, produzindo evidências auditáveis sobre propriedades de segurança relacionadas à memória e à aritmética, respeitados os limites do modelo e do desenrolamento?

Essa questão é desdobrada em três questões experimentais, respondidas no Capítulo~\ref{cap:experimentos}:

\begin{description}[leftmargin=1.5cm, style=nextline]
\item[\textbf{QE1}] Qual é o tempo de verificação por módulo em \textit{hardware} de prateleira?

\item[\textbf{QE2}] O BMC expõe violações alcançáveis no modelo em código que os testes unitários existentes do PX4 não referenciam?

\item[\textbf{QE3}] Quais configurações são necessárias para tornar cada módulo tratável pelo verificador?
\end{description}

\section{Delimitação: método e estudo de caso}
\label{sec:delimitacao}

Esta pesquisa tem como contribuição principal uma metodologia de verificação formal para firmware de sistemas críticos. O PX4 Autopilot, no ramo \textit{main} e no commit \texttt{8be2590c15}, constitui o estudo de caso utilizado para aplicar e avaliar o procedimento. Essa distinção permite separar a contribuição metodológica dos resultados obtidos especificamente nessa plataforma.

A metodologia reúne etapas que podem ser aplicadas a diferentes bases de código C/C++ de grande porte: seleção de unidades verificáveis (i), execução do verificador (iv), análise dos contraexemplos com conferência do código de produção (v), reporte à comunidade responsável pelo projeto (vi), validação cruzada entre ferramentas e verificação de \textit{patches} (vii), além da portabilidade entre projetos e da integração contínua (viii). Essas seis etapas são consideradas agnósticas ao domínio e estão descritas no Capítulo~\ref{cap:metodologia} (Figura~\ref{fig:metodologia-visao-geral}).

Duas etapas, por sua vez, foram particularizadas para firmware embarcado: a instrumentação com \textit{stubs} de \textit{hardware} (ii), utilizada para abstrair barramentos, sensores e chamadas ao sistema operacional de tempo real; e a especificação de propriedades derivadas de \textit{datasheets} (iii), que traduz limites declarados pelos fabricantes em asserções sobre variáveis do programa. Embora essas etapas exijam informações próprias do contexto embarcado, a estrutura geral da metodologia pode ser adaptada a outros sistemas.

Essa adaptação, contudo, não foi demonstrada em diferentes domínios nesta pesquisa. A evidência está concentrada no estudo de caso do PX4 e em uma replicação mínima da lógica equivalente no ArduPilot, descrita no Capítulo~\ref{cap:metodologia}, Etapa~viii. A aplicação a sistemas industriais, automotivos ou dispositivos médicos embarcados exigiria adaptar as dependências abstraídas e as propriedades verificadas às características de cada alvo. Portanto, a transferência da metodologia para esses contextos permanece uma possibilidade a ser investigada.

A avaliação experimental foi organizada em duas campanhas. A primeira compreende os 12 módulos da amostra principal, analisados por meio de \textit{harnesses} e \textit{stubs}. A segunda aplica o mesmo método a dez módulos complementares --- oito módulos adicionais do PX4, o driver RoboClaw e o driver \texttt{batmon} ---, com o objetivo de avaliar a generalidade da abordagem além da amostra principal. Seus desfechos são heterogêneos: a maioria confirma padrões já conhecidos ou não produz violação, enquanto o driver RoboClaw apresentou um defeito reproduzido no código de produção. Essa segunda campanha não integra as tabelas consolidadas da primeira nem os resultados do artigo derivado; o driver \texttt{batmon}, em particular, foi conferido em um \textit{commit} posterior do ramo \textit{main} (\texttt{c9b9e4e25e}). A apresentação separada preserva a distinção entre a amostra principal e as investigações complementares.

\section{Objetivos}
\label{sec:objetivos}

O objetivo geral desta pesquisa é desenvolver e avaliar uma metodologia de verificação formal baseada em \textit{Bounded Model Checking} para identificar e triar candidatos a violações de propriedades de segurança em módulos críticos do firmware PX4 Autopilot, utilizando o ESBMC como ferramenta principal.

Para alcançar esse objetivo, foram definidos três objetivos específicos, relacionados às questões experimentais:

\begin{description}[leftmargin=1.5cm, style=nextline, font=\normalfont\bfseries]

\item[O1] Propor uma metodologia sistemática e replicável para a verificação formal de firmware embarcado em C/C++, contemplando a seleção de módulos críticos, a instrumentação com \textit{stubs} para abstração de dependências de \textit{hardware}, a especificação de propriedades de segurança derivadas de \textit{datasheets} e a análise dos contraexemplos gerados pelo ESBMC~\cite{Cordeiro2012}. Esse objetivo está relacionado principalmente à \textbf{QE3}.

\item[O2] Aplicar a metodologia ao PX4 Autopilot em duas campanhas experimentais, verificando módulos de diferentes naturezas e classificando os candidatos encontrados como violações alcançáveis no modelo, falsos positivos, especificações disputadas ou defeitos reproduzidos no código de produção. Os achados também são submetidos à avaliação da comunidade PX4, por meio do reporte de contraexemplos e propostas de correção. Esse objetivo está relacionado à \textbf{QE2}.

\item[O3] Avaliar a viabilidade da abordagem por meio da medição do tempo de verificação de cada módulo e do registro das configurações necessárias para tornar a análise tratável. Esse objetivo está relacionado à \textbf{QE1} e à \textbf{QE3}.

\end{description}

A Tabela~\ref{tab:rastreabilidade} relaciona os objetivos específicos às etapas da metodologia, aos resultados utilizados em sua avaliação e às respectivas questões experimentais.

\begin{table}[H]
\centering
\caption[Rastreabilidade entre objetivos, metodologia, resultados e questões]{Rastreabilidade entre objetivos, metodologia, resultados e questões experimentais.}
\label{tab:rastreabilidade}
\footnotesize
\begin{tabular}{@{}llll@{}}
\toprule
\textbf{Objetivo} & \textbf{Metodologia} &
\textbf{Resultados} & \textbf{Questão} \\
\midrule
O1 & Cap.~\ref{cap:metodologia} (etapas i--viii) &
Cap.~\ref{cap:experimentos} & QE3 \\
O2 & Seções~\ref{sec:selecao-modulos} e
\ref{sec:caso-gps} &
Cap.~\ref{cap:experimentos} & QE2 \\
O3 & Cap.~\ref{cap:metodologia} (etapa iv) &
Cap.~\ref{cap:experimentos} & QE1, QE3 \\
\bottomrule
\end{tabular}
\end{table}

\section{Justificativa}
\label{sec:justificativa}

A verificação de software embarcado apresenta desafios relacionados às restrições de memória, processamento e tempo de resposta impostas pelo ambiente de execução~\cite{Mzid2024}. No caso do PX4, os alvos considerados utilizam o modo \textit{flat} do NuttX, no qual não há isolamento de memória entre os componentes da aplicação. Embora microcontroladores STM32 disponham de unidade de proteção de memória (MPU), sua utilização depende dos modos \textit{protected} ou \textit{kernel} do NuttX, e não do modo \textit{flat} adotado nos alvos analisados~\cite{NuttX}. Essa característica reforça a relevância de investigar possíveis violações de propriedades de memória e aritmética nos módulos do sistema de controle de voo.

Nesse contexto, o BMC complementa os testes tradicionais ao analisar as execuções representadas no modelo até um limite de desenrolamento definido. A técnica pode revelar violações associadas a combinações de entradas e condições não contempladas nos casos de teste elaborados manualmente~\cite{Biere1999,Beyer2022}. Seus resultados, entretanto, dependem das propriedades especificadas e da fidelidade do modelo, razão pela qual a análise dos contraexemplos e a conferência com o código de produção são partes essenciais do procedimento.

A pesquisa também se justifica por sua aplicação no projeto SAGRES, voltado ao monitoramento ambiental em regiões remotas da Amazônia. Nessas localidades, a recuperação de equipamentos após falhas pode ser complexa, e a interrupção de uma operação pode comprometer a continuidade da coleta de dados. A verificação de propriedades de segurança dos componentes do PX4 permite investigar possíveis problemas no firmware antes de sua utilização operacional.

Por fim, a documentação do procedimento e dos artefatos experimentais permite que outras pesquisas examinem e reproduzam as etapas realizadas. A contribuição não consiste em afirmar que a metodologia elimina todos os defeitos ou garante a segurança integral do sistema, mas em oferecer um processo sistemático para selecionar módulos, instrumentá-los, verificar propriedades delimitadas e analisar as evidências produzidas.

\section{Contribuições Científicas}
\label{sec:contribuicoes}

Este trabalho apresenta cinco contribuições principais. As três primeiras, relacionadas à primeira campanha, estão alinhadas, em conteúdo, às contribuições do artigo derivado~\cite{Rodrigues2026}; a demonstração formal da inalcançabilidade do estado violador do GPS no driver de produção, incluída na terceira, constitui um resultado próprio desta dissertação. A quarta amplia a investigação para a segunda campanha, e a quinta disponibiliza o artefato público de reprodução.

\begin{enumerate}

\item \textbf{Metodologia de verificação em oito etapas:} definição de um procedimento para verificação formal de firmware, com seis etapas agnósticas ao domínio e duas específicas para sistemas embarcados: instrumentação com \textit{stubs} de \textit{hardware} e especificação de propriedades derivadas de \textit{datasheets}. O procedimento também contempla validação cruzada entre ferramentas e verificação de \textit{patches} (Capítulo~\ref{cap:metodologia}).

\item \artigo{\textbf{Evidência experimental delimitada pelo modelo:} os \textit{harnesses} de dez módulos do PX4 foram verificados sem violações das propriedades instrumentadas de segurança de memória e de \textit{overflow} aritmético, até o limite de desenrolamento adotado. O tempo total da campanha com os 12 módulos foi de aproximadamente 159 segundos, reduzido a 94 segundos quando desconsiderado o GPS uBlox (Capítulo~\ref{cap:experimentos}).}

\item \novo{\textbf{Triagem documentada dos achados da campanha:} classificação dos três contraexemplos produzidos em dois módulos, distinguindo uma violação alcançável no modelo, um falso positivo e uma especificação disputada. No driver GPS uBlox, foi identificada uma violação alcançável no modelo (CWE-191 $\rightarrow$ CWE-119, Issue~\#26866), para a qual foi proposta uma guarda no PR~\#27576. A análise desta dissertação demonstrou, por prova indutiva com ESBMC e CBMC, que o estado violador não ocorre no driver de produção, a partir da análise do laço real de \texttt{dumpGpsData()} (Seções~\ref{subsec:patch-gps} e \ref{subsec:gps-inducao}). No BMI088, o candidato relacionado ao contador de FIFO foi classificado como falso positivo, pois o \textit{harness} não reproduzia uma guarda presente no código de produção. Já a observação de temperatura (Issue~\#26865, PR~\#27575) foi classificada como especificação disputada, uma vez que a interpretação da propriedade baseada no \textit{datasheet} não foi aceita pelos mantenedores. Essa observação também foi reproduzida no ArduPilot (Issue~\#33272).}

\item \novo{\textbf{Segunda campanha --- aplicação complementar do método:} aplicação da metodologia a dez módulos complementares (oito módulos adicionais do PX4, o driver RoboClaw e o driver \texttt{batmon}), para avaliar a generalidade da abordagem além da amostra principal. A investigação do RoboClaw identificou um defeito na rotina de terminação de voo, em que o driver envia ao controlador do motor valores de byte superiores a 127, fora da faixa aceita pelo protocolo de comunicação; a ocorrência foi reproduzida por \textit{replay} no código original e observada em bancada, numa Pixhawk~6X~RT com o controlador emulado. No driver \texttt{batmon}, foi identificada uma escrita fora dos limites do arranjo \texttt{voltage\_cell\_v[14]} quando a contagem de células atinge 16 (\textit{commit} \texttt{c9b9e4e25e}). Os desfechos dos oito módulos adicionais foram heterogêneos: duplicatas de defeitos já corrigidos, um falso positivo, dois com \textit{issue} de baixa severidade e um sem violação.}

\item \artigo{\textbf{Artefato público para reprodução:} disponibilização dos \textit{harnesses}, registros de execução e scripts utilizados na pesquisa, permitindo consultar e reproduzir os experimentos (\url{https://github.com/marcela2881/dissertacao-px4-esbmc}; Zenodo, DOI \texttt{10.5281/zenodo.23282236}).}

\end{enumerate}

\section{Organização do Trabalho}
\label{sec:organizacao}

O documento está estruturado em seis capítulos. O Capítulo~\ref{cap:fundamentacao} apresenta a fundamentação teórica sobre verificação formal, \textit{model checking}, \textit{bounded model checking}, solvers SMT, ESBMC e arquitetura do PX4 Autopilot. O Capítulo~\ref{cap:relacionados} discute trabalhos relacionados à aplicação de verificação formal em sistemas embarcados e veículos autônomos.

O Capítulo~\ref{cap:metodologia} descreve o procedimento adotado, incluindo os critérios de seleção dos módulos, a criação de \textit{stubs}, a definição das propriedades de segurança e a configuração do ESBMC. O Capítulo~\ref{cap:experimentos} apresenta os experimentos e resultados, incluindo a análise dos contraexemplos, a classificação dos achados e a avaliação da viabilidade da abordagem. Por fim, o Capítulo~\ref{cap:conclusao} reúne as conclusões, contribuições, limitações e possibilidades de trabalhos futuros.

O documento é complementado por três apêndices: o Apêndice~\ref{apend:producoes} registra a produção científica derivada desta dissertação; o Apêndice~\ref{apend:harness-gps} apresenta o \textit{harness} do GPS uBlox; e o Apêndice~\ref{apend:contraexemplos} reúne os contraexemplos gerados pelo ESBMC para os três candidatos submetidos à triagem.

\chapter{Fundamentação Teórica}
\label{cap:fundamentacao}

Este capítulo reúne os conceitos que sustentam as decisões metodológicas do
trabalho: o funcionamento do \textit{Bounded Model Checking}, as razões para o
uso do ESBMC em firmware embarcado e a arquitetura do PX4 Autopilot, alvo da
verificação.

\section{Verificação Formal de Software}

A verificação formal emprega métodos matemáticos para avaliar se um modelo do
sistema satisfaz propriedades previamente especificadas~\cite{Clarke1986}.
Diferentemente dos testes convencionais, que examinam o comportamento em cenários
selecionados, essas técnicas analisam sistematicamente os comportamentos
representados pelo modelo, de acordo com as hipóteses e os limites estabelecidos.

Este trabalho concentra-se em propriedades de \textit{safety} --- condições que
não devem ocorrer, como acessos fora dos limites de \textit{arrays} ou
desreferência de ponteiros nulos. A escolha decorre do domínio: em firmware
embarcado sem proteção de memória robusta, uma violação de \textit{safety} pode
corromper o estado global e causar falha catastrófica em voo. Além disso, o
ESBMC verifica automaticamente diversas propriedades de \textit{safety} --- como
limites de \textit{arrays} e validade de ponteiros --- e permite que a
pesquisadora especifique propriedades adicionais por meio de instruções
\texttt{assert}, de modo que a análise cubra tanto condições de segurança de
memória quanto comportamentos próprios de cada módulo. Propriedades de
\textit{liveness} --- que garantem que algo desejável acabe ocorrendo --- ficaram
fora do escopo: o suporte do ESBMC a elas limita-se à terminação e à detecção de
laços não terminantes, sem verificação geral de LTL~\cite{Menezes2024}.
%% LC (resolvido): afirmacao sobre liveness recalibrada ao escopo real do ESBMC
%% (terminacao/lacos nao terminantes, nao LTL geral); citacao Clarke1986 -> Menezes2024.

%% LC (resolvido): DART (Godefroid2005) e teste concolico; execucao simbolica passou
%% a citar o artigo original, King1976.
Entre as técnicas de verificação formal destacam-se o \textit{Model
Checking}~\cite{Clarke1981}, que explora automaticamente o espaço de estados; o
\textit{Theorem Proving}~\cite{Huth2004,Tang2025}, que constrói provas com auxílio de
assistentes de prova; a \textit{Abstract Interpretation}~\cite{Cousot1977}, que
aproxima o comportamento por domínios abstratos; e a \textit{Symbolic
Execution}~\cite{King1976}, que executa o programa com valores simbólicos. Pela
automação e pela adequação a software embarcado, adota-se aqui o Model Checking.

\section{Model Checking}

O Model Checking, introduzido por Clarke e Emerson~\cite{Clarke1981} e, de forma
independente, por Queille e Sifakis~\cite{Queille1982}, constrói um modelo
matemático do sistema --- tipicamente uma máquina de estados finita --- e
verifica algoritmicamente se ele satisfaz uma especificação em lógica
temporal~\cite{Clarke1986}. Como esse modelo é uma abstração do sistema real, a
validade da conclusão depende da fidelidade da representação construída.

\subsection{Processo de Model Checking}

O processo tem três etapas. Na modelagem, o sistema é representado como uma
estrutura de Kripke, em que os vértices são estados e as arestas, transições. As
propriedades são escritas em lógicas temporais como LTL (\textit{Linear Temporal
Logic})~\cite{Pnueli1977} ou CTL (\textit{Computation Tree Logic})~\cite{Clarke1986};
uma propriedade LTL típica, $\square(\textit{request} \rightarrow
\Diamond\textit{response})$, lê-se ``toda requisição é eventualmente
respondida'', com $\square$ denotando ``sempre'' e $\Diamond$, ``eventualmente''.
Por fim, o verificador percorre os estados e, se algum violar a propriedade,
retorna um contraexemplo --- a sequência de transições que leva à
violação~\cite{Clarke1986}. Essa geração automática de contraexemplos é o que
torna o Model Checking atraente para sistemas complexos.

%% LC (resolvido): algoritmo restrito a invariantes/safety -- renomeado e com
%% ressalva de que propriedades temporais gerais exigem rotulacao CTL ou automato
%% de Buchi (Clarke1986). LC (resolvido): acrescentado o mapa de predecessores.
O Algoritmo~\ref{alg:mc} apresenta, de forma simplificada, a busca de invariantes
no espaço de estados; é uma ilustração do princípio, e não o procedimento
completo das demais verificações temporais. Propriedades temporais gerais não se
decidem olhando um estado isolado: exigem a rotulação de CTL ou a busca de um
ciclo aceitante em um autômato de Büchi~\cite{Clarke1986}. O caso aqui ilustrado
--- \textit{safety}/invariantes --- é o empregado nesta pesquisa.

\begin{algorithm}
\caption{Algoritmo de verificação de invariantes por busca no espaço de estados}
\label{alg:mc}
\begin{algorithmic}[1]
\Require Modelo $M = (S, S_0, T, L)$, propriedade $\phi$
\Ensure \textit{verdadeiro} se $M \models \phi$, contraexemplo caso contrário
\State $estados \leftarrow S_0$
\State $visitados \leftarrow \emptyset$
\While{$estados \neq \emptyset$}
    \State $s \leftarrow \text{remover}(estados)$
    \If{$s \notin visitados$}
        \State $visitados \leftarrow visitados \cup \{s\}$
        \If{$s \not\models \phi$}
            \State \Return contraexemplo até $s$
        \EndIf
        \State $estados \leftarrow estados \cup \text{sucessores}(s, T)$
    \EndIf
\EndWhile
\State \Return \textit{verdadeiro}
\end{algorithmic}
\end{algorithm}

O conjunto $estados$ representa os estados ainda não explorados; a estrutura usada
para selecionar o próximo elemento determina a estratégia de busca --- uma fila
leva à busca em largura, uma pilha à busca em profundidade ---, e ambas preservam
a corretude para invariantes, desde que todos os estados alcançáveis sejam
explorados. Na prática, mantém-se um mapa de predecessores durante a busca: ao
encontrar um estado violador, o contraexemplo é reconstruído seguindo esse mapa,
do estado inicial até ele.

\subsection{Explosão de Estados}

O principal obstáculo do Model Checking é a explosão combinatória: o número de
estados cresce exponencialmente com o número de variáveis ($2^n$ para $n$
variáveis booleanas). O driver do IMU BMI088, por exemplo, manipula leituras
brutas do sensor e valores convertidos em ponto flutuante de 32 bits; um modelo que
tratasse como simbólicas apenas dez dessas variáveis de 32 bits já ultrapassaria
$2^{320}$ configurações, inviabilizando a exploração exaustiva.
%% LC (resolvido): hipotese do 2^320 declarada (dez variaveis de 32 bits).
Diversas técnicas atacam esse problema ---
BDDs (\textit{Binary Decision Diagrams}), abstração, redução de ordem parcial e o
\textit{Bounded Model Checking}, que limita a profundidade de
exploração~\cite{Biere1999}.

\section{Bounded Model Checking}

O \textit{Bounded Model Checking} (BMC), proposto por Biere et
al.~\cite{Biere1999}, restringe a busca por contraexemplos a execuções de até $k$
transições, com $k$ definido pelo usuário, em vez de explorar todos os estados
alcançáveis.

\subsection{Formulação de BMC}

Dado o estado inicial $I$, a relação de transição $T$ e uma propriedade de
\textit{safety} $\phi$, o BMC codifica a busca por contraexemplos como um
problema de satisfatibilidade:

\begin{equation}
BMC_{\phi}(k) = I(s_0) \wedge \bigwedge_{i=0}^{k-1} T(s_i, s_{i+1}) \wedge
\bigvee_{i=0}^{k} \neg\phi(s_i)
\label{eq:bmc}
\end{equation}

O termo $I(s_0)$ fixa as condições iniciais; a conjunção central codifica $k$
transições válidas; e a disjunção final verifica se algum estado $s_i$ viola
$\phi$. Se a fórmula for satisfatível, o solver devolve uma atribuição que
constitui um contraexemplo concreto~\cite{Biere1999}.

%% LC (resolvido): incluir exemplo completo (laco C -> desenrolamento -> SSA ->
%% VCC -> consulta SMT). LC (resolvido): relacionar a formula do BMC a codificacao
%% do ESBMC para programas C -- ambos atendidos pelo exemplo abaixo.
Em um programa C, essa codificação parte do código real. Considere um laço que
percorre um \textit{array} de $N$ posições (Listagem~\ref{lst:exemplo-bmc}):

\begin{lstlisting}[caption={Laço em C cuja segurança de acesso o ESBMC verifica.},
  label={lst:exemplo-bmc}, numbers=left, frame=single, language=C]
int soma = 0;
for (int i = 0; i < n; i++)
    soma += v[i];      // acesso a v[i]: vale 0 <= i < N ?
\end{lstlisting}

O vetor \texttt{v} tem \texttt{N} posições, e o laço acessa \texttt{v[i]} para
\texttt{i} de $0$ a \texttt{n}$-1$. Considerando \texttt{n} não negativo, a
condição de segurança é que o número de iterações não ultrapasse o tamanho do
vetor, isto é, $0 \leq n \leq N$; para que a análise seja conclusiva, o limite de
desenrolamento também deve ser suficiente para representar todas as iterações
possíveis. O ESBMC desenrola o laço $k$ vezes, substituindo cada iteração por uma
cópia condicionada a \texttt{i < n}; converte o resultado à forma SSA
(\textit{Static Single Assignment}), em que cada atribuição gera uma nova versão da
variável ($soma_1, soma_2, \dots$); e, para cada acesso \texttt{v[i]}, gera uma
condição de verificação (VCC) de limite ($0 \leq i < N$). As VCCs são submetidas ao
solver SMT: se \texttt{n} puder exceder \texttt{N}, a negação da condição de
segurança pode ser satisfatível, permitindo ao solver produzir uma atribuição que
representa o acesso fora dos limites do vetor; caso contrário, o resultado é UNSAT
até $k$.

\subsection{Tratamento de Laços, Completude e $k$-Indução}

Aplicar BMC a programas exige desenrolar (\textit{unwinding}) os laços até $k$.
Uma \textit{assertion} ao fim de cada laço verifica se ele poderia executar além
de $k$; quando ela falha, a análise daquele laço deixa de ser conclusiva, e é
preciso elevar $k$ ou recorrer à $k$-indução --- foi o que ocorreu com o MAVLink,
verificado com $k = 300$ (Capítulo~\ref{cap:experimentos}). Com a opção
\texttt{--no-unwinding-assertions}, laços acima de $k$ são truncados
silenciosamente~\cite{Gadelha2017}.
%% LC (resolvido): Biere1999 nao sustenta o comportamento da flag do ESBMC;
%% citacao trocada para Gadelha2017.

Sob uma codificação fiel do modelo e com o limite de desenrolamento adequadamente
tratado, o BMC é \textit{sound} e completo até $k$: todo contraexemplo encontrado
representa uma violação da propriedade \emph{no modelo}, e toda violação alcançável
em até $k$ passos é identificada; a ausência de contraexemplos, porém, indica
apenas que nenhuma violação
foi encontrada dentro do limite de exploração, e não demonstra, isoladamente, que
a propriedade se mantém em execuções que ultrapassem esse limite~\cite{Biere1999}. A correspondência entre esse contraexemplo e o código de
produção não é automática --- depende da fidelidade do \textit{harness} ---, e é o
que a etapa de conferência (etapa~(v), Capítulo~\ref{cap:metodologia})
estabelece. Para estender a garantia além de $k$, a $k$-indução~\cite{Sheeran2000}
combina um caso base (a propriedade vale nos primeiros $k$ estados) com um passo
indutivo (se vale em $k$ estados consecutivos quaisquer, vale no seguinte); o
ESBMC a implementa automaticamente~\cite{Gadelha2017}. A propriedade só é
considerada comprovada quando o caso base e o passo indutivo são ambos verificados
com sucesso. Distinta dessa $k$-indução automática, a prova de inalcançabilidade do GPS uBlox apresentada nesta
dissertação (Seção~\ref{subsec:gps-inducao}) é uma indução manual de um único
passo sobre um invariante do laço.
%% LC (resolvido): k-inducao e de Sheeran2000 (nao Biere1999), implementada no
%% ESBMC por Gadelha2017; distincao entre inducao manual de um passo e a k-inducao
%% automatica da ferramenta acrescentada.

\section{Satisfiability Modulo Theories}
\label{sec:smt}

O \textit{Satisfiability Modulo Theories} (SMT) generaliza a satisfatibilidade
booleana (SAT) ao raciocinar sobre teorias matemáticas
específicas~\cite{Barrett2018}. Os solvers SMT combinam procedimentos de decisão
para teorias como aritmética, \textit{arrays} e \textit{bit-vectors} com
algoritmos SAT eficientes, resolvendo fórmulas impraticáveis em lógica
proposicional pura. Entre os solvers atuais estão o Z3~\cite{deMoura2008}, o
Bitwuzla~\cite{Niemetz2023} e o cvc5~\cite{Barbosa2022}; o
Boolector~\cite{Brummayer2009}, especializado em \textit{bit-vectors} e
\textit{arrays}, é o solver usado em todas as verificações deste trabalho
(versão~3.2.4).
%% LC (resolvido): definicao de SMT estava sem citacao -- acrescentada a
%% referencia padrao Barrett2018 (Handbook of Model Checking, 2018).

\subsection{Teorias Relevantes e Vantagens sobre SAT}

A teoria de \textit{bit-vectors} modela inteiros de largura fixa, capturando a
semântica de tipos como \texttt{int32\_t} e \texttt{uint8\_t}, inclusive
\textit{overflow} e \textit{underflow}. A de \textit{arrays} representa estruturas
indexadas por operações de leitura (\textit{select}) e escrita (\textit{store}),
essencial para modelar memória e \textit{buffers}. Para dados de sensores, a
teoria de ponto flutuante implementa a aritmética IEEE-754, com valores especiais
como NaN e infinito. O Boolector, porém, não possui teoria nativa de ponto
flutuante: o ESBMC contorna isso codificando cada operação IEEE-754 em operações
de \textit{bit-vector} (\textit{bit-blasting}), de modo que solvers restritos a
\textit{bit-vectors}, como o Boolector, verifiquem programas com ponto
flutuante~\cite{Gadelha2020}.
%% LC (resolvido): Boolector nao tem teoria nativa de ponto flutuante --
%% explicada a codificacao por bit-blasting do ESBMC (Gadelha2020).

Frente à codificação direta em SAT, os solvers SMT são mais expressivos ---
representam restrições como $x + y > 10$ sem codificação bit a bit --- e exploram
a estrutura matemática das teorias, ganhando eficiência~\cite{deMoura2008}. No
PX4, essa expressividade é essencial: os módulos operam com ponto flutuante
IEEE-754, inteiros sem sinal sujeitos a \textit{underflow} e \textit{arrays} de
tamanho fixo para \textit{buffers} FIFO --- classes que os solvers SMT tratam com
precisão (o ponto flutuante, no caso do Boolector, pela codificação em
\textit{bit-vectors} descrita acima).

\section{ESBMC}

O ESBMC (\textit{Efficient SMT-Based Context-Bounded Model Checker}) verifica
programas em C e C++, com \textit{front-ends} adicionais para
Python~\cite{Farias2024}, Kotlin~\cite{Menezes2022} e
Solidity~\cite{Menezes2024}. Mantido por uma colaboração internacional que inclui
pesquisadores da Universidade Federal do Amazonas, adota uma abordagem baseada em
SMT, combinando a eficiência dos procedimentos SAT com o raciocínio sobre teorias
matemáticas --- aritmética, \textit{bit-vectors} e \textit{arrays} ---, o que dá
suporte mais preciso à verificação de aritmética e estruturas de dados
complexas~\cite{Cordeiro2012}.
%% LC (resolvido): sigla ESBMC uniformizada para "Efficient SMT-Based
%% Context-Bounded Model Checker" nos quatro lugares do documento.

\subsection{Pipeline de Verificação}

O fluxo do ESBMC tem quatro estágios (Figura~\ref{fig:esbmc_pipeline}). O
código-fonte passa pelo \textit{front-end} da linguagem --- Clang para C/C++,
Soot para Kotlin, CPython para Python ---, gerando uma representação intermediária
(programa GOTO). A execução simbólica converte o GOTO em fórmulas que representam
as execuções até a profundidade $k$, das quais se geram as condições de
verificação (VCCs). As VCCs, após simplificação, vão ao solver SMT (aqui, o
Boolector~3.2.4): SAT indica um contraexemplo; UNSAT, que a propriedade vale até
$k$~\cite{Menezes2024}.

\begin{figure}[ht]
  \centering
  \resizebox{\textwidth}{!}{%
  \begin{tikzpicture}[
    font=\linespread{1}\footnotesize, >={Stealth[length=2mm]},
    bloco/.style={draw=palEdge, fill=palFill, rounded corners=2pt, align=center,
      inner sep=4pt, minimum height=9mm},
    fe/.style={draw=palEdge, fill=palFill, rounded corners=1.5pt, align=center,
      text width=2.3cm, inner sep=2pt, minimum height=8mm},
    sat/.style={draw=palEdgeDark, line width=0.8pt, fill=palStrong, rounded corners=2pt, align=center, text width=2.6cm, inner sep=3pt},
    unsat/.style={draw=palEdge, fill=palPaper, rounded corners=2pt, align=center, text width=2.6cm, inner sep=3pt}]
    \node[fe] (c) {\textbf{C/C++}\\Clang (AST)};
    \node[fe, below=1.5mm of c] (k) {\textbf{Java/Kotlin}\\Soot (Jimple)};
    \node[fe, below=1.5mm of k] (s) {\textbf{Solidity}\\gramática};
    \node[fe, below=1.5mm of s] (p) {\textbf{Python}\\AST do CPython};
    \node[above=1mm of c, font=\small\bfseries] (fet) {\textit{Front-ends}};
    \begin{scope}[on background layer]
      \node[draw=palEdge, fill=palPaper, rounded corners=3pt, fit=(fet)(c)(p), inner sep=3pt] (fes) {};
    \end{scope}
    \node[bloco, right=9mm of fes, text width=2.5cm] (goto) {\textbf{Programa GOTO}\\representação intermediária};
    \node[bloco, right=7mm of goto, text width=2.7cm] (sym) {\textbf{Execução simbólica}\\BMC, $k$-indução,\\interpretação abstrata,\\\textit{slicing}};
    \node[bloco, right=7mm of sym, text width=2.6cm] (vcc) {\textbf{Condições de verificação}\\(VCCs)};
    \node[bloco, above=8mm of vcc, text width=4.2cm, fill=palStrong, draw=palEdgeDark] (prop) {\textbf{Propriedades}\\limites de arranjos, ponteiros,\\\textit{overflow}, divisão por zero,\\\textit{assertions} do usuário};
    \node[bloco, right=7mm of vcc, text width=2.2cm] (smt) {\textbf{Solver SMT}\\Boolector, Z3,\\Bitwuzla, cvc5,\\MathSAT, Yices};
    \node[sat, right=9mm of smt, yshift=9mm] (sa) {\textbf{SAT}\\contraexemplo\\(violação no modelo)};
    \node[unsat, right=9mm of smt, yshift=-9mm] (un) {\textbf{UNSAT}\\propriedade válida\\até o limite $k$};
    \draw[->] (fes.east |- goto) -- (goto);
    \draw[->] (goto) -- (sym);
    \draw[->] (sym) -- (vcc);
    \draw[->] (prop) -- (vcc);
    \draw[->] (vcc) -- (smt);
    \draw[->] (smt.east) -- ++(4mm,0) |- (sa.west);
    \draw[->] (smt.east) -- ++(4mm,0) |- (un.west);
  \end{tikzpicture}}
  \caption{Pipeline de verificação do ESBMC: o código-fonte é processado pelos
           \textit{front-ends} correspondentes, transformado em programa GOTO e submetido
           à execução simbólica para geração das condições de verificação, que
           são resolvidas por um solver SMT: SAT indica um contraexemplo
           (violação no modelo); UNSAT, que a propriedade vale até o limite $k$.
           Elaborado pela autora com base em~\cite{ESBMC_GitHub}.}
  \label{fig:esbmc_pipeline}
\end{figure}

\subsection{Classes de Propriedades Verificadas}

O ESBMC é capaz de detectar \textit{buffer overflows}, desreferência de ponteiros
nulos, \textit{use-after-free} e vazamentos de memória; \textit{overflow} e
\textit{underflow} de inteiros; operações de ponto flutuante que produzem NaN ou
infinito; e, em programas concorrentes, \textit{data races} e
\textit{deadlocks}~\cite{Menezes2024}. Cada classe é ativada por opções próprias
da ferramenta; nesta pesquisa foram habilitadas as verificações de segurança de
memória (limites de \textit{arrays} e validade de ponteiros) e de \textit{overflow}
aritmético, conforme a configuração detalhada no Capítulo~\ref{cap:experimentos}.
O usuário pode ainda inserir \textit{assertions} próprias no código.

%% LC (resolvido): lacuna de fundamentacao -- acrescentada subsecao sobre a CWE
%% (relacao com CVE, hierarquia, e as entradas usadas no Cap. de resultados).
\subsection{Classificação de Vulnerabilidades: a CWE}

A \textit{Common Weakness Enumeration} (CWE) é uma taxonomia mantida pelo MITRE
que cataloga tipos de fraquezas de software --- padrões recorrentes de erro que
podem originar vulnerabilidades --- independentemente de qualquer ocorrência
específica~\cite{MITRE_CWE}. Ela se distingue da \textit{Common Vulnerabilities
and Exposures} (CVE), que identifica instâncias públicas e concretas em produtos
específicos e costuma referenciar uma ou mais entradas CWE para indicar a causa
raiz da falha. A CWE organiza suas entradas em uma hierarquia de generalidade, na
qual uma fraqueza mais específica é filha de uma mais abstrata; por exemplo,
CWE-787 (\textit{Out-of-bounds Write}) é filha de CWE-119 (\textit{Improper
Restriction of Operations within the Bounds of a Memory Buffer})~\cite{MITRE_CWE}.

Esta pesquisa emprega seis entradas da CWE para classificar os achados do
Capítulo~\ref{cap:resultados}: CWE-20 (\textit{Improper Input Validation}),
CWE-1284 (\textit{Improper Validation of Specified Quantity in Input}), subclasse
de CWE-20, CWE-119 (\textit{Improper Restriction of Operations within the Bounds of
a Memory Buffer}), CWE-190 (\textit{Integer Overflow or Wraparound}), CWE-191
(\textit{Integer Underflow (Wrap or Wraparound)}) e CWE-787 (\textit{Out-of-bounds
Write})~\cite{MITRE_CWE}. A nomenclatura da CWE-191 usa ``\textit{underflow}'' num
sentido específico de segurança --- uma subtração sem sinal que dá a volta para um
valor grande ---, distinto do uso em análise numérica (resultado pequeno demais
para ser representado); essa distinção é retomada no
Capítulo~\ref{cap:resultados} ao classificar o achado no GPS uBlox.

\subsection{Extensões Recentes}

Trabalhos recentes ampliaram o ESBMC ou o empregaram em novos domínios. Entre as
extensões da própria ferramenta, Monteiro et al. ampliaram o suporte a C++, com
\textit{templates} e herança virtual~\cite{Monteiro2022} --- pertinente ao PX4,
escrito em C++17. Entre os desenvolvimentos e aplicações relacionados, o FuSeBMC
combina verificação formal e geração de testes por \textit{fuzzing}, usando o ESBMC
como um de seus motores~\cite{Alshmrany2022}, e Wu et al. aplicaram o ESBMC a
componentes do Arm Confidential Computing~\cite{Wu2024}.

%% LC (pendente -- resolver ao editar o Cap. 3): o paragrafo de Extensoes Recentes
%% se repete no Cap. 3; manter uma ocorrencia e usar referencia cruzada.

\section{PX4 Autopilot}

O PX4 Autopilot é uma plataforma de software \textit{open-source} para controle de
veículos aéreos não tripulados, sob licença BSD~\cite{PX4_GitHub,
PX4_Documentation}. Suporta diversas configurações de aeronave --- multirotores,
asas fixas, VTOL e plataformas híbridas --- e é amplamente utilizado em aplicações
comerciais, científicas e recreativas.

\subsection{Arquitetura de Software}

A arquitetura do PX4 segue o padrão \textit{publish-subscribe}: os módulos
comunicam-se pelo \textit{middleware} uORB (\textit{micro Object Request Broker}),
que desacopla os componentes e facilita o desenvolvimento modular e a
manutenção~\cite{PX4_Documentation}. Os principais subsistemas incluem os
\textit{drivers} de sensores (GPS, IMU, magnetômetro e barômetro), a estimação de
estado por filtros de Kalman estendidos, os controladores de atitude e posição, a
navegação com planejamento de trajetória e os mecanismos de \textit{failsafe},
além da interface de comunicação MAVLink~\cite{PX4_Documentation}. A
Figura~\ref{fig:px4_arquitetura} resume essa organização.

\begin{figure}[ht]
  \centering
  \resizebox{0.95\textwidth}{!}{%
  \begin{tikzpicture}[
    font=\footnotesize, >={Stealth[length=2mm]},
    verif/.style={draw=palEdgeDark, fill=palStrong, line width=0.9pt, rounded corners=3pt,
      align=center, text width=3cm, minimum height=11mm, inner sep=3pt},
    plano/.style={draw=palEdge, fill=palPaper, rounded corners=3pt,
      align=center, text width=3cm, minimum height=11mm, inner sep=3pt},
    tabp/.style={draw=palEdge, fill=palFill, rounded corners=2pt,
      align=center, text width=2.1cm, minimum height=4.5mm, inner sep=1pt,
      font=\scriptsize\bfseries},
    tabs/.style={draw=palEdge, fill=palFill, rounded corners=2pt,
      align=center, text width=2.1cm, minimum height=4.5mm, inner sep=1pt,
      font=\scriptsize\bfseries},
    bus/.style={draw=palEdgeDark, fill=palFill, line width=1pt, rounded corners=4pt,
      align=center, text width=2.6cm, minimum height=18mm},
    grp/.style={draw=palEdge, dashed, rounded corners=5pt, inner sep=6pt}]

    \node[verif] (imu) {\textbf{IMU Driver}\\(BMI088)};
    \node[verif, below=10mm of imu] (gps) {\textbf{GPS Driver}\\(uBlox)};
    \node[verif, below=10mm of gps] (rc) {\textbf{RC Input / Expo}};
    \node[tabp, above=1.5mm of imu] (timu) {Publisher};
    \node[tabp, above=1.5mm of gps] (tgps) {Publisher};
    \node[tabp, above=1.5mm of rc]  (trc)  {Publisher};

    \node[bus, right=30mm of gps] (uorb) {\textbf{uORB}\\barramento de\\mensagens};

    \node[plano, right=30mm of uorb] (fc) {\textbf{Controle de voo}};
    \node[plano, above=10mm of fc]  (ekf) {\textbf{EKF / Estimador de estado}};
    \node[verif, below=10mm of fc]  (mav) {\textbf{MAVLink}};
    \node[tabs, above=1.5mm of ekf] (tekf) {Subscriber};
    \node[tabs, above=1.5mm of fc]  (tfc)  {Subscriber};
    \node[tabs, above=1.5mm of mav] (tmav) {Subscriber};

    \begin{scope}[on background layer]
      \node[grp, fit=(timu)(imu)(trc)(rc)]   (gpub) {};
      \node[grp, fit=(tekf)(ekf)(tmav)(mav)] (gsub) {};
    \end{scope}
    \node[font=\bfseries, above=1.5mm of gpub] {Publishers};
    \node[font=\bfseries, above=1.5mm of gsub] {Subscribers};

    \draw[->] (imu.east) -- node[above, font=\scriptsize\itshape] {orb\_publish()} (uorb.west);
    \draw[->] (gps.east) -- (uorb.west);
    \draw[->] (rc.east)  -- (uorb.west);
    \draw[->] (uorb.east) -- (ekf.west);
    \draw[->] (uorb.east) -- (fc.west);
    \draw[->] (uorb.east) -- (mav.west);

    \node[verif, text width=5.2cm, minimum height=7mm, below=26mm of uorb] (leg)
      {\textbf{Exemplos de módulos verificados}};
  \end{tikzpicture}}
  \caption{Arquitetura de comunicação do PX4 Autopilot segundo o padrão
    \textit{publish-subscribe}: \textit{drivers} de sensores (IMU BMI088, GPS uBlox)
    e módulos de entrada e telemetria (Expo RC, MAVLink) publicam e assinam tópicos
    por meio do \textit{middleware} uORB. Com realce (tom mais escuro e borda mais
    espessa), exemplos de módulos verificados neste trabalho; a figura é
    ilustrativa e não representa a
    totalidade dos módulos analisados (a amostra completa consta do
    Capítulo~\ref{cap:metodologia}). Elaborado pela autora com base
    em~\cite{PX4_Documentation}.}
  \label{fig:px4_arquitetura}
\end{figure}

\section{Ambiente de Execução}
\label{sec:ambiente-execucao}

O PX4 roda sobre o NuttX, sistema operacional de tempo real compatível com POSIX
e voltado a microcontroladores de recursos limitados~\cite{NuttX}, com
escalonamento preemptivo e primitivas de sincronização. Os alvos considerados
neste trabalho, baseados em microcontroladores STM32, dispõem de uma unidade de
proteção de memória (MPU), mas essa proteção só é ativada nos modos
\textit{protected} ou \textit{kernel} do NuttX; como foram compilados no modo
\textit{flat}, não há isolamento entre os módulos --- todos compartilham o mesmo
espaço de endereçamento, de modo que um defeito de memória em qualquer módulo pode
comprometer o estado global~\cite{NuttX}.
%% LC (resolvido): STM32 possui MPU -- esclarecido que ela so e usada nos modos
%% protected/kernel; nos alvos deste trabalho (modo flat) nao ha isolamento.

\subsection{Desafios para Verificação Formal}

O PX4 reúne características que tornam sua verificação ao mesmo tempo desafiadora e
necessária. O uso extensivo de ponto flutuante em navegação e controle exige
modelagem precisa da semântica IEEE-754, sob risco de não detectar defeitos
ligados a NaN ou a imprecisões numéricas; os laços de controle executam a taxas da
ordem de centenas de hertz, impondo restrições temporais rigorosas; e a interação
com hardware sujeito a falhas --- sensores que retornam dados corrompidos,
\textit{timeouts} de comunicação --- exige robustez adicional do
software~\cite{Wanner2024}.

Defeitos em sistemas de controle de voo podem resultar na perda do equipamento
ou, em cenários graves, em danos a terceiros~\cite{Siddiqi2022, Mekdad2023}. A
certificação de software aeronáutico apoia-se tradicionalmente em atividades de
teste (DO-178C); o suplemento de métodos formais DO-333 reconhece a verificação
formal como meio de conformidade: conforme os objetivos e critérios aplicáveis, a
análise formal pode cumprir parte dos objetivos de verificação e, em determinados
casos, substituir revisões, análises ou testes específicos~\cite{RTCA_DO178C,
RTCA_DO333}. Trabalhos
anteriores em sistemas embarcados \textit{safety-critical} estabeleceram
metodologias aplicáveis a VANTs~\cite{Mzid2024}, e a combinação de Model Checking
com análise de código embarcado tem se mostrado eficaz na detecção de defeitos que
escapam a abordagens baseadas apenas em testes~\cite{Khairullah2022, Sousa2026}.
%% LC (resolvido): alegacao de certificacao reformulada com o suplemento de
%% metodos formais DO-333 (nao substitui testes automaticamente). LC (resolvido):
%% Chatterjee2022 nao sustentava a frase -- trocado por Sousa2026.

\section{Resumo}

Este capítulo reuniu os fundamentos do trabalho. A verificação formal foi
apresentada como complemento aos testes, com foco em propriedades de
\textit{safety}, cuja violação, no domínio de voo, pode ser catastrófica. Entre as
técnicas, adotou-se o Model Checking, pela automação e pela geração de
contraexemplos; o \textit{Bounded Model Checking}, por limitar a busca a $k$
passos, é \textit{sound} e completo apenas até esse limite --- restrição atenuada
pela $k$-indução, aplicada adiante ao GPS uBlox. A expressividade dos solvers SMT
(\textit{bit-vectors}, \textit{arrays} e IEEE-754) é o que viabiliza verificar
firmware embarcado com precisão, e o ESBMC consolida essas capacidades. Por fim, a
arquitetura modular do PX4, articulada pelo \textit{middleware} uORB, é o que
torna viável a verificação isolada de cada módulo --- premissa central da
metodologia desenvolvida nos próximos capítulos.

\chapter{Trabalhos Relacionados}
\label{cap:relacionados}

Este capítulo apresenta trabalhos que aplicaram verificação formal a sistemas
embarcados e a software de veículos autônomos, posicionando esta pesquisa em
relação ao estado da arte.

\section{Protocolo de Levantamento}
\label{sec:rel-protocolo}

O levantamento foi exploratório, e não uma revisão sistemática: o objetivo foi
posicionar a pesquisa, não mapear exaustivamente o campo. As bases consultadas
foram IEEE Xplore, ACM Digital Library e Google Scholar --- esta usada sobretudo
para rastrear citações a partir dos trabalhos já identificados. As buscas
combinaram termos de três eixos: a técnica (\textit{bounded model checking},
\textit{formal verification}), o alvo (\textit{embedded firmware}, \textit{UAV},
\textit{autopilot}) e a ferramenta (ESBMC, CBMC, Frama-C), privilegiando
publicações a partir de 2019, exceto os trabalhos seminais da área.

Foram incluídos trabalhos que aplicam verificação formal a código C/C++, que
tratam de sistemas embarcados críticos ou de veículos autônomos, ou que levantam
o estado da arte nesses domínios. Excluíram-se os restritos a verificação em
nível de modelo sem correspondência com código executável e os que tratam de
segurança de UAVs sob perspectiva exclusivamente de rede ou ataque físico, sem
relação com o firmware embarcado.
%% LC (acknowledged): protocolo exploratorio; nao ha registro de strings/datas/
%% contagens -- retirada a promessa de reproduzir a busca.

\section{Verificação Formal de Sistemas Embarcados}
\label{sec:rel-embarcados}

Ferramentas de BMC como CBMC~\cite{Kroening2014} e ESBMC~\cite{Menezes2024}
permitem analisar programas C/C++ pela exploração simbólica de seus caminhos de
execução. O CBMC foi uma das primeiras a aplicar BMC diretamente a programas C,
adotando uma abordagem baseada em SAT; o ESBMC adota uma abordagem baseada em
SMT~\cite{deMoura2008}, que incorpora teorias como vetores de bits,
\textit{arrays} e ponto flutuante. Aplicações recentes demonstram o uso dessas
ferramentas na identificação de defeitos em projetos de código aberto de grande
escala~\cite{Sousa2026}. As duas são usadas nesta pesquisa, a segunda na
validação cruzada.

\textit{Surveys} recentes~\cite{Beyer2022, Beyer2024} indicam que boa parte dos
trabalhos de verificação ainda se concentra em programas sintéticos ou
\textit{benchmarks} padronizados, e não em sistemas reais de produção --- lacuna
que motiva diretamente este trabalho. Cabe a ressalva de que essa observação
decorre de um levantamento exploratório, não de uma medição sistemática.
%% LC (resolvido): afirmacao de lacuna atenuada -- nao atribuida aos relatorios
%% de competicao, que nao a sustentam.

\section{Aplicações Práticas de Verificação Formal}
\label{sec:rel-praticas}

Sousa et al.~\cite{Sousa2026} aplicaram o ESBMC a projetos \textit{open-source}
em C de grande escala, com uma metodologia sistemática de pré-processamento
automatizado e priorização por função para guiar o verificador nas partes mais
críticas, identificando defeitos previamente desconhecidos. É a abordagem mais
próxima desta pesquisa; a diferença é que os projetos de Sousa et al. são
aplicações genéricas, enquanto este estudo foca em firmware de controle de voo
com restrições de tempo real --- tratadas aqui pela seleção de módulos
verificáveis isoladamente, com \textit{stubs} que abstraem dependências de
hardware e temporização, sem comprometer a análise das propriedades de
\textit{safety} internas a cada módulo.
%% LC (resolvido): Sousa2023 era o preprint de Sousa2026 -- citacoes unificadas
%% em Sousa2026.

Monteiro et al.~\cite{Monteiro2022} estenderam o ESBMC para C++, com suporte a
\textit{templates} e herança virtual --- diretamente aplicável ao PX4, escrito em
C++17. Wu et al.~\cite{Wu2024} verificaram componentes de arquiteturas Arm
Confidential Computing, demonstrando a escalabilidade da ferramenta em sistemas
complexos. Em outra direção, Menezes et al.~\cite{Menezes2025GCSE}, em
colaboração com a Intel, propuseram a técnica VO-GCSE (\textit{Verification
Optimization through Global Common Subexpression Elimination}), que reduz o tempo
de verificação ao eliminar subexpressões redundantes em modelos simbólicos ---
artigo da trilha de demonstrações do FSE 2025. A otimização não foi empregada
nesta pesquisa, mas é relevante para módulos de firmware com estruturas
repetitivas como as do PX4. Na fronteira entre verificação e teste, Alshmrany et
al.~\cite{Alshmrany2022} propuseram o FuSeBMC, que combina BMC e \textit{fuzzing}
para ampliar a cobertura na geração automática de casos de teste.
%% LC (resolvido): informado que a otimizacao VO-GCSE nao foi usada neste
%% trabalho. LC (resolvido): FuSeBMC retirado deste paragrafo (discutido na
%% subsecao de abordagens nao-BMC, como desenvolvimento relacionado).

\section{Verificação de Sistemas Autônomos}
\label{sec:rel-autonomos}

Wanner et al.~\cite{Wanner2024} revisam a certificação e a confiabilidade de
sistemas aviônicos de UAVs, destacando a necessidade de verificação rigorosa, mas
sem aplicá-la ao código-fonte. Siddiqi et al.~\cite{Siddiqi2022} e Mekdad et
al.~\cite{Mekdad2023} identificam fragilidades em UAVs sob a perspectiva de
ataques externos e de rede, sem abordar o firmware embarcado.
Luckcuck et al.~\cite{Luckcuck2019} levantam, em \textit{survey} sistemático, os
formalismos e desafios da verificação de sistemas robóticos autônomos,
evidenciando que a verificação de firmware de autopiloto permanece pouco
explorada.

Taylor et al.~\cite{Taylor2021} propõem o Avis, um \textit{model checking}
\textit{in-situ} para VANTs que executa o firmware em simulação e injeta falhas
de sensores para observar o comportamento do veículo. É o trabalho mais próximo
quanto ao alvo, porém verifica o sistema em execução, e não o código-fonte por
análise estática, como o BMC aqui adotado. Bernardeschi et
al.~\cite{Bernardeschi2023} combinam verificação formal e co-simulação no
controle cooperativo de drones, com especificações em lógica temporal, mas não
tratam de defeitos no nível dos \textit{drivers}. Pollien et
al.~\cite{Pollien2021} verificam com Frama-C a biblioteca matemática do
autopiloto \textit{open-source} Paparazzi, provando ausência de erros de execução
em funções de conversão de ponto flutuante; a tese que consolida o
trabalho~\cite{Pollien2023} combina Frama-C e o assistente de provas Coq. Trata-se
de verificação dedutiva, que exige anotações manuais extensas, diferentemente do
BMC automatizado adotado aqui.

Di Sorbo et al.~\cite{deSorbo2022} caracterizam automaticamente problemas de
\textit{safety} reportados em plataformas de software para UAVs, revelando
padrões recorrentes de defeitos. Mzid~\cite{Mzid2024} propõe padrões de projeto
para sistemas embarcados de tempo real \textit{safety-critical}, sem aplicação
empírica a controle de voo. Khatiri et al.~\cite{Khatiri2024} propõem o
Aerialist, que automatiza a geração e a execução de cenários de teste no PX4 por
simulação; embora focado em teste, complementa esta pesquisa ao reforçar o valor
de abordagens automatizadas de garantia de qualidade em firmware de autopiloto.

\subsection{Abordagens não Baseadas em \textit{Model Checking}}
\label{subsec:rel-nao-bmc}

As seções anteriores discutem, em sua maior parte, a própria família de técnicas
desta pesquisa. Convém situá-la também em relação às alternativas, porque a
escolha do BMC implica renúncias que só se tornam visíveis por contraste.

\textbf{Verificação dedutiva e análise estática \textit{sound}.} A verificação
dedutiva estabelece a correção por prova, e não por exploração de estados --- é a
abordagem de Pollien et al.~\cite{Pollien2021} sobre o Paparazzi. Quando as
especificações e as hipóteses do modelo são adequadas, essa verificação pode
estabelecer propriedades para além de um limite finito de execução, diferentemente
do BMC convencional, cuja análise está condicionada à profundidade de
desenrolamento; em contrapartida, exige anotar o código com pré-condições,
pós-condições e invariantes; em uma
base de cerca de 710 mil linhas, escrita por centenas de colaboradores e sem
anotações, esse custo é proibitivo. A interpretação abstrata é \textit{sound} por
construção e escala bem, mas produz alarmes que o analista descarta um a um, sem
o contraexemplo concreto que o BMC entrega --- justamente o que permite a
conferência contra o código de produção descrita na
Seção~\ref{subsec:etapa-replay}. O Astrée, de Cousot et
al.~\cite{Cousot2005Astree}, é o representante mais conhecido dessa família, com
histórico em código aeronáutico de produção, como o da Airbus~\cite{Bertrane2010};
a ferramenta comercial Polyspace segue a mesma linha. Já a linguagem SPARK, um
subconjunto de Ada, permite provar estaticamente a ausência de erros de execução
na compilação~\cite{McCormick2015}. As três compartilham o custo estrutural que
motiva a escolha do BMC: exigem disciplina de anotação ou de linguagem que uma
base pré-existente e não anotada, como a do PX4, não satisfaz.

\textbf{Síntese de código verificado.} Em vez de verificar código existente, o
programa HACMS da DARPA construiu o firmware de VANTs a partir de componentes
concebidos para serem verificados --- entre eles o microkernel seL4 e o
compilador CompCert ---, produzindo o autopiloto SMACCMPilot, que resistiu a um
\textit{red team} com seis semanas e acesso total ao código~\cite{Fisher2017}. O
contraste com esta pesquisa é nítido: o HACMS parte de uma base construída para a
verificação, enquanto o objeto deste trabalho é um firmware de produção
pré-existente e não anotado, sobre o qual o BMC é aplicável sem reescrita.

\textbf{Teste dirigido a VANT.} No extremo oposto estão as técnicas dinâmicas.
Kim et al.~\cite{Kim2021PGFuzz} propõem o PGFUZZ, \textit{fuzzing} guiado por
política aplicado a ArduPilot, PX4 e Paparazzi. O RVFuzzer, de T.~Kim et
al.~\cite{Kim2019RVFuzzer}, usa realimentação do controlador para gerar valores
de parâmetros --- enviados pela estação de solo --- que levam o veículo a estados
inseguros, encontrando defeitos de validação de entrada em autopilotos de
produção, incluindo o PX4. Han et al.~\cite{Han2024RangeFuzz} detectam defeitos de
especificação de faixa em módulos de controle de voo. Os três encontram defeitos
reais, por execução, não por análise. A relação com este trabalho é de
complementaridade: técnicas dinâmicas exercitam o sistema integrado, com
escalonador e hardware presentes, mas cobrem apenas as execuções que ocorrem; o
BMC explora exaustivamente os caminhos representados no modelo até o limite de
desenrolamento, sobre um recorte abstraído por \textit{stubs}. O argumento é desenvolvido na
Seção~\ref{sec:discussao-resultados}.

\textbf{Contexto normativo.} Na aviação certificada, a DO-178C rege a garantia de
software aeronáutico crítico~\cite{RTCA_DO178C}, e seu suplemento DO-333
estabelece as condições sob as quais métodos formais podem substituir, total ou
parcialmente, atividades de teste exigidas pela norma~\cite{RTCA_DO333}. É nesse
regime que Astrée, Polyspace e SPARK acumularam histórico de uso. O PX4, por não
ser certificado sob a DO-178C, está fora desse arcabouço: a verificação aqui
conduzida não substitui nem satisfaz requisito de certificação algum, e seus
resultados são evidência técnica independente, não conformidade normativa.

\section{Posicionamento deste Trabalho}
\label{sec:rel-posicionamento}

A Tabela~\ref{tab:trabalhos-relacionados} sintetiza os principais trabalhos
relacionados quanto à técnica, ao sistema-alvo, ao tipo de estudo, ao objeto
avaliado e à ferramenta.

\begin{table}[p]
\centering
\caption{Síntese comparativa dos trabalhos relacionados em verificação formal}
\label{tab:trabalhos-relacionados}
\scriptsize
\renewcommand{\arraystretch}{0.92}
\setlength{\tabcolsep}{2pt}
\begin{tabular}{|p{2.3cm}|p{2.0cm}|p{2.2cm}|p{1.7cm}|p{2.1cm}|p{1.9cm}|}
\hline
\textbf{Trabalho} & \textbf{Técnica} & \textbf{Sistema-alvo} &
\textbf{Tipo de estudo} & \textbf{Objeto} & \textbf{Ferramenta} \\
\hline
Beyer (2022)        & Survey MC       & Múltiplas linguagens    & Survey     & ---                   & Vários    \\ \hline
Beyer (2024)        & SV-COMP         & Benchmark C             & Competição & Benchmark             & Vários    \\ \hline
Kroening (2014)     & BMC             & Programas C             & Ferramenta & Benchmark             & CBMC      \\ \hline
Menezes (2024)      & BMC             & Programas C             & Ferramenta & Benchmark             & ESBMC     \\ \hline
Sousa (2026)        & BMC             & Open-source C           & Empírico   & Sistema de produção   & ESBMC     \\ \hline
Monteiro (2022)     & BMC             & Programas C++           & Ferramenta & Código sintético      & ESBMC     \\ \hline
Wu (2024)           & BMC             & Arm CCA                 & Empírico   & Sistema de produção   & ESBMC     \\ \hline
Alshmrany (2022)    & BMC + Testes    & Programas C             & Ferramenta & Benchmark             & FuSeBMC   \\ \hline
Menezes (2025)      & BMC + GCSE      & Programas C             & Demonstração & Benchmark e indústria & ESBMC   \\ \hline
Fisher (2017)       & Síntese verif.  & \textbf{UAV SMACCMPilot}& Empírico   & Sistema de produção   & seL4/CompCert \\ \hline
Pollien (2021)      & Dedutivo        & UAV Paparazzi           & Empírico   & Sistema de produção   & Frama-C   \\ \hline
Cousot (2005)       & Interp. abstrata& Código aeroespacial     & Ferramenta & Sistema de produção   & Astrée    \\ \hline
McCormick (2015)    & Dedutivo/ling.  & Ada embarcado           & Conceitual & ---                   & SPARK     \\ \hline
Mzid (2024)         & Padrões         & Sistemas RT críticos    & Conceitual & ---                   & --        \\ \hline
Taylor (2021)       & MC in-situ      & \textbf{UAV runtime}    & Empírico   & Sistema de produção   & Avis      \\ \hline
Bernardeschi (2023) & FV + co-sim     & \textbf{Drones cooper.} & Empírico   & Modelo e co-simulação & NuSMV     \\ \hline
Luckcuck (2019)     & Survey FV       & Robótica autônoma       & Survey     & ---                   & Vários    \\ \hline
Wanner (2024)       & Survey          & \textbf{UAV aviônico}   & Survey     & ---                   & --        \\ \hline
Siddiqi (2022)      & Survey          & \textbf{UAV segurança}  & Survey     & ---                   & --        \\ \hline
Mekdad (2023)       & Survey          & \textbf{UAV redes}      & Survey     & ---                   & --        \\ \hline
Kim (2021)          & Fuzzing         & \textbf{PX4, ArduPilot} & Empírico   & Sistema de produção   & PGFUZZ    \\ \hline
Kim (2019)          & Fuzzing         & \textbf{PX4, ArduPilot} & Empírico   & Sistema de produção   & RVFuzzer  \\ \hline
Han (2024)          & Fuzzing         & \textbf{Controle de voo}& Empírico   & Sistema de produção   & ---       \\ \hline
Khatiri (2024)      & Simulação       & \textbf{PX4 VANTs}      & Ferramenta & Simulação             & Aerialist \\ \hline
Di Sorbo (2022)     & Análise SW      & \textbf{UAV software}   & Empírico   & Relatos de defeito    & --        \\ \hline
\textbf{Este trabalho} & \textbf{BMC} & \textbf{PX4 firmware}   & \textbf{Empírico} & \textbf{Código-fonte (1 VAM)$^\dagger$} & \textbf{ESBMC (+ CBMC)} \\ \hline
\end{tabular}

\scriptsize
\begin{flushleft}
$^\dagger$VAM = violação alcançável no modelo, na acepção fixada no
Capítulo~\ref{cap:introducao}. Não se reivindica defeito confirmado: a análise do
driver de produção mostrou o estado violador inalcançável
(Seção~\ref{subsec:gps-inducao}). Ver Tabela~\ref{tab:resultados}.
\end{flushleft}
\end{table}

Este trabalho diferencia-se por aplicar BMC diretamente ao código-fonte de
\textit{drivers} de um sistema de produção real e amplamente utilizado (PX4
Autopilot, ramo \textit{main}, \textit{commit} \texttt{8be2590c15}). Enquanto
Taylor et al.~\cite{Taylor2021} verificam o comportamento em execução e
Bernardeschi et al.~\cite{Bernardeschi2023} focam a lógica de controle
cooperativo, este estudo verifica \textit{drivers} críticos no nível do
código-fonte, identificando candidatos a defeito de memória e aritmética em
caminhos representados no modelo de verificação, cuja exploração não é garantida
pelos testes convencionais do projeto. Em relação ao
fluxo de BMC de Sousa et al.~\cite{Sousa2026}, do qual herda as etapas gerais, a
contribuição metodológica está em dois passos próprios de firmware: a abstração
da E/S de hardware por \textit{stubs} e a especificação de \textit{assertions}
derivadas de \textit{datasheets} e do código-fonte.

Diferentemente dos \textit{surveys} sobre segurança de
UAVs~\cite{Wanner2024, Siddiqi2022, Mekdad2023}, este trabalho oferece evidência
empírica da \textbf{aplicabilidade} da verificação formal a firmware de drones em
produção: uma violação alcançável no modelo (GPS uBlox), reportada à comunidade
PX4 (Issue~\#26866), proposta como endurecimento defensivo (PR~\#27576) e mostrada
inalcançável no driver de produção, e dois candidatos triados, respectivamente,
como especificação disputada (temperatura do BMI088) e falso positivo (FIFO do
BMI088), com o procedimento de triagem documentado. A aplicação ao SAGRES e ao
monitoramento ambiental na Amazônia é a motivação do estudo de caso, e não um
fator que altere o método.
%% LC (resolvido): retirada a sugestao de que o monitoramento ambiental influencia
%% o metodo -- apresentado como motivacao, nao como fator metodologico.

\section{Resumo}
\label{sec:rel-resumo}

Este capítulo posicionou a pesquisa no estado da arte. Ferramentas de BMC como
ESBMC e CBMC têm eficácia crescente na detecção de defeitos em código C/C++ de
produção, mas a maior parte dos trabalhos ainda se concentra em \textit{benchmarks}
sintéticos ou aplicações genéricas. Entre os sistemas autônomos, as abordagens
existentes tratam de verificação em execução, co-simulação ou comportamento
emergente, sem alcançar o nível dos \textit{drivers} de firmware. As alternativas
ao \textit{model checking} --- verificação dedutiva, análise estática
\textit{sound}, síntese de código verificado e teste dirigido --- foram discutidas
na Seção~\ref{subsec:rel-nao-bmc}, e sua relação com esta pesquisa é de
complementaridade.

O trabalho de Pollien et al.~\cite{Pollien2021} sobre o Paparazzi é o precedente
mais próximo em verificação formal de código de autopiloto, mas usa técnicas
dedutivas com anotações manuais, ao passo que esta pesquisa emprega BMC
automatizado com geração de contraexemplos concretos. A combinação de verificação
automatizada, conferência contra o código de produção, validação cruzada com o
CBMC e reporte direto à comunidade \textit{open-source} é a abordagem proposta
para firmware de autopiloto de VANTs, endereçando uma lacuna observada no
levantamento.

\chapter{Metodologia}
\label{cap:metodologia}

%% LC 01 (resolvido 19/10/2026): "Incluir tabela: ferramentas, versoes, flags e
%% hipoteses de modelagem para cada etapa." -> Tabela tab:etapas-ferramentas.

Este capítulo apresenta a metodologia empregada na verificação formal de módulos críticos do PX4 Autopilot com o ESBMC~\cite{Menezes2024}. O procedimento está organizado em oito etapas, conforme a numeração adotada no artigo derivado: (i) seleção dos módulos críticos; (ii) instrumentação para abstrair dependências de \textit{hardware}; (iii) especificação das propriedades de segurança; (iv) configuração e execução do ESBMC; (v) análise dos contraexemplos e conferência com o código de produção; (vi) reporte à comunidade PX4; (vii) validação cruzada com o CBMC e verificação de \textit{patches}; e (viii) reprodução em outro autopiloto (ArduPilot) e integração contínua. O módulo GPS uBlox é utilizado como exemplo para apresentar, de forma concreta, a aplicação dessas etapas.

\corr{A metodologia foi aplicada em duas campanhas, que totalizam 22 módulos distintos. A \textbf{primeira campanha}, apresentada no artigo derivado, abrangeu doze módulos relacionados a sensoriamento, estimação, navegação e comunicação. Nessa etapa, os \textit{harnesses} foram construídos a partir da transcrição de trechos do código do PX4. A \textbf{segunda campanha}, desenvolvida no âmbito desta dissertação, acrescentou dez módulos em três frentes: oito candidatos adicionais em módulos de entrada e de sistema, o driver de saída RoboClaw e o driver de bateria batmon. Nessa segunda etapa, a verificação e a conferência foram realizadas sobre o código-fonte \emph{original} do PX4, com \textit{replay} sob sanitizadores (Seção~\ref{sec:segunda-campanha}). As duas campanhas utilizaram o commit \texttt{8be2590c15}, exceto o batmon, cuja análise foi conferida contra o \texttt{main} de 13/10/2026 (commit \texttt{c9b9e4e25e}).}

As etapas~(i) e~(iv)--(vi) seguem o fluxo de Sousa et
al.~\cite{Sousa2026}; as etapas~(ii) e~(iii) são específicas de \textit{firmware},
e as etapas~(vii) e~(viii) são acréscimos deste trabalho. Dentro da etapa~(v), a
conferência contra o código de produção é a que separa violação alcançável no
modelo de falso positivo: foi ela que reclassificou o candidato do contador de
FIFO do BMI088 (Seção~\ref{subsec:imu-fifo}), que já havia sido proposto
\textit{upstream}, como falso positivo.

%% LC 02 (resolvido 19/10/2026): ordem real da primeira campanha.
\corr{Embora as etapas sejam apresentadas em ordem lógica, sua execução não seguiu essa sequência em todos os casos da primeira campanha. Nos casos do GPS e do contador de FIFO do BMI088, o reporte (vi) ocorreu imediatamente após a execução do ESBMC (iv), pois as \textit{issues} foram abertas para solicitar a triagem da comunidade. A conferência com o código de produção (v) foi realizada posteriormente, a partir da análise do código e das manifestações dos mantenedores. Assim, nesses casos, a sequência efetiva foi (i)--(iv), (vi), (v) e (vii). Essa inversão é considerada na discussão das ameaças à validade (Seção~\ref{sec:ameacas-validade}). Na segunda campanha, adotou-se a ordem prevista: nenhum candidato foi reportado antes da conferência com o código de produção e da execução do \textit{replay} (Seção~\ref{sec:segunda-campanha}).}

\corr{A Tabela~\ref{tab:etapas-ferramentas} resume, para cada etapa, as
ferramentas, versões, opções e hipóteses de modelagem empregadas.}

\begin{table}[!htbp]
\centering
\caption{\corr{Ferramentas, versões, opções e hipóteses de modelagem por etapa}}
\label{tab:etapas-ferramentas}
\scriptsize
\renewcommand{\arraystretch}{1.1}
\setlength{\tabcolsep}{3pt}
\newcommand{\cel}{\raggedright\hspace{0pt}}
\begin{tabular}{|p{2.0cm}|p{3.7cm}|p{4.9cm}|p{4.1cm}|}
\hline
\cel\textbf{Etapa} & \cel\textbf{Ferramentas e versões} & \cel\textbf{Opções} & \cel\textbf{Hipóteses de modelagem} \tabularnewline
\hline
\cel\corr{(i) Seleção} & \cel\corr{leitura do código e da documentação do PX4 (commit \texttt{8be2590c15})} & \cel\corr{---} & \cel\corr{módulo isolável por \textit{stubs}; amostra não estatística} \tabularnewline
\hline
\cel\corr{(ii) \textit{Stubs}} & \cel\corr{C/C++; \texttt{nondet\_T()}, \texttt{\_\_ESBMC\_assume} (1.ª campanha) e \texttt{\_\_VERIFIER\_assume} (2.ª)} & \cel\corr{---} & \cel\corr{E/S de \textit{hardware} sobreaproximada por valores não determinísticos; pré-condições da Tabela~\ref{tab:hipoteses-harness}} \tabularnewline
\hline
\cel\corr{(iii) Propriedades} & \cel\corr{\textit{assertions} no \textit{harness}} & \cel\corr{---} & \cel\corr{tradução de documento (\textit{datasheet}, manual, código) em predicado sobre variável de programa} \tabularnewline
\hline
\cel\corr{(iv) ESBMC} & \cel\corr{1.ª: ESBMC 8.1.0, Boolector 3.2.4, i7-13700K, Ubuntu 24.04. 2.ª: ESBMC 8.1.0--8.5.0, Boolector/Z3, máquina da autora, Ubuntu 24.04} & \cel\corr{1.ª: \texttt{--unwind 10 --no-unwinding-assertions --overflow-check}. 2.ª: \texttt{--unwind} por teste, com e sem \textit{unwinding assertions}, \texttt{--overflow-check --ub-shift-check}} & \cel\corr{limite $k$; sem \textit{unwinding assertions}, laços maiores que $k$ são truncados} \tabularnewline
\hline
\cel\corr{(v) Contraexemplos} & \cel\corr{1.ª: leitura do código. 2.ª: \textit{replay} com g++~13.3.0 e ASan/UBSan} & \cel\corr{\texttt{-std=c++17 -O0} \texttt{-fsanitize=address,undefined}} & \cel\corr{E/S simulada de forma determinística} \tabularnewline
\hline
\cel\corr{(vi) Reporte} & \cel\corr{GitHub (PX4-Autopilot)} & \cel\corr{---} & \cel\corr{---} \tabularnewline
\hline
\cel\corr{(vii) Validação cruzada} & \cel\corr{CBMC 5.95.1 (MiniSat)} & \cel\corr{opções equivalentes (Tabela~\ref{tab:equivalencia-opcoes})} & \cel\corr{mesmo \textit{harness} e mesmas pré-condições do ESBMC} \tabularnewline
\hline
\cel\corr{(viii) ArduPilot e CI} & \cel\corr{CBMC 5.95.1; GitHub Actions (\texttt{ubuntu-latest})} & \cel\corr{\texttt{--unwind 10 --signed-overflow-check --unsigned-overflow-check --bounds-check}} & \cel\corr{lógica equivalente no driver do ArduPilot} \tabularnewline
\hline
\end{tabular}
\end{table}

\section{Visão Geral da Abordagem}
\label{sec:visao-geral}

\corr{A Figura~\ref{fig:metodologia-visao-geral} sintetiza as oito etapas da metodologia e indica como elas foram aplicadas nas duas campanhas. O fluxo destaca a conferência com o código de produção, necessária para distinguir violações reproduzíveis de resultados decorrentes das limitações do modelo.} \corr{As
etapas~(i) e~(iv)--(vi) seguem o fluxo de Sousa et al.~\cite{Sousa2026}; as
etapas~(ii) e~(iii) são específicas de \textit{firmware}, e as etapas~(vii)
e~(viii) são acréscimos deste trabalho. A conferência (etapa~v) é feita por
leitura do código na primeira campanha e por \textit{replay} do código real sob
sanitizadores na segunda. Apenas defeitos reproduzidos e violações alcançáveis
no modelo seguem para o reporte~(vi); falsos positivos e resultados sem violação
encerram o fluxo do módulo.}

%% LC 03 (resolvido 19/10/2026): seta "sem violacao" -> reporte removida.
\begin{figure}[!tb]
\centering
\resizebox{0.92\linewidth}{!}{%
\begin{tikzpicture}[
  font=\linespread{1}\footnotesize,
  >={Stealth[length=2mm]},
  node distance=4mm,
  base/.style={draw=palEdge, rounded corners=2pt, align=center, inner sep=3pt,
    fill=palPaper},
  etapa/.style={base, text width=5.8cm, minimum height=7mm},
  firm/.style={etapa, fill=palFill},
  nova/.style={base, text width=4.8cm, minimum height=7mm, draw=palEdgeDark,
    fill=palStrong},
  dir/.style={base, text width=4.8cm, minimum height=7mm},
  conf/.style={base, text width=6.2cm, fill=palFill, draw=palEdgeDark, very thick},
  dec/.style={draw=palEdge, diamond, aspect=2.4, inner sep=1pt, fill=palPaper,
    align=center},
  desf/.style={base, text width=3.0cm, minimum height=10mm, fill=palStrong,
    draw=palEdgeDark},
  outn/.style={base, text width=3.0cm, minimum height=10mm, fill=palPaper},
  rot/.style={font=\scriptsize\itshape, align=center}]

  \node[etapa] (e1) {(i) Seleção de módulos críticos};
  \node[firm, below=of e1] (e2) {(ii) Instrumentação com \textit{stubs} de \textit{hardware}};
  \node[firm, below=of e2] (e3) {(iii) Especificação de propriedades\\(\textit{datasheet}, documentação, código)};
  \node[etapa, below=of e3] (e4) {(iv) Configuração e execução do ESBMC};
  \node[etapa, below=of e4] (e5) {(v) Análise de contraexemplos};

  \node[conf, below=6mm of e5] (conf) {\textbf{Conferência contra o código de produção}\\[1pt]
    1.\textsuperscript{a} campanha: leitura do código\\
    2.\textsuperscript{a} campanha: \textit{replay} do código real {\footnotesize+} ASan/UBSan};
  \node[dec, below=6mm of conf] (dec) {Reproduz em\\produção?};

  \node[desf, below=10mm of dec] (def) {Defeito reproduzido / VAM};
  \node[outn, left=20mm of def] (fp) {Falso positivo (modelagem / ferramenta)};
  \node[outn, right=20mm of def] (ok) {Sem violação até o limite $k$};

  \node[dir, right=2.4cm of e1.north east, anchor=north west] (e6) {(vi) Reporte à comunidade PX4\\(condicionado à revisão)};
  \node[nova, below=of e6] (e7) {(vii) Validação cruzada (CBMC)\\e verificação de \textit{patch}};
  \node[nova, below=of e7] (e8) {(viii) Reprodução (ArduPilot)\\e integração contínua};

  \foreach \a/\b in {e1/e2,e2/e3,e3/e4,e4/e5,e5/conf,conf/dec,e6/e7,e7/e8}
    \draw[->] (\a) -- (\b);

  \draw[->] (dec) -- node[rot, right] {sim} (def);
  \draw[->] (dec.west) -| node[rot, pos=0.25, above] {guarda\\omitida} (fp.north);
  \draw[->] (dec.east) -| node[rot, pos=0.25, above] {nenhuma\\violação} (ok.north);

  \draw[->] (def.south) -- ++(0,-5mm) -| ([xshift=7mm]e6.east) -- (e6.east);

  \node[firm, text width=3mm, minimum height=3mm, below=12mm of fp.south west,
    anchor=north west] (l1) {};
  \node[right=1mm of l1, font=\scriptsize, anchor=west] (t1) {etapas de \textit{firmware}};
  \node[nova, text width=3mm, minimum height=3mm, right=5mm of t1] (l2) {};
  \node[right=1mm of l2, font=\scriptsize, anchor=west] (t2) {acréscimos deste trabalho};
  \node[conf, text width=3mm, minimum height=3mm, right=5mm of t2] (l3) {};
  \node[right=1mm of l3, font=\scriptsize, anchor=west] {conferência (1.\textsuperscript{a} e 2.\textsuperscript{a} campanhas)};
\end{tikzpicture}}
\caption[Metodologia de verificação em oito etapas]{Metodologia de verificação em oito etapas, aplicada em duas campanhas. Elaborado pela autora.}
\label{fig:metodologia-visao-geral}
\end{figure}

%% Prof. Lucas: "O que voce define como verificacao modular?" -- CORRIGIDO
\revisao{Foi adotada a verificação modular --- abordagem na qual o sistema é
decomposto em unidades funcionais (módulos) verificadas de forma independente, e
não como um único artefato monolítico. Cada módulo é isolado de suas dependências
por meio de abstrações (\textit{stubs}), e as propriedades de segurança são
verificadas no escopo local da unidade, o que reduz o tamanho do espaço de
estados analisado em cada execução do verificador~\cite{Cordeiro2012}.}

\revisao{A escolha decorre da dimensão da base de código.}
A árvore \texttt{src/} do PX4 soma aproximadamente 710 mil linhas de
código C/C++ (arquivos-fonte e cabeçalhos, contadas com \texttt{wc -l} no
commit 8be2590c15)\revisao{.
Essa dimensão torna a verificação monolítica impraticável devido à
explosão combinatória do espaço de estados~\cite{Biere1999}.}

%% Prof. Lucas: "O que voce define como remessa?" -- CORRIGIDO
\corr{A verificação da primeira campanha foi organizada em três rodadas. Neste
trabalho, uma rodada designa um conjunto de módulos verificados em um mesmo ciclo
de trabalho, agrupados por afinidade funcional ou por prioridade de criticidade.
A organização em rodadas permite escalonar o esforço de verificação, priorizando
primeiro os componentes cuja falha tem maior impacto sobre a segurança de voo e
ampliando a cobertura em ciclos subsequentes.}

%% LC 04 (resolvido 19/10/2026): definicao de harness em destaque.
\corr{\paragraph{Definição de \textit{harness}.} Neste trabalho, um
\textbf{\textit{harness} de verificação} é um arquivo C/C++ escrito
especificamente para a verificação, que (a)~declara as entradas do código sob
análise como valores não determinísticos, (b)~restringe essas entradas por
pré-condições (\textit{assumptions}), (c)~invoca as funções do módulo e
(d)~expressa as propriedades de segurança como \textit{assertions}. O
\textit{harness} tem uma função de entrada própria (\texttt{main()} na primeira
campanha; \texttt{harness\_main()} na segunda, porque o código original do PX4
já define uma \texttt{main} por meio de \texttt{ModuleBase}) e seleciona, por
uma variável de escolha, qual dos testes do arquivo será verificado.}

O ESBMC suporta verificação modular por contratos de função, nos
quais pré-condições e pós-condições são especificadas formalmente para cada
função\footnote{Documentação oficial do ESBMC sobre contratos de função: \url{https://esbmc.github.io/docs/function-contracts/}}. Neste trabalho, optou-se por uma abordagem
complementar: em vez de contratos explícitos, as propriedades foram expressas
diretamente como \textit{assertions} nos \textit{harnesses}.
%% LC 05 (resolvido 19/10/2026): justificativa reescrita por campanha.
\corr{Duas razões justificam essa escolha. Primeiro, as propriedades de
interesse vêm de documentos externos ao código (\textit{datasheets}, manuais e
especificações de protocolo) e se expressam naturalmente como condições sobre
o resultado de uma chamada, que é o que uma \textit{assertion} no
\textit{harness} captura. Segundo, um contrato precisaria ser anotado na própria
função verificada: na primeira campanha, as funções transcritas no
\textit{harness} são curtas e não justificariam a sobrecarga de contratos; na
segunda, anotar contratos exigiria modificar o código-fonte do PX4, o que essa
campanha justamente evita.}

\subsection{Etapa~(i) --- Seleção de Módulos Críticos}

A seleção dos módulos considerou dois critérios complementares: a relevância funcional para a segurança de voo e a viabilidade técnica da verificação isolada. O primeiro foi avaliado pela posição do módulo na cadeia de dados e de controle. Assim, receberam prioridade os drivers de sensores que fornecem dados ao estimador de estado (EKF2), os componentes de comunicação que processam comandos externos e os módulos de navegação responsáveis pela definição de trajetórias autônomas. Falhas nesses elementos podem afetar diretamente as decisões de voo. O segundo critério considerou a possibilidade de isolar o módulo das demais partes do \textit{firmware} por meio de abstrações, sem descaracterizar a lógica sob análise. \corr{Na segunda campanha, esses mesmos critérios orientaram a seleção de módulos de entrada, de sistema e de saída (Seção~\ref{sec:segunda-campanha}).}

%% LC 06 (resolvido 19/10/2026): "escrito manualmente" -> "assistencia de IA
%% generativa e revisadas pela autora".
\revisao{Cabe justificar a divisão em rodadas, já que o tempo total de execução do
verificador é da ordem de poucos minutos e, isoladamente, não a justificaria. O custo
dominante desta metodologia não é o tempo de máquina, e sim o \textbf{esforço humano} que o
antecede e o sucede: leitura do módulo, identificação das dependências a abstrair, construção
do \textit{harness}, derivação das propriedades a partir da documentação técnica e,
sobretudo, triagem dos contraexemplos.} \corr{Os \textit{harnesses} da primeira campanha somam
aproximadamente $2.800$ linhas, construídas com assistência de IA generativa e
revisadas pela autora, conforme a Declaração de Uso de Inteligência Artificial,
nas páginas iniciais desta dissertação, ao passo que a campanha completa executa
em cerca de dois minutos e meio.}

\revisao{É esse desequilíbrio que organiza o ciclo. Verificar em rodadas permite que o
aprendizado obtido em um conjunto de módulos --- quais abstrações são adequadas, quais
configurações do verificador cada perfil de módulo exige, quais padrões de contraexemplo são
recorrentes --- seja incorporado antes de investir esforço no conjunto seguinte. A necessidade
de \texttt{--unwind 300} no MAVLink e de \texttt{--no-library} no VTOL, por exemplo, só se
tornou conhecida durante a primeira rodada e informou a configuração das seguintes. A divisão
é, portanto, um mecanismo de organização do trabalho e de priorização por criticidade, e não
uma consequência do custo computacional.}

A organização em rodadas reflete essa priorização: a primeira rodada concentra
os drivers de sensores e comunicação de maior criticidade, cuja falha pode
comprometer toda a missão; as rodadas subsequentes ampliam a cobertura para
monitoramento, estimação de estado e navegação. Esse escalonamento permite
identificar \revisao{candidatos} de alto impacto antes de investir tempo de
verificação em módulos auxiliares.

\subsection{Etapa~(ii) --- Instrumentação com \textit{Stubs}}

%% LC 23 (resolvido 19/10/2026): etapa reescrita distinguindo as duas campanhas.
%% LC 22 (resolvido 19/10/2026): explicacao de nao determinismo mantida so aqui.
\corr{A instrumentação tem por finalidade abstrair as entradas e saídas de \textit{hardware}, as chamadas ao sistema operacional NuttX e as dependências de outros módulos. Para isso, utiliza-se a primitiva \texttt{nondet\_T()}, em que \texttt{T} representa o tipo do dado. Essa primitiva permite ao verificador considerar qualquer valor representável por esse tipo, incluindo condições extremas, dados inválidos e respostas inesperadas do \textit{hardware}, sem exigir a enumeração manual de casos de teste~\cite{Menezes2024}. Por exemplo, \texttt{nondet\_uint8()} representa qualquer valor entre $0$ e $255$.}

\corr{As duas campanhas aplicam essa abstração de formas diferentes. Na
\textbf{primeira}, os \textit{harnesses} não contêm funções de \textit{stub}: os
bytes que o driver leria do barramento (SPI, I2C ou UART) são declarados como
valores não determinísticos na própria função de teste e passados diretamente ao
trecho do driver transcrito no \textit{harness}
(Listagens~\ref{lst:stub_i2c} e~\ref{lst:main_imu}). Na \textbf{segunda}, como o
código verificado é o original, as funções de que ele depende precisam existir:
foram escritos \textit{stubs} de fato --- \texttt{read}, \texttt{write} e
\texttt{select} da UART, o barramento I2C, o \textit{middleware} uORB, o tempo, os
parâmetros e o registro de mensagens ---, que preservam a assinatura original e
devolvem valores não determinísticos (Seção~\ref{sec:segunda-campanha}).}

%% LC 07 (resolvido 19/10/2026): pre-condicoes tabuladas e justificadas.
\corr{Os valores não determinísticos, porém, não ficam irrestritos: cada
\textit{harness} restringe parte das entradas por pré-condições
(\texttt{\_\_ESBMC\_assume}), seja para refletir a faixa que o código de produção
de fato recebe, seja para concentrar a busca numa região de interesse. Essas
restrições delimitam o alcance de cada veredito e estão tabuladas, com a
justificativa de cada uma, na Tabela~\ref{tab:hipoteses-harness}. Pré-condições
estreitas demais podem mascarar defeitos; pré-condições largas demais admitem
estados que o driver real rejeita e produzem falsos positivos --- o mecanismo do
candidato do contador de FIFO. Por isso, toda violação passa pela conferência
contra o código de produção da etapa~(v).}

\begin{table}[!htb]
\centering
\caption[Pré-condições dos \textit{harnesses} da primeira campanha]{\corr{Principais pré-condições dos \textit{harnesses} da primeira campanha e sua justificativa}}
\label{tab:hipoteses-harness}
\footnotesize
\begin{tabular}{|p{2.3cm}|p{5.6cm}|p{6.0cm}|}
\hline
\textbf{Módulo} & \textbf{Pré-condição (\texttt{\_\_ESBMC\_assume})} & \textbf{Justificativa} \\
\hline
\corr{GPS uBlox} & \corr{\texttt{0 < input\_len <= GPS\_DUMP\_DATA\_SIZE + 10}; \texttt{dump\_buffer.len >= GPS\_DUMP\_DATA\_SIZE - 5}} & \corr{tamanho do vetor de entrada do \textit{harness}; estado inicial próximo da fronteira (discutido na Seção~\ref{sec:exemplo-ilustrativo})} \\
\hline
\corr{IMU BMI088} & \corr{eixos do acelerômetro e do giroscópio em \texttt{[INT16\_MIN, INT16\_MAX]}; bytes de temperatura e FIFO irrestritos} & \corr{faixa do tipo; não restringe nada além do tipo} \\
\hline
\corr{MAVLink} & \corr{\texttt{0 < packet\_len <= BUFFER\_SIZE + 50}} & \corr{admite, de propósito, pacotes maiores que o \textit{buffer} (280 bytes), para exercitar a guarda de capacidade} \\
\hline
\corr{Expo RC} & \corr{entrada em $[-1;\,1]$ e expoente em $[0;\,1]$, finitos} & \corr{faixa normalizada do controle remoto e do parâmetro de expo} \\
\hline
\corr{RC Update} & \corr{\texttt{min} em $[800;1200]$, \texttt{trim} em $[1400;1600]$, \texttt{max} em $[1800;2200]$, \texttt{min < trim < max}} & \corr{faixas de calibração de canal em microssegundos} \\
\hline
\corr{PWM Output} & \corr{saída em $[\texttt{PWM\_MOTOR\_MIN}; \texttt{PWM\_MOTOR\_MAX}]$ no teste nominal; $[-10^4;\,10^4]$ no teste de robustez} & \corr{faixa de motor; teste de robustez com entrada fora da faixa} \\
\hline
\corr{EKF2} & \corr{variâncias finitas, não negativas e limitadas por \texttt{VAR\_VEL\_MAX}; \texttt{min\_var <= max\_var}} & \corr{limites das variâncias usados no estimador} \\
\hline
\corr{Magnetômetro} & \corr{variâncias finitas, em $[0;\,\texttt{MAX\_VAR\_ALLOWED}]$, com soma positiva; ângulos em $[0;\,\pi]$} & \corr{evita divisão por zero no ganho; faixa de ângulo} \\
\hline
\corr{Navigator} & \corr{índices \texttt{<= MAX\_MISSION\_ITEMS}; distância em $[0;\,10^5]$ m; raio de aceitação em $[0{,}1;\,1000]$ m} & \corr{tamanho máximo de missão; faixas físicas plausíveis} \\
\hline
\corr{VTOL} & \corr{densidade do ar em $[0{,}1;\,3]$; pesos em $[0{,}1;\,100]$; velocidade de transição em $[5;\,50]$ m/s} & \corr{faixas físicas plausíveis} \\
\hline
\corr{Barômetro} & \corr{pressão bruta em $[50000;\,110000]$ Pa; \textit{offsets} limitados; prioridade em $[-200;\,200]$} & \corr{faixa de operação de barômetros; teste de prioridade além da faixa válida} \\
\hline
\corr{Battery Monitor} & \corr{SoC em $[0;\,1]$; limiares ordenados; tensão de célula em $[3{,}0;\,4{,}5]$ V; 1 a 12 células} & \corr{faixas de bateria LiPo e de configuração} \\
\hline
\end{tabular}
\end{table}

%% LC 08 (resolvido 19/10/2026): citacao Godefroid2005 removida.
\corr{Para cada módulo, um arquivo de \textit{harness} centraliza as entradas
não determinísticas, as pré-condições, as funções de teste e a função de
entrada.}

\subsection{Etapa~(iii) --- Especificação de Propriedades}

As propriedades verificadas foram organizadas em duas categorias: automáticas e customizadas. As propriedades automáticas são habilitadas por opções de linha de comando do ESBMC e abrangem classes gerais de defeitos, como acessos fora dos limites de vetores e \textit{buffers}, desreferências de ponteiros nulos e \textit{overflows} aritméticos em inteiros~\cite{Cordeiro2012}. Essas verificações dispensam \textit{assertions} específicas para cada módulo e podem ser aplicadas de maneira uniforme, desde que as opções correspondentes estejam habilitadas.

%% LC 09 (resolvido 19/10/2026): rastreabilidade atenuada (1.a) e afirmada (2.a).
As propriedades customizadas capturam requisitos específicos de cada módulo,
derivados de datasheets de fabricantes e da especificação do protocolo PX4. São
expressas como assertions inseridas no harness, imediatamente após a chamada às
funções verificadas. \corr{A rastreabilidade entre cada \textit{assertion} e o
requisito de origem é registrada na Tabela~\ref{tab:proveniencia}. Nos
\textit{harnesses} da primeira campanha, os comentários junto às funções de
teste descrevem a especificação em linguagem natural (por exemplo,
``ESPECIFICAÇÃO: Converter dados raw de temperatura para Celsius''), mas não
indicam a seção nem a página do documento. Na segunda campanha, cada
\textit{harness} anota, no cabeçalho, a origem de cada pré-condição e de cada
propriedade, com seção e página do documento quando ele existe (por exemplo,
o manual da RoboClaw, Rev.~6.0, pp.~48--49).}

\subsection{Etapa~(iv) --- Configuração e Execução do ESBMC}

A configuração do ESBMC depende, entre outros fatores, do limite de desenrolamento e do \textit{solver} SMT selecionado. O parâmetro \texttt{--unwind} determina quantas iterações de cada laço serão exploradas. Limites menores reduzem o custo da análise, mas podem deixar de alcançar defeitos que ocorram em iterações posteriores. Já a opção \texttt{--default-solver} seleciona o \textit{solver} utilizado e pode afetar tanto o tempo de execução quanto a obtenção de um resultado dentro do tempo-limite. Nesta pesquisa, o Z3 não produziu um veredito dentro do \textit{timeout} nos testes indicados para GPS, MAVLink e EKF2 com \texttt{--unwind 10}; por esse motivo, o Boolector foi adotado nas execuções correspondentes da primeira campanha (Seção~\ref{sec:ameacas-validade}).
%% LC 10 (resolvido 19/10/2026): bit-blasting de ponto flutuante do ESBMC.
\corr{Como o Boolector não suporta a teoria de ponto flutuante do padrão
SMT-LIB, o ESBMC codifica as operações IEEE-754 de precisão simples e dupla como
operações sobre vetores de bits (\textit{bit-blasting} de ponto flutuante),
preservando a semântica de arredondamento, os valores especiais (NaN, infinitos)
e os números subnormais; com solvers que suportam essa teoria, como o Z3, a
codificação pode usá-la diretamente~\cite{Gadelha2020}. Essa codificação
precisa é essencial para os módulos de sensores e de estimação, que usam
aritmética de ponto flutuante.}

A configuração é iterativa: parte-se de parâmetros padrão e ajusta-se conforme
os resultados. Uma violação de asserção de desenrolamento (\textit{unwinding
assertion}) indica que o limite atual é insuficiente para cobrir os laços do
módulo, exigindo aumento do parâmetro \texttt{--unwind}, como ocorreu
no MAVLink ($k = 10 \rightarrow 300$). Quando as funções da biblioteca matemática
tornam o modelo intratável, a opção \texttt{--no-library} desativa a biblioteca C
abstrata do ESBMC: funções como \texttt{sinf} e \texttt{sqrtf} passam a
retornar valores não determinísticos, uma sobreaproximação que não deixa escapar
violações das demais verificações, mas deixa de verificar as propriedades
internas à biblioteca --- caso do VTOL~\cite{Menezes2024}.

\subsection{Etapa~(v) --- Análise de Contraexemplos e Conferência contra o Código de Produção}

Quando o ESBMC reporta uma violação, o contraexemplo gerado
apresenta valores concretos de entrada que conduzem ao estado de falha.
\revisao{A análise começa pela leitura do traço: identifica-se qual propriedade foi
violada, em que linha, e qual atribuição de valores conduziu ao estado. Nessa
etapa o contraexemplo ainda é um \textbf{candidato}, na acepção fixada no
Capítulo~\ref{cap:introducao} --- não há, até aqui, evidência de que o estado
seja alcançável fora do ambiente de verificação.}

\revisao{A classificação e a avaliação de impacto ocorrem depois, com base nas
consequências da falha no contexto de voo. Um \textit{overflow} aritmético em
variável de controle de atitude tem impacto distinto de uma violação de asserção
em caminho inalcançável em condições normais. Essa triagem guia quais resultados
seguem para as etapas seguintes e com qual prioridade.}

\paragraph{Conferência contra o código de produção.}
\label{subsec:etapa-replay}

\revisao{Esta conferência decide se um candidato é violação alcançável no modelo ou falso
positivo, e é a que sustenta a distinção que organiza os resultados deste
trabalho. Um contraexemplo é sempre válido \emph{para o modelo codificado}; o que
se verifica é se o estado violado persiste quando o
\textit{harness} sai de cena e o código de produção entra em seu lugar.}
Ela pode ser feita de duas formas: por \textit{replay} ou por leitura do
código de produção.

%% LC 11 (resolvido 19/10/2026): protocolo reescrito pelo que foi executado na
%% segunda campanha (candidatos extras, RoboClaw e batmon).
\corr{No \textit{replay}, o arquivo-fonte \emph{original} do driver é compilado
com um compilador comum --- g++~13.3.0, com \texttt{-std=c++17 -O0} --- junto com
\textit{stubs} determinísticos que substituem o \textit{hardware} e injetam os
valores exatos do contraexemplo: a UART simulada registra cada pacote enviado e
devolve as respostas escolhidas pelo programa de \textit{replay}. O programa é
executado, sob os sanitizadores de endereço (ASan) e de comportamento indefinido
(UBSan), e registra o efeito observável do caminho do contraexemplo, como o byte
efetivamente enviado ao atuador ou o acesso a um índice fora dos limites.}

\corr{O critério de decisão é binário. Se a execução do código original
reproduz o efeito previsto pelo contraexemplo --- por exemplo, um byte fora da
faixa especificada no pacote enviado ---, o candidato é classificado como
\textbf{violação alcançável no modelo} e, conforme a análise de alcançabilidade
no uso real, como defeito. Se não reproduz, investiga-se a razão: se o código de
produção contém guarda que o \textit{harness} omitiu, o candidato é
\textbf{falso positivo}. Duas verificações de sanidade acompanham o
\textit{replay}: confirmar que a execução alcança a função analisada, e não
termina antes, no código de preparação; e confirmar que o cenário crítico é
derivado da entrada injetada, e não fixado por atribuição direta no programa de
\textit{replay} --- caso em que se estaria demonstrando aritmética, e não
alcançabilidade.}

Nos três contraexemplos da primeira campanha, a conferência foi
feita por leitura do código de produção: o candidato do contador de FIFO foi
refutado pela inspeção do driver do BMI088, e a inalcançabilidade do estado
violador do GPS veio da revisão do mantenedor e da análise do código-fonte,
depois verificada por prova indutiva (Seção~\ref{subsec:gps-inducao}).
\corr{Na segunda campanha, o \textit{replay} com sanitizadores foi executado para
os candidatos alcançáveis dos oito módulos adicionais, do driver da RoboClaw e do
batmon (Seção~\ref{sec:segunda-campanha}); os programas de \textit{replay} e seus
registros estão no repositório do trabalho (\texttt{candidatos\_extras\_v5/},
\texttt{testes\_hardware/roboclaw/} e \texttt{novo\_achado\_batmon/}). Registra-se,
como limitação, que o \textit{replay} executável não foi aplicado aos candidatos
do contador de FIFO e do GPS da primeira campanha, cuja conferência se deu apenas
por leitura do código de produção; essa diferença de procedimento é retomada nas
ameaças à validade (Seção~\ref{sec:ameacas-validade}).}

\subsection{Etapa~(vi) --- Reporte \corr{à} Comunidade PX4}

%% LC 12 (resolvido 19/10/2026): desfecho de cada reporte e momento do reporte.
\corr{Os candidatos que sobrevivem à conferência são reportados como
\textit{issues} no repositório oficial do PX4 no GitHub, acompanhados de:
descrição do achado, contraexemplo com valores concretos, trecho de código
afetado, referência ao documento que define o comportamento esperado e, quando
houver, proposta de correção. Esse formato permite aos mantenedores reproduzir o
problema sem acesso ao ambiente de verificação~\cite{PX4_GitHub}.}

\corr{O reporte não garante aceitação --- e, nesta pesquisa, os da primeira
campanha não foram aceitos como defeitos. A observação de temperatura do BMI088
(\textit{Issue}~\#26865) foi recusada pelos mantenedores, que não reconheceram a
faixa operacional do \textit{datasheet} como restrição de software
(Seção~\ref{subsec:imu-temperatura}); o PR~\#27575, que tratava da temperatura e
do FIFO, também não foi aceito; e a guarda proposta para o GPS
(\textit{Issue}~\#26866, PR~\#27576) foi enquadrada pelo mantenedor como
endurecimento defensivo, depois que a revisão mostrou o estado violador
inalcançável no driver de produção. A revisão pela comunidade é, assim, a única
etapa capaz de elevar um achado à condição de defeito confirmado, e também a que
pode recusá-lo.}

\corr{O momento do reporte pesou nesses desfechos. Para o GPS e para o contador
de FIFO, a \textit{issue} foi aberta antes da conferência contra o código de
produção, pedindo a própria triagem; a conferência veio depois, e no caso do FIFO
rebaixou o candidato a falso positivo. A lição incorporada à segunda campanha foi
reportar apenas depois da conferência, do \textit{replay}, da consulta à
especificação com página citada e da verificação de que o código ainda é o mesmo
na branch principal e de que não há \textit{issue} aberta sobre o assunto. Em
coerência com essa lição e com a orientação da pesquisa, nenhum achado da segunda
campanha foi reportado sem a prévia revisão do orientador. No caso do RoboClaw,
em que o driver havia sido reescrito na branch principal, o relato indicou
explicitamente o ramo de \textit{release} em que o código analisado permanece
(Seção~\ref{subsec:2c-roboclaw}).}

\subsection{Etapa~(vii) --- \revisao{Validação Cruzada com Verificador Independente} e Verificação de \textit{Patch}}
\label{subsec:etapa-validacao-cruzada}

\revisao{Os candidatos são submetidos a um segundo
verificador, o CBMC~5.95.1~\cite{Kroening2014}, sobre o mesmo módulo e com opções
equivalentes. O objetivo é distinguir defeito de artefato de \textit{toolchain}:
se apenas um dos verificadores reporta a violação, a hipótese de artefato ganha
força; se ambos reportam, ela perde.}

\revisao{Cabe uma ressalva de escopo, aplicada de forma sistemática no
Capítulo~\ref{cap:experimentos}: a concordância entre dois verificadores sobre o
mesmo \textit{harness} corrobora a codificação do problema, e não a
alcançabilidade do estado no driver de produção. Os dois partem da mesma
abstração e das mesmas pré-condições. Acresce que os verificadores podem
classificar o mesmo defeito sob propriedades formais distintas, conforme tratem a
conversão entre tipos com e sem sinal --- o que é informação adicional, e não
contradição.} \corr{Na segunda campanha, a comparação foi estendida a todos os
\textit{harnesses}, e não apenas aos candidatos, para medir também a
aplicabilidade de cada ferramenta ao código do PX4
(Seção~\ref{subsec:2c-protocolo}).}

Quando uma correção é proposta, o mesmo \textit{harness} é reverificado com
o \textit{patch} aplicado, nas mesmas pré-condições: sem o \textit{patch}, a
violação deve aparecer; com ele, deve desaparecer. A comparação só isola o efeito
da guarda se as pré-condições forem as mesmas nos dois casos
(Seção~\ref{subsec:patch-gps}).

%% LC 13 (resolvido 19/10/2026): etapa renomeada; reproducao so com CBMC.
\subsection{Etapa~(viii) --- \corr{Reprodução em Outro Autopiloto e Integração Contínua}}
\label{subsec:etapa-ardupilot-ci}

Quando a lógica de um achado tem equivalente em outro autopiloto, ele é
conferido também lá. \corr{Na primeira campanha, essa reprodução foi feita
apenas com o CBMC~5.95.1, e não com o ESBMC; trata-se, portanto, de uma
reprodução da observação em outro código, e não de uma verificação completa do
outro autopiloto.} A observação de temperatura do BMI088 foi reproduzida
no driver equivalente do ArduPilot~\cite{ArduPilot}
(\texttt{AP\_InertialSensor\_BMI088.cpp}, linhas 383--385; \textit{Issue}
\#33272), onde um mantenedor encerrou a \textit{issue} citando a faixa de
medição de $-104$ a $150\,^{\circ}$C do \textit{datasheet}. A reprodução serve
apenas como verificação de consistência: os dois projetos recusaram a
observação.
\corr{A repetição dessa reprodução com o ESBMC~8.5, para que a etapa empregasse
as duas ferramentas também no ArduPilot, fica como refinamento futuro; como
executada, a reprodução usou somente o CBMC.}

%% LC 14 (resolvido 12/11/2026): workflow lido no fork (esbmc-verify.yml).
%% Observacao de fidelidade: o workflow instala e roda o CBMC (nao o ESBMC) sobre
%% versoes reduzidas dos harnesses do IMU e do GPS -- descrito como smoke check.
%% So o tempo medido de uma rodada permanece pendente (aba Actions).
\corr{Por fim, um \textit{workflow} do GitHub Actions
(\texttt{esbmc-verify.yml}), em um \textit{fork} do PX4
(\url{https://github.com/marcela2881/PX4-Autopilot}), executa a cada
\textit{push} e \textit{pull request} para a \textit{branch} \texttt{main} uma
verificação de fumaça (\textit{smoke check}) com o CBMC, em \textit{runner}
\texttt{ubuntu-latest}, sobre versões reduzidas dos \textit{harnesses} do IMU
BMI088 e do GPS uBlox (\texttt{--unwind 10}, com checagem de \textit{overflow}
com e sem sinal e de limites). O \textit{workflow} reexecuta essa verificação a
cada alteração do código, servindo de salvaguarda contra regressões; emprega o
CBMC, e não o ESBMC, por ser mais simples de instalar no ambiente de integração
contínua.}

\section{Exemplo Ilustrativo: GPS uBlox}
\label{sec:exemplo-ilustrativo}

A Listagem~\ref{lst:exemplo-harness-gps} apresenta o fragmento central do harness construído para o GPS uBlox, destacando os elementos essenciais da metodologia: a declaração da entrada não determinística (\revisao{\texttt{nondet\_size\_t()}}), a \textit{assumption} que restringe o domínio (\texttt{\_\_ESBMC\_assume}) e a \textit{assertion} que expressa a propriedade de segurança verificada. O código completo de todos os harnesses e os logs de execução do ESBMC estão disponíveis em \url{https://github.com/marcela2881/dissertacao-px4-esbmc}.

\begin{lstlisting}[
  caption={\revisao{Fragmento do harness de verificação do GPS uBlox, função \texttt{test\_gps\_real\_bit\_operation}, transcrito de \texttt{test\_gps\_driver.cpp}.}},
  label={lst:exemplo-harness-gps},
  numbers=left,
  frame=single,
  language=C]
// Entrada simbolica - o solver escolhe o valor
size_t input_len = nondet_size_t();

// Restricao do dominio: o buffer de entrada tem
// GPS_DUMP_DATA_SIZE + 10 = 210 bytes
__ESBMC_assume(input_len > 0 && input_len <= GPS_DUMP_DATA_SIZE + 10);

uint8_t input_data[GPS_DUMP_DATA_SIZE + 10];
gps_dump_s dump_buffer;

// Estado inicial do buffer, proximo da fronteira
dump_buffer.len = nondet_uint8();
__ESBMC_assume(dump_buffer.len >= GPS_DUMP_DATA_SIZE - 5);

// Execucao da funcao real do PX4 sobre a entrada simbolica
dumpGpsData(input_data, input_len, gps_dump_comm_mode_t::Full,
            /* msg_to_gps_device = */ true,
            &dump_buffer, gps_dump_comm_mode_t::Full);

// Propriedade de seguranca: o comprimento acumulado nao deve
// ultrapassar GPS_DUMP_DATA_SIZE, considerado o bit 7 de sinalizacao
assert(dump_buffer.len == 0 ||
       dump_buffer.len <= GPS_DUMP_DATA_SIZE ||
       (dump_buffer.len & 0x7F) <= GPS_DUMP_DATA_SIZE);
\end{lstlisting}

\revisao{As duas pré-condições delimitam o espaço de entradas e, com ele, o alcance
do veredito. A primeira restringe \texttt{input\_len} ao tamanho do vetor de
entrada declarado no próprio \textit{harness}, $210$ bytes, evitando que o
verificador reporte acessos fora dos limites do vetor de teste em vez do defeito
sob investigação. A segunda posiciona \texttt{dump\_buffer.len} a no máximo cinco
unidades da fronteira de $200$ bytes, concentrando a busca na vizinhança em que a
subtração sem sinal pode dar a volta. Trata-se de decisão de modelagem que
restringe a busca: violações que exigissem estados iniciais distantes da
fronteira não seriam alcançadas por este \textit{harness}.}

%% LC 15 (resolvido 19/10/2026): tautologia declarada; propriedade de memoria do
%% memcpy e a verificada.
\corr{A \textit{assertion} final da Listagem~\ref{lst:exemplo-harness-gps} é,
na prática, uma tautologia: o terceiro termo, \texttt{(dump\_buffer.len \& 0x7F)
<= GPS\_DUMP\_DATA\_SIZE}, é sempre verdadeiro, pois o valor mascarado com
\texttt{0x7F} nunca passa de $127$, que é menor que $200$. A \textit{assertion}
não detecta, portanto, nenhuma violação. A violação reportada pelo ESBMC vem da
verificação automática de limites dentro do \texttt{memcpy}
(\textit{dereference failure}: acesso fora dos limites do objeto), que é a
propriedade efetivamente verificada neste teste. Uma propriedade de saída
significativa seria \texttt{dump\_buffer.len <= GPS\_DUMP\_DATA\_SIZE} após a
chamada, sem os termos disjuntivos.}

\corr{Para separar o que é defeito do que é artefato dessa modelagem, o
\textit{harness} foi reexecutado com a propriedade de saída na forma não
tautológica (\texttt{dump\_buffer.len <= GPS\_DUMP\_DATA\_SIZE}) e partindo de um
estado inicial válido (\texttt{dump\_buffer.len < GPS\_DUMP\_DATA\_SIZE}). Nessas
condições o ESBMC~8.2.0 (solver Z3) verifica \textsc{verification successful} em
cerca de $389$~s: a lógica transcrita do \texttt{dumpGpsData} respeita o limite a
partir de qualquer estado válido. A violação observada na versão exploratória
decorre, portanto, exclusivamente do estado inicial acima da capacidade admitido
pela pré-condição --- estado inalcançável no driver de produção, como confirma a
prova indutiva da Seção~\ref{subsec:gps-inducao}. Os dois resultados são
consistentes: o modelo só viola a propriedade quando parte de um estado que o
código de produção nunca alcança, o que sustenta a classificação do achado como
violação alcançável no modelo, e não como defeito.}

%% LC 16 (resolvido 19/10/2026): violacao e admitida pela pre-condicao.
\corr{A segunda pré-condição merece uma ressalva adicional. Como ela admite
\texttt{dump\_buffer.len} entre $195$ e $255$, o \textit{harness} já parte de
estados em que o preenchimento está \emph{acima} da capacidade de $200$ bytes ---
o contraexemplo usa $207$. A violação, portanto, não é derivada a partir de um
estado inicial válido: ela é admitida pela pré-condição. O que o contraexemplo
mostra é que, \emph{se} o preenchimento passar da capacidade, a guarda do trecho
transcrito falha; não mostra que esse estado é alcançável. Foi exatamente essa
pergunta que a revisão do mantenedor e a prova indutiva da
Seção~\ref{subsec:gps-inducao} responderam negativamente para o driver de
produção.}

%% LC 17 (resolvido 19/10/2026): retirado "fiel"; duas diferencas declaradas.
\corr{O \textit{harness} também difere do driver em dois pontos, que impedem
chamá-lo de transcrição fiel. Primeiro, a capacidade: $200$ bytes no
\textit{harness} (\texttt{GPS\_DUMP\_DATA\_SIZE}) e $79$ no PX4
(\texttt{sizeof(dump\_data->data)}, em \texttt{msg/GpsDump.msg}). Segundo, os
tipos dos operandos: no \textit{harness}, a subtração é feita em \texttt{int}
(macro \texttt{int} menos \texttt{uint8\_t} promovido), e o resultado negativo é
convertido para \texttt{size\_t}; no PX4, a subtração já é em \texttt{size\_t}.
O mecanismo do \textit{underflow} é análogo nos dois, mas os valores e o
caminho não são idênticos; por isso, o resultado vale para o modelo, e a
conclusão sobre o driver depende da conferência da etapa~(v). Essa limitação,
comum a toda a primeira campanha, motivou a segunda campanha, em que o código
verificado é o original (Seção~\ref{sec:segunda-campanha}).}

Para concretizar a aplicação da metodologia, esta seção percorre
\revisao{as oito etapas} tendo o driver GPS uBlox como objeto de verificação. O driver,
com \revisao{$1.712$ linhas de código} (o \textit{harness}
correspondente tem 294 linhas), é responsável pelo parsing de mensagens
do protocolo UBX recebidas via UART e pela atualização das estruturas de dados de
posição e velocidade consumidas pelo estimador EKF2.

\textbf{Seleção.} O GPS uBlox foi incluído na primeira rodada por sua posição
crítica na cadeia de navegação: coordenadas corrompidas podem comprometer o
Return-to-Launch e outras funções de failsafe que dependem de posição
absoluta~\cite{Wanner2024}.

\textbf{Instrumentação.} \corr{Os bytes que o driver receberia da UART foram
declarados como valores não determinísticos na função de teste e passados ao
trecho transcrito do driver}, simulando tanto pacotes UBX bem formados
quanto dados corrompidos ou truncados. A função \texttt{dumpGpsData()}, que copia
pacotes recebidos para um buffer de diagnóstico, foi incluída no harness por
conter aritmética sobre tipos sem sinal em sua verificação de espaço disponível.

\textbf{Propriedades.} Além das verificações automáticas de segurança de
memória, foi especificada a propriedade~\eqref{eq:prop_gps}: o comprimento do
buffer de diagnóstico não deve superar \texttt{GPS\_DUMP\_DATA\_SIZE} (200
bytes no \textit{harness}; no driver de produção, o registro tem 79 bytes)
após qualquer operação de escrita.

\textbf{Execução.} A verificação foi executada com Boolector 3.2.4,
\texttt{--unwind 10} e \texttt{--overflow-check}, \revisao{concluindo em $64{,}9$
segundos, como mediana de três execuções no servidor descrito na
Seção~\ref{sec:config-experimental}}.

\textbf{Análise.} \revisao{O ESBMC parou na primeira violação --- um
underflow aritmético: quando \texttt{dump\_data->len} supera 200, a expressão
\texttt{GPS\_DUMP\_DATA\_SIZE - dump\_data->len}, em que
\texttt{GPS\_DUMP\_DATA\_SIZE} é um \texttt{\#define} de valor 200 (do tipo
\texttt{int}) e \texttt{len}, um \texttt{uint8\_t} promovido a \texttt{int},
produz um valor negativo ($200 - 207 = -7$ no contraexemplo), que é convertido
para \texttt{size\_t} --- um valor enorme --- na comparação com
\texttt{write\_len}, tornando a guarda
de proteção sempre falsa. O \texttt{memcpy} subsequente copia dados sem
restrição de limite, violando os limites do objeto alocado. No driver
de produção (\texttt{gps.cpp}, linha 741), a expressão correspondente é
\texttt{sizeof(dump\_data->data) - dump\_data->len}, avaliada em
\texttt{size\_t} (sem sinal), sobre um buffer de 79 bytes
(\texttt{msg/GpsDump.msg}).}
\revisao{O contraexemplo, produzido pela função
\texttt{test\_gps\_real\_bit\_operation} (\texttt{test\_choice = 3}), tem
\texttt{input\_len = 10} e \texttt{dump\_buffer.len = 207}. A revisão do
mantenedor do PX4 no PR~\#27576 e a análise do código de produção mostraram
que o estado violador não ocorre no driver real: o laço de escrita limita
\texttt{write\_len} ao espaço livre e zera \texttt{len} após publicar o
registro. Essa inalcançabilidade foi depois verificada formalmente por prova
indutiva (Seção~\ref{subsec:gps-inducao}).}

\textbf{Reporte.} O \revisao{achado} foi reportado à comunidade PX4 via Issue
\#26866\footnote{\url{https://github.com/PX4/PX4-Autopilot/issues/26866}},
acompanhado do contraexemplo completo, e a correção foi proposta como uma
guarda no PR~\#27576, que o mantenedor enquadrou como
endurecimento defensivo~\cite{PX4_GitHub}.

\textbf{Validação cruzada e \textit{patch}.} O CBMC~5.95.1 sinaliza a
mesma violação apenas com \texttt{--unwind} 2, 3 e 5; com \texttt{--unwind 10},
não terminou em 30 minutos. Com as mesmas pré-condições, o ESBMC e o CBMC
acusam a violação sem o \textit{patch} e não a acusam com o \textit{patch} do PR
nem com a forma sugerida pelo mantenedor (Seção~\ref{subsec:patch-gps}).
\corr{Na comparação refeita em 19/10/2026 com o ESBMC~8.5.0 e o tempo-limite de
$600$~s, o ESBMC voltou a encontrar a violação em cerca de $137$~s (mediana de
três execuções) e o CBMC não terminou dentro do limite
(Seção~\ref{subsec:2c-protocolo}).}

\textbf{Integração contínua.} O \textit{harness} do GPS roda no
\textit{workflow} de integração contínua do \textit{fork} a cada \textit{push}
(Seção~\ref{subsec:etapa-ardupilot-ci}). A etapa de reprodução no ArduPilot não
se aplica a este achado, já que ela foi feita apenas para a observação de
temperatura do BMI088.

\section{Seleção dos Módulos para Verificação}
\label{sec:selecao-modulos}

\revisao{O procedimento de seleção combinou duas etapas de natureza distinta, e convém explicitá-las separadamente. A primeira é o \textbf{levantamento do conjunto candidato}, que não foi algorítmico: os módulos foram indicados pela equipe do projeto SAGRES em reunião, com base na experiência operacional da equipe sobre quais componentes são mais determinantes para missões de sensoriamento remoto em áreas remotas. Trata-se, portanto, de elicitação junto a especialistas do domínio, e não de um estudo sistemático de criticidade --- distinção que se registra aqui para que o leitor possa avaliar o alcance da amostra (Seção~\ref{sec:ameacas-validade}, no Capítulo~\ref{cap:experimentos}).}

\revisao{A segunda etapa é a \textbf{filtragem por critérios objetivos}, essa sim sistematizável e reproduzível, formalizada no Algoritmo~\ref{alg:selecao}.} \revisao{Cada candidato indicado foi submetido pela autora a três verificações sequenciais: leitura do código-fonte, para avaliar a complexidade interna e a superfície de entrada não confiável; leitura da documentação e do datasheet, para identificar requisitos de segurança explícitos; e avaliação da viabilidade técnica de isolamento por meio de \textit{stubs}. Apenas os módulos que atendem aos três critérios são incluídos no conjunto final.}

\begin{algorithm}[H]
\caption{Procedimento de seleção de módulos críticos do PX4}
\label{alg:selecao}
\begin{algorithmic}[1]
\Require Repositório PX4-Autopilot (ramo main, commit 8be2590c15), datasheets dos componentes, indicações do coordenador SAGRES
\Ensure Conjunto $M$ de módulos selecionados para verificação
\State $M \gets \emptyset$
\State $\text{candidatos} \gets$ módulos indicados pelo coordenador SAGRES
\ForAll{$m \in \text{candidatos}$}
    \State ler código-fonte de $m$ \Comment{análise da lógica interna}
    \State $\text{superficie}(m) \gets$ pontos de entrada de dado não confiável em $m$
    \State ler datasheet e documentação associados a $m$
    \State $\text{requisitos}(m) \gets$ requisitos de segurança extraídos dos documentos
    \State $\text{isolavel}(m) \gets$ avaliação da viabilidade de abstração via \textit{stubs}
    \If{$\text{superficie}(m) \neq \emptyset \,\land\, \text{requisitos}(m) \neq \emptyset \,\land\, \text{isolavel}(m) = \text{verdadeiro}$}
        \State $M \gets M \cup \{m\}$
    \EndIf
\EndFor
\State \Return $M$
\end{algorithmic}
\end{algorithm}

O tamanho do conjunto candidato indicado pela equipe e os eventuais
descartes em cada critério não foram registrados durante a pesquisa; por isso,
não se pode afirmar quantos candidatos o filtro eliminou, e a etapa de
levantamento do conjunto candidato não é inteiramente reproduzível. Registra-se
essa limitação junto com outra, apontada também no artigo derivado: selecionar
módulos pela viabilidade de isolamento por \textit{stubs} sub-representa os
módulos com forte acoplamento ao \textit{hardware} ou ao escalonador, de modo
que a amostra não é estatisticamente representativa do PX4
(Seção~\ref{sec:ameacas-validade}).

A aplicação do procedimento ao PX4 Autopilot (ramo main, commit 8be2590c15) resultou em um conjunto de 12 módulos, organizados em três rodadas conforme descrito na Seção~\ref{sec:visao-geral}. Os quatro módulos da primeira rodada foram selecionados por sua criticidade direta para a segurança de voo: drivers de sensores (IMU BMI088 e GPS uBlox), \corr{comunicação (MAVLink, \textit{buffer} de transmissão)} e controle de entrada (Expo RC). As rodadas subsequentes ampliam a cobertura para monitoramento de sistema, navegação e estimação de estado.

%% LC 19 (resolvido 19/10/2026): citacao deSorbo2022 removida.
\corr{Foram priorizados componentes cuja falha pode resultar na perda de
controle do VANT ou na execução incorreta de missões autônomas.}
Os resultados quantitativos da verificação de cada módulo \revisao{(veredito, achados e tempo de execução)} são apresentados no Capítulo~\ref{cap:experimentos} (Experimentos e Resultados), juntamente com a discussão dos contraexemplos e das decisões de configuração do ESBMC específicas para cada módulo.

O driver IMU BMI088 integra acelerômetro e
giroscópio via SPI~\cite{PX4_Documentation}, realizando a leitura de
registradores, a conversão de dados brutos para unidades físicas (m/s² e rad/s),
 a aplicação de calibração e a leitura de temperatura. Dados incorretos da IMU
propagam-se para a estimação de atitude e podem causar instabilidade de
voo~\cite{Siddiqi2022}. O driver GPS uBlox realiza o parsing de mensagens do
protocolo UBX via UART~\cite{PX4_Documentation}; falhas no parser podem
corromper as coordenadas e comprometer funções de failsafe, como Return-to-Launch.
%% LC 20 (resolvido 19/10/2026): renomeado "MAVLink (buffer de transmissao)".
\corr{O módulo MAVLink implementa o protocolo de comunicação entre o autopiloto
e a estação de controle~\cite{Ficco2022, Hamza2024}. O \textit{harness} deste
trabalho não verifica o \textit{parser} de mensagens recebidas, e sim o
\textit{buffer} de transmissão (\texttt{send\_bytes}, em
\texttt{src/modules/mavlink/mavlink\_main.cpp}): o preenchimento não deve exceder
o tamanho máximo de pacote. Por isso, o módulo é referido neste texto como
``MAVLink (\textit{buffer} de transmissão)''.}
%% LC 21 (resolvido 19/10/2026): citacao trocada pelo arquivo-fonte.
\corr{O módulo Expo implementa a curva exponencial aplicada aos comandos do
controle remoto, com aritmética de ponto flutuante; a função verificada é
\texttt{math::expo()}, definida em
\texttt{src/lib/mathlib/math/Functions.hpp} no commit \texttt{8be2590c15}.}
Os módulos das rodadas seguintes ampliam a
cobertura para monitoramento, navegação e estimação de estado. A
Figura~\ref{fig:px4-modular} mostra a organização modular em que eles se
inserem.

\begin{figure}[htb]
    \centering
    \resizebox{\textwidth}{!}{%
    \begin{tikzpicture}[
      font=\linespread{1}\footnotesize, >={Stealth[length=2mm]},
      ver/.style={draw=palEdgeDark, fill=palStrong, very thick, rounded corners=2pt,
        align=center, text width=3.2cm, minimum height=9mm},
      nver/.style={draw=palEdge, fill=palPaper, rounded corners=2pt, align=center,
        text width=3.2cm, minimum height=9mm},
      bus/.style={draw=palEdgeDark, very thick, fill=palFill, rounded corners=3pt, align=center,
        text width=2.6cm, minimum height=3.6cm},
      tit/.style={font=\small\bfseries}]
      \node[bus] (u) {\textbf{uORB}\\(barramento de mensagens\\\textit{publish--subscribe})};
      \node[ver, left=22mm of u, yshift=24mm] (imu) {\textbf{IMU BMI088}};
      \node[ver, below=3mm of imu] (gps) {\textbf{GPS uBlox}};
      \node[ver, below=3mm of gps] (mag) {\textbf{Magnetômetro}\\e \textbf{Barômetro}};
      \node[ver, below=3mm of mag] (rc) {\textbf{RC Update}\\e \textbf{Expo RC}};
      \node[ver, below=3mm of rc] (bat) {\textbf{Battery Monitor}};
      \node[tit, above=3mm of imu] {Publicadores};
      \node[ver, right=22mm of u, yshift=24mm] (ekf) {\textbf{EKF2}\\(covariância)};
      \node[ver, below=3mm of ekf] (nav) {\textbf{Navigator/Mission}};
      \node[nver, below=3mm of nav] (ctl) {Controle de posição\\e de atitude};
      \node[ver, below=3mm of ctl] (vtol) {\textbf{VTOL}};
      \node[ver, below=3mm of vtol] (pwm) {\textbf{PWM Output}};
      \node[tit, above=3mm of ekf] {Assinantes};
      \foreach \n in {imu,gps,mag,rc,bat} \draw[->] (\n.east) -- (u.west);
      \foreach \n in {ekf,nav,ctl,vtol,pwm} \draw[->] (u.east) -- (\n.west);
      \node[ver, below=10mm of u, text width=3.6cm] (mav) {\textbf{MAVLink}\\(enlace com a estação de solo)};
      \draw[<->] (u.south) -- (mav.north);
      \node[left=1mm of u.north west, anchor=south east, font=\scriptsize\itshape] {\texttt{orb\_publish()}};
      \node[right=1mm of u.north east, anchor=south west, font=\scriptsize\itshape] {\texttt{orb\_copy()}};
      \node[ver, text width=4mm, minimum height=4mm, below=6mm of mav.south west, xshift=-52mm, anchor=north west] (l1) {};
      \node[right=1mm of l1, anchor=west, font=\scriptsize] (t1) {módulo verificado neste trabalho (12 módulos)};
      \node[nver, text width=4mm, minimum height=4mm, right=6mm of t1] (l2) {};
      \node[right=1mm of l2, anchor=west, font=\scriptsize] {módulo não verificado};
    \end{tikzpicture}}
    \caption[Organização modular do PX4 Autopilot]{Organização modular do PX4 Autopilot, com destaque para os doze módulos da primeira campanha. Elaborado pela autora com base em~\cite{PX4_Documentation}.}
    \label{fig:px4-modular}
\end{figure}

\section{Instrumentação e Preparação do Código}
\label{sec:instrumentacao}

A verificação formal de firmware embarcado apresenta desafios relacionados a
dependências de hardware. O PX4 utiliza intensivamente a HAL do NuttX e as bibliotecas de comunicação I2C, SPI e UART~\cite{PX4_Documentation}. Executar o verificador sobre o código original resultaria em falhas de compilação pela ausência dessas dependências no ambiente de verificação.

%% LC 22 (resolvido 19/10/2026): explicacao de nao determinismo removida daqui.
\corr{A Etapa~(ii) descreveu como essas dependências são abstraídas por valores
não determinísticos. Esta seção mostra como isso aparece nos \textit{harnesses}
da primeira campanha; a instrumentação da segunda campanha, com \textit{stubs}
de fato, é descrita na Seção~\ref{sec:segunda-campanha}.}

Nos \textit{harnesses} depositados, a abstração da E/S de
\textit{hardware} fica na própria função de teste: os bytes que o driver leria do
barramento (SPI, I2C ou UART) são declarados como valores não determinísticos e
passados diretamente ao trecho do driver sob verificação. A
Listagem~\ref{lst:stub_i2c} mostra a função \texttt{updateTemperature()} como
transcrita no \textit{harness} do IMU, e a Listagem~\ref{lst:main_imu}, a função
de teste que a exercita.

\begin{lstlisting}[
  caption={Função \texttt{updateTemperature()} transcrita no \textit{harness} do IMU BMI088 (\texttt{test\_bmi088\_imu.cpp}, comentários omitidos)},
  label={lst:stub_i2c},
  numbers=left,
  frame=single,
  language=C,
  breaklines=true,
  basicstyle=\ttfamily\small]
float updateTemperature(uint8_t temp_msb, uint8_t temp_lsb) {
    uint16_t Temp_uint11 = (temp_msb * 8) + (temp_lsb / 32);
    int16_t Temp_int11 = 0;

    if (Temp_uint11 > 1023) {
        Temp_int11 = Temp_uint11 - 2048;
    } else {
        Temp_int11 = Temp_uint11;
    }

    float temperature = (Temp_int11 * 0.125f) + 23.0f;
    return temperature;
}
\end{lstlisting}

\begin{lstlisting}[
  caption={Função de teste da temperatura no \textit{harness} do IMU BMI088 (\texttt{test\_bmi088\_imu.cpp}, comentários omitidos)},
  label={lst:main_imu},
  numbers=left,
  frame=single,
  language=C,
  breaklines=true,
  basicstyle=\ttfamily\small]
void test_temperature_calculation() {
    uint8_t temp_msb = nondet_uint8();
    uint8_t temp_lsb = nondet_uint8();

    float temperature = updateTemperature(temp_msb, temp_lsb);

    assert(temperature >= -40.0f && temperature <= 85.0f);

    assert(!isnan(temperature));
    assert(!isinf(temperature));
}
\end{lstlisting}

A função da Listagem~\ref{lst:stub_i2c} é uma transcrição do trecho
correspondente do driver BMI088 do PX4, copiada para dentro do
\textit{harness}, e não a função compilada a partir do código de produção. É assim
em todos os \textit{harnesses} \corr{da primeira campanha}: cada um extrai trechos do código-fonte do PX4,
que não foi modificado. Isso tem uma consequência que se registra aqui e nas
ameaças à validade (Seção~\ref{sec:ameacas-validade}): o que se verifica é a
cópia da lógica, e uma guarda presente no driver e omitida no trecho transcrito
produz um falso positivo --- exatamente o mecanismo do candidato do contador de
FIFO (Seção~\ref{subsec:imu-fifo}). É por isso que a conferência contra o código
de produção, na etapa~(v), é obrigatória. \corr{Essa limitação é a principal
motivação da segunda campanha.}

\section{Propriedades Verificadas}
\label{sec:propriedades}

As propriedades de segurança foram organizadas em duas categorias: automáticas,
verificadas pelo ESBMC a partir das flags de configuração, e customizadas,
expressas por assertions inseridas no código.

\subsection{Propriedades Automáticas}

O ESBMC verifica automaticamente diversas classes de defeitos quando habilitado
pelas flags de configuração. Seguindo a formalização de Cordeiro et al.~\cite{Cordeiro2012}, uma propriedade de safety $\phi$ é violada quando existe um estado alcançável $s_i$ tal que $s_i \not\models \phi$. \revisao{Em segurança de memória, o verificador detecta acessos fora dos limites de arranjos e desreferências de ponteiros nulos. Em aritmética inteira, identifica \textit{overflows} e \textit{underflows} de inteiros com sinal. Verifica ainda divisões por zero}~\cite{Menezes2024}.

A Tabela~\ref{tab:propriedades-automaticas} detalha as principais classes de defeitos verificadas automaticamente pelo ESBMC nesta pesquisa, suas descrições e as \textit{flags} de linha de comando associadas, conforme a documentação da ferramenta~\cite{Menezes2024}.

\begin{table}[!htb]
\centering
\caption{Propriedades automáticas verificadas pelo ESBMC neste trabalho}
\label{tab:propriedades-automaticas}
\footnotesize
\begin{tabular}{|p{3cm}|p{7cm}|p{3cm}|}
\hline
\textbf{Classe} & \textbf{Descrição} & \textbf{Flag ESBMC} \\
\hline
Acesso fora dos limites & Leitura/escrita em endereços inválidos de arrays e buffers & padrão (\texttt{--no-bounds-check} desliga) \\
\hline
Ponteiro nulo & Desreferência de ponteiro igual a \texttt{NULL} & padrão (\texttt{--no-pointer-check} desliga) \\
\hline
Overflow aritmético & Operações com inteiros com sinal que excedem a capacidade do tipo & \texttt{--overflow-check} \\
\hline
Divisão por zero & Operações de divisão ou módulo com divisor nulo & padrão (\texttt{--no-div-by-zero-check} desliga) \\
\hline
\corr{Deslocamento indefinido} & \corr{Deslocamento de bits com quantidade negativa ou maior que a largura do tipo (só na segunda campanha)} & \corr{\texttt{--ub-shift-check}} \\
\hline
\corr{Valores NaN} & \corr{Operação de ponto flutuante que produz NaN (só nos testes indicados da segunda campanha)} & \corr{\texttt{--nan-check}} \\
\hline
\end{tabular}

\footnotesize
\begin{flushleft}
\revisao{Acesso fora dos limites, ponteiro nulo e divisão por zero são verificados por padrão pelo ESBMC~8.1, que só oferece opções para desligá-los; a
verificação de \textit{overflow} aritmético é habilitada explicitamente na
invocação (Listagem~\ref{lst:comando-esbmc}).} \corr{As duas últimas linhas
valem apenas para a segunda campanha, com o ESBMC~8.5.}
\end{flushleft}
\end{table}

As verificações de NaN (\texttt{--nan-check}) e de vazamento de memória
(\texttt{--memory-leak-check}) não foram habilitadas na invocação da
Listagem~\ref{lst:comando-esbmc}, como no artigo derivado. O enunciado de que
os \textit{harnesses} de dez módulos não apresentam violações vale, portanto, para
as quatro classes da tabela e para as \textit{assertions} customizadas. NaN e
infinito são verificados apenas onde o \textit{harness} tem \textit{assertion}
explícita (\texttt{!isnan}, \texttt{!isinf}), o que ocorre em oito dos doze
módulos: IMU BMI088, Expo RC, EKF2, magnetômetro, barômetro, bateria, RC Update e
VTOL.

\revisao{Essas verificações são aplicadas uniformemente a todos os 12 módulos
analisados, sem necessidade de instrumentação específica por módulo. A flag
\texttt{--overflow-check} foi habilitada em todas as execuções, dado o uso
intensivo de aritmética inteira nos drivers de sensores do PX4.}
Ela não cobre o \textit{wraparound} de inteiros sem sinal, que a norma da
linguagem define e que só seria sinalizado com \texttt{--unsigned-overflow-check},
não usada nesta campanha, nem a conversão de um valor negativo para um tipo sem
sinal. Por isso, o \textit{underflow} do GPS não aparece como
\textit{overflow} aritmético, e sim como acesso fora dos limites dentro do
\texttt{memcpy} (Seção~\ref{sec:exemplo-ilustrativo}).

\subsection{Assertions Customizadas}

Além das verificações automáticas, foram adicionadas assertions baseadas em
especificações de fabricantes e requisitos do PX4~\cite{PX4_Documentation}.
Formalmente, cada assertion customizada expressa uma propriedade
$\phi_i$ da forma $\phi_i(s) \equiv \texttt{expr}$, onde \texttt{expr} é uma
condição booleana sobre o estado $s$ do programa~\cite{Cordeiro2012}.

Cada propriedade customizada é apresentada em três níveis de abstração: (i) \textbf{especificação em linguagem natural}, extraída do datasheet do fabricante ou da documentação do protocolo; (ii) \textbf{formalização}, expressando a propriedade como predicado lógico sobre o estado do programa; e (iii) \textbf{implementação}, por meio de \textit{assertion} inserida no \textit{harness} de verificação.

\textbf{Propriedade 1 --- Faixa operacional de temperatura do BMI088.}

\textit{Especificação em linguagem natural:} conforme o datasheet do fabricante Bosch Sensortec~\cite{Bosch2021}, a temperatura lida do sensor de temperatura integrado ao BMI088 deve estar contida na faixa operacional do componente, definida entre $-40\,^{\circ}$C e $+85\,^{\circ}$C. Valores fora dessa faixa indicam leitura inválida e não devem ser utilizados como entrada para os cálculos de compensação térmica do EKF2.

\textit{Formalização:}

\begin{equation}
\phi_{\text{temp}}(s) \equiv -40{,}0 \leq \texttt{temperature}(s) \leq 85{,}0
\label{eq:prop_temp}
\end{equation}

\revisao{Registre-se desde já que a passagem de uma faixa operacional de componente
físico para uma restrição sobre variável de programa é um passo de modelagem
contestável: o \textit{datasheet} declara em que faixa o sensor opera dentro da
especificação, e não que o software esteja proibido de representar uma leitura
fora dela. Foi precisamente esse passo que os mantenedores do PX4 recusaram, e a
discussão consta da Seção~\ref{subsec:imu-temperatura}. A propriedade é mantida
aqui como exemplo justamente por isso: ela ilustra a categoria de propriedade
cuja tradução exige mais cuidado, discutida na
Seção~\ref{subsec:proveniencia}.} Como a conversão afim do driver não
limita o valor, qualquer leitura bruta fora da faixa viola $\phi_{\text{temp}}$
por construção; por isso, como no artigo derivado, $\phi_{\text{temp}}$ é tratada
não como garantia do \textit{hardware}, mas como uma sonda de cobertura que
sinaliza a ausência de validação da faixa de entrada no driver.

\textit{Implementação:}

\begin{lstlisting}[
  numbers=left,
  frame=single,
  language=C,
  basicstyle=\ttfamily\small]
float temperature = updateTemperature(temp_msb, temp_lsb);
assert(temperature >= -40.0f && temperature <= 85.0f);
\end{lstlisting}

\textbf{Propriedade 2 --- Integridade do buffer de diagnóstico do GPS uBlox.}

\textit{Especificação em linguagem natural:} o nível de preenchimento do buffer de diagnóstico usado pela função \texttt{dumpGpsData()} deve permanecer dentro da capacidade do registro. No driver de produção, essa capacidade é \texttt{sizeof(dump\_data->data)}, 79 bytes (\texttt{msg/GpsDump.msg}); no \textit{harness}, ela é modelada pela macro \texttt{GPS\_DUMP\_DATA\_SIZE}, de valor 200, que não existe no código do PX4. Após qualquer operação de escrita no buffer, o comprimento acumulado \texttt{dump\_data.len} não deve exceder esse limite, sob pena de acesso fora dos limites da memória alocada.

\textit{Formalização:}

\begin{equation}
\phi_{\text{gps}}(s) \equiv \texttt{dump\_data.len}(s) \leq
\texttt{GPS\_DUMP\_DATA\_SIZE}
\label{eq:prop_gps}
\end{equation}

Como \texttt{len} é entrada livre, verificar $\phi_{\text{gps}}$ na entrada da
função seria trivialmente violado pela própria entrada. A verificação que importa
é a do acesso no ponto de uso, o \texttt{memcpy}: a violação reportada pelo ESBMC
é um acesso fora dos limites dentro do \texttt{memcpy} (\textit{dereference
failure}), e não a \textit{assertion} abaixo, que é a propriedade de mais alto
nível do \textit{harness}. \corr{Como discutido na
Seção~\ref{sec:exemplo-ilustrativo}, essa \textit{assertion} é tautológica na
forma depositada.}

\textit{Implementação:}

\begin{lstlisting}[
  numbers=left,
  frame=single,
  language=C,
  basicstyle=\ttfamily\small]
dumpGpsData(input_data, input_len, gps_dump_comm_mode_t::Full,
            /* msg_to_gps_device = */ true,
            &dump_buffer, gps_dump_comm_mode_t::Full);
assert(dump_buffer.len == 0 ||
       dump_buffer.len <= GPS_DUMP_DATA_SIZE ||
       (dump_buffer.len & 0x7F) <= GPS_DUMP_DATA_SIZE);
\end{lstlisting}

\textbf{\revisao{Propriedade 3 --- Ganho de Kalman no intervalo unitário.}}

\revisao{\textit{Especificação em linguagem natural:} no filtro de Kalman
implementado pelo módulo de fusão do magnetômetro, o ganho $K$ pondera quanto da
inovação da medição é incorporado à estimativa.} No \textit{harness}, ele
é calculado como $K = P/(P + R)$, em que $P$ é a variância do estado e $R$, a da
observação.  \revisao{Por definição, trata-se de um
fator de ponderação entre a predição e a observação, e portanto deve permanecer
no intervalo $[0;\,1]$. Valor fora desse intervalo indica erro numérico na
atualização da covariância e faz a estimativa divergir da observação em vez de
convergir para ela.}

\revisao{Diferentemente das duas anteriores, esta propriedade não deriva de
\textit{datasheet} nem de constante do código-fonte, mas da definição matemática
do estimador. É, por isso, a categoria de propriedade mais robusta das três:
independe de fidelidade entre documento e implementação e não envolve traduzir
grandeza física em variável de programa.}

\revisao{\textit{Formalização:}}

\begin{equation}
\phi_{\text{kalman}}(s) \equiv 0 \leq K(s) \leq 1
\label{eq:prop_kalman}
\end{equation}

\revisao{\textit{Implementação:}}

\begin{lstlisting}[
  numbers=left,
  frame=single,
  language=C,
  basicstyle=\ttfamily\small]
float compute_kalman_gain(float state_var, float obs_var)
{
    if (state_var + obs_var <= 0.0f) return 0.0f;
    return state_var / (state_var + obs_var);
}

float gain = compute_kalman_gain(state_var, obs_var);
assert(gain >= 0.0f);
assert(gain <= 1.0f);
assert(!isnan(gain));
\end{lstlisting}

As variâncias de entrada são não determinísticas, restritas por
\textit{assumptions} a valores finitos, não negativos e com soma positiva.

%% LC 24 (resolvido 19/10/2026): funcao escrita no harness, modela
%% VehicleMagnetometer.cpp.
\corr{A função \texttt{compute\_kalman\_gain} não é código do PX4: ela foi escrita
no próprio \textit{harness} (\texttt{test\_vehicle\_magnetometer.cpp}) para
modelar o cálculo do ganho feito em
\texttt{src/modules/sensors/vehicle\_magnetometer/VehicleMagnetometer.cpp}.
O veredito desta propriedade vale, portanto, para o modelo escrito no
\textit{harness}, e não diretamente para o código do PX4.}

No MAVLink, verificou-se que o preenchimento do buffer de transmissão
nunca excede sua capacidade, de 280 bytes, o tamanho máximo de pacote do
protocolo, mesmo com pacotes de entrada maiores que ela e após envios
sucessivos.

\subsection{\revisao{Proveniência das Propriedades}}
\label{subsec:proveniencia}

\revisao{A rastreabilidade entre cada propriedade verificada e o documento que a
originou é parte da metodologia: sem ela, não há como um terceiro auditar se a propriedade
traduz corretamente o requisito, nem como a comunidade avaliar um contraexemplo. A
Tabela~\ref{tab:proveniencia} registra essa correspondência para todas as propriedades
customizadas deste trabalho.}

%% LC 25 (resolvido parcialmente 19/10/2026): arquivo-fonte e funcao do harness
%% acrescentados; numeros de linha do PX4 e pagina do datasheet seguem pendentes.
%% LC 26 (resolvido 19/10/2026): Battery Monitor e Barometer acrescentados.
\begin{table}[!htb]
\centering
\caption[Proveniência das propriedades customizadas (primeira campanha)]{\revisao{Proveniência das propriedades customizadas verificadas} \corr{(primeira campanha)}}
\label{tab:proveniencia}
\footnotesize
\begin{tabular}{|p{3.0cm}|p{2.2cm}|p{2.5cm}|p{6.0cm}|}
\hline
\textbf{Propriedade} & \textbf{Módulo} & \textbf{Origem} & \textbf{Referência exata} \\
\hline
Faixa operacional de temperatura & IMU BMI088 & Datasheet do fabricante &
Bosch Sensortec, BST-BMI088-DS000-19, rev.~1.9~\cite{Bosch2021}; \corr{\textit{harness}: \texttt{test\_temperature\_calculation}} \\
\hline
Capacidade do buffer FIFO & IMU BMI088 & Código-fonte do PX4 &
Constante \texttt{FIFO::SIZE = 1024} em \texttt{Bosch\_BMI088\_ Accelerometer\_Registers.hpp}, \corr{l.~152}; \corr{\textit{harness}: \texttt{test\_fifo\_count\_calculation}} \\
\hline
Limite do buffer de diagnóstico & GPS uBlox & Código-fonte do PX4 &
\texttt{sizeof(dump\_data->data)}, 79 bytes, em \texttt{msg/GpsDump.msg}; no \textit{harness}, macro \texttt{GPS\_DUMP\_ DATA\_SIZE = 200}; \corr{\texttt{src/drivers/gps/gps.cpp}, l.~741} \\
\hline
Capacidade do buffer de transmissão & MAVLink & Especificação de protocolo &
Tamanho máximo de pacote MAVLink, 280 bytes (\texttt{MAVLINK\_MAX\_ PACKET\_LEN})~\cite{PX4_Documentation}; \corr{\texttt{Mavlink::send\_bytes} em \texttt{src/modules/mavlink/mavlink\_main.cpp}, l.~852} \\
\hline
Faixa de entrada do rádio & Expo RC, RC Update & Documentação do PX4 &
Intervalo normalizado $[-1;\,1]$~\cite{PX4_Documentation}; \corr{\texttt{math::expo()} em \texttt{src/lib/mathlib/math/Functions.hpp}, l.~84; \texttt{src/modules/rc\_update/rc\_update.cpp}} \\
\hline
Limites de PWM & PWM Output & Documentação do PX4 &
Faixas de motor e servo~\cite{PX4_Documentation}; \corr{\texttt{src/drivers/pwm\_out/PWMOut.cpp}; constantes \texttt{PWM\_MOTOR\_MIN}/\texttt{PWM\_MOTOR\_MAX} do \textit{harness}} \\
\hline
Variância não negativa & EKF2, Magnetômetro & Requisito matemático &
Propriedade de matriz de covariância; \corr{\texttt{src/modules/ekf2/EKF/covariance.cpp}} \\
\hline
Ganho de Kalman em $[0;\,1]$ & Magnetômetro & Requisito matemático &
Definição do filtro de Kalman; \corr{modelado no \textit{harness} a partir de \texttt{VehicleMagnetometer.cpp}} \\
\hline
Índice de missão válido & Navigator & Código-fonte do PX4 &
Limites de \textit{waypoint}; \corr{\texttt{src/modules/navigator/mission.cpp}} \\
\hline
\textit{Thrust} em $[0;\,0{,}9]$ & VTOL & Documentação do PX4 &
Limite de transição~\cite{PX4_Documentation}; \corr{\texttt{src/modules/vtol\_att\_control/vtol\_type.cpp}} \\
\hline
\corr{Fator de escala de tensão em $[1{,}0;\,1{,}3]$; avisos e falhas coerentes com os limiares} & \corr{Battery Monitor} & \corr{Código-fonte do PX4} & \corr{\texttt{src/lib/battery/battery.cpp} (\texttt{determineWarning}, \texttt{determineFaults}, \texttt{computeScale})} \\
\hline
\corr{Prioridade em $[0;\,100]$; pressão corrigida finita e em $[30000;\,120000]$ Pa} & \corr{Barômetro} & \corr{Código-fonte do PX4} & \corr{\texttt{src/lib/sensor\_calibration/Barometer.cpp}} \\
\hline
\end{tabular}
\end{table}

As \textit{assertions} de cada linha da tabela estão no
\textit{harness} do módulo correspondente, no repositório de artefatos, em
\texttt{resultados/<módulo>/test\_*.cpp}: \texttt{IMU\_BMI088},
\texttt{GPS}, \texttt{MAVLink}, \texttt{Expo\_RC}, \texttt{RC\_Update},
\texttt{PWM\_Output}, \texttt{EKF2\_Covariance}, \texttt{Magnetometer},
\texttt{Navigator\_Mission}, \texttt{VTOL}\corr{, \texttt{Battery\_Monitor} e
\texttt{Barometer}.}

\corr{No \textit{datasheet} do BMI088 (BST-BMI088-DS000-19, rev.~1.9, janeiro de 2024), a faixa operacional de $-40$ a $+85\,^{\circ}$C está na Seção~1.1, Tabela~1 (parâmetro \textit{Operating Temperature}, $T_A$), p.~7; a conversão do valor lido para graus Celsius, com resolução de $0{,}125\,^{\circ}$C por unidade, está na Seção~5.3.7 (registradores \texttt{0x22}--\texttt{0x23}), p.~28.}

\revisao{A distinção entre as três origens não é apenas documental: ela determina a
força da conclusão. Propriedades derivadas do \textbf{código-fonte} são as mais seguras, pois
a constante verificada é a mesma que o compilador usa. Propriedades derivadas da
\textbf{documentação ou de especificação de protocolo} dependem da fidelidade entre documento
e implementação. Propriedades derivadas de \textbf{datasheet} são as mais delicadas, porque o
documento descreve o comportamento do componente físico, e não necessariamente uma restrição
sobre a variável de software que o representa --- ponto discutido em detalhe na
Seção~\ref{subsec:imu-temperatura} do Capítulo~\ref{cap:experimentos}.}

\section{Análise de Contraexemplos}
\label{sec:contraexemplos}

%% LC 27 (resolvido 19/10/2026): secao reduzida a escala de classificacao.
%% LC 28 (resolvido 19/10/2026): falha reportada e limite do memcpy.
\corr{O procedimento de análise foi descrito na etapa~(v)
(Seção~\ref{subsec:etapa-replay}). Esta seção registra apenas a escala usada
para classificar os resultados, fixada no Capítulo~\ref{cap:introducao}:}
\textbf{violação alcançável no modelo},
quando a violação decorre dos estados admitidos pelas pré-condições do
\textit{harness} (caso do GPS); e
\textbf{falso positivo}, quando o candidato decorre de guarda presente no código
de produção e omitida no \textit{harness} (caso do contador de FIFO). A
observação de temperatura ficou fora dessas duas categorias: a violação é real
no modelo, mas a própria propriedade, derivada do \textit{datasheet}, foi
contestada pelos mantenedores. \corr{No caso do GPS, a falha reportada pelo
ESBMC não é a violação da \textit{assertion} da propriedade~\eqref{eq:prop_gps},
e sim a verificação automática de limites dentro do \texttt{memcpy}
(\textit{dereference failure}), consequência do \textit{underflow} descrito na
Seção~\ref{sec:exemplo-ilustrativo}.}

\corr{Na segunda campanha, essa mesma escala foi aplicada, com uma distinção que
a primeira já exigia: um \textbf{falso positivo de modelagem} decorre de uma
guarda de produção omitida no \textit{harness} (o contador de FIFO na primeira
campanha; o \texttt{bst} na segunda), ao passo que um \textbf{falso positivo da
ferramenta} decorre do verificador, e não do código. Os desfechos da segunda
campanha e os valores concretos de cada contraexemplo são apresentados no
Capítulo~\ref{cap:experimentos} e no Apêndice~\ref{apend:contraexemplos}.}

\section{\corr{Segunda Campanha: Código de Produção, \textit{Replay} e Sanitizadores}}
\label{sec:segunda-campanha}

\corr{A primeira campanha mostrou a limitação de verificar trechos transcritos:
o veredito vale para a cópia, e a conclusão sobre o driver depende de uma
conferência que, nos casos do FIFO e do GPS, inverteu o resultado. A segunda
campanha foi desenhada para reduzir essa distância. Ela mantém as oito etapas,
mas endurece a conferência contra o código de produção (etapa~v): cada candidato
alcançável é submetido a \textit{replay} do \emph{código real} do PX4 sob
sanitizadores dinâmicos (ASan/UBSan), e a validação cruzada com o CBMC passa a
cobrir todos os \textit{harnesses}, não apenas os candidatos. Dois achados
desta campanha foram reportados à comunidade em 12/11/2026, depois da revisão do
orientador: o do RoboClaw (\textit{Issue}~\#29785) e o do batmon
(\textit{Issue}~\#29787). Ambos aguardavam triagem na data de fechamento deste
texto.}

\corr{A campanha reúne dez módulos complementares aos doze da primeira, em três
frentes: (A)~oito candidatos adicionais em módulos de entrada e de sistema;
(B)~o driver de saída \textbf{RoboClaw}; e (C)~o driver de bateria
\textbf{batmon}. As frentes compartilham o protocolo de evidência --- verificação
com ESBMC e CBMC, \textit{replay} do código real e sanitizadores ---, mas
diferem no commit e no desfecho, como se detalha a seguir.}

\subsection{\corr{Frente A --- Oito candidatos adicionais}}
\label{subsec:2c-candidatos}

\corr{Oito módulos foram verificados no mesmo commit da primeira campanha
(\texttt{8be2590c15}), com ESBMC~8.1.0 e~8.5.0, CBMC~5.95.1 e os sanitizadores de
endereço (ASan) e de comportamento indefinido (UBSan). Um único script
(\texttt{rodar\_tudo.sh}) regenera todos os registros em qualquer máquina com as
três ferramentas, cada log com cabeçalho (data, máquina, pasta, comando e
versões) e código de saída diferente de zero se algum resultado divergir do
esperado. A Tabela~\ref{tab:candidatos-extras} resume os desfechos.}

\begin{table}[!htb]
\centering
\caption{\corr{Candidatos adicionais da segunda campanha (commit \texttt{8be2590c15})}}
\label{tab:candidatos-extras}
\scriptsize
\renewcommand{\arraystretch}{0.95}
\setlength{\tabcolsep}{3pt}
\begin{tabular}{|p{2.4cm}|p{5.2cm}|p{5.0cm}|}
\hline
\textbf{Módulo} & \textbf{Resultado no commit testado} & \textbf{Situação e desfecho} \\
\hline
\corr{\texttt{sumd.cpp}} & \corr{com quadro legítimo de 32 canais, \texttt{sumd\_decode()} escreve o 64.º byte em \texttt{sumd\_data[64]}, sobre \texttt{crc16\_high}} & \corr{\textbf{continua no \texttt{main}}; defeito reproduzido por \textit{replay} da função real; relatado por terceiro em 08/10/2026 (PR~\#29031, refeito como \#29117), fechado sem \textit{merge}; não reportado de novo} \\
\hline
\corr{\texttt{MicroStrain.cpp}} & \corr{construtor: \texttt{max\_len} pode ser 128 e \texttt{\_port[128]='\textbackslash0'} escreve fora do vetor de 128} & \corr{\textbf{continua no \texttt{main}}; defeito reproduzido (ASan); não reportado, pelo impacto desprezível (exige caminho de dispositivo com 128 caracteres ou mais)} \\
\hline
\corr{\texttt{externalChecks.cpp}} & \corr{\texttt{num\_events} (0--255) indexa \texttt{events[5]}} & \corr{era real; \textbf{corrigido \textit{upstream}} em \texttt{ca1d36476a} (PR~\#28598)} \\
\hline
\corr{\texttt{ft7\_technologies.cpp}} & \corr{contadores de byte indexam \texttt{\_readbuf[30]} sem limite} & \corr{era real; \textbf{corrigido \textit{upstream}} em \texttt{3acf7935a0} (PR~\#28600)} \\
\hline
\corr{\texttt{IridiumSBD.cpp}} & \corr{a guarda de espaço compara com \texttt{\_packet\_length}, mas o \texttt{memcpy} copia \texttt{buflen}} & \corr{era real; \textbf{corrigido \textit{upstream}} em \texttt{0c903d47d6} (PR~\#28599)} \\
\hline
\corr{\texttt{bst.cpp}} & \corr{violação apenas no \textit{harness} que omitia duas guardas já presentes no commit (de \texttt{e8e86a2e0f})} & \corr{\textbf{falso positivo de modelagem}; o \textit{harness} fiel verifica \textsc{successful}} \\
\hline
\corr{\texttt{logger.cpp}} & \corr{nenhuma violação de memória em \texttt{write\_info\_multiple()} para qualquer \texttt{key\_len}} & \corr{sem defeito de memória; o truncamento é alcançável, mas gera registro com \texttt{key\_len} errado, não escrita fora dos limites} \\
\hline
\corr{\texttt{Common.cpp}} & \corr{violação só com índice que a guarda de produção nunca permite} & \corr{\textbf{retirado}} \\
\hline
\end{tabular}
\end{table}

\corr{Dois pontos merecem registro. Primeiro, três candidatos
(\texttt{externalChecks}, \texttt{ft7} e \texttt{IridiumSBD}) eram violações reais
no commit testado e foram \emph{corrigidos pelo próprio PX4}, por outras pessoas,
em setembro de 2026; as correções \textit{upstream} corroboram, de forma
independente, que os achados eram defeitos reais, e por isso nenhuma \textit{issue}
foi aberta para eles. O \texttt{sumd} teve trajetória parecida: o mesmo defeito
foi relatado por outra pessoa, mas o PR foi fechado sem \textit{merge}, e o código
permanece igual no \texttt{main}. Segundo, a verificação com o CBMC usou as mesmas
pré-condições e limites de desdobramento do ESBMC, com \textit{unwinding
assertions} e \textit{bounds} que cobrem os laços (por exemplo, 260 no
\texttt{externalChecks} e 345 no \texttt{IridiumSBD}); as duas ferramentas
acusam as mesmas violações. O caso do \texttt{sumd} precisou de
\texttt{--unwind 70}; como o modo \texttt{--multi-property} do ESBMC ultrapassa
quinze minutos nesse limite, a contagem de propriedades foi tomada do CBMC.}

\subsection{\corr{Frente B --- Driver de saída RoboClaw}}
\label{subsec:2c-roboclaw}

\corr{O driver \texttt{roboclaw} converte os comandos de controle em pacotes do
protocolo \textit{Packet Serial} enviados por UART aos controladores de motor.
Dois defeitos foram verificados no commit \texttt{8be2590c15}, com ESBMC~8.2.0
(solver Z3) e CBMC~5.95.1 (Tabela~\ref{tab:roboclaw-resultados}).}

\begin{table}[!htb]
\centering
\caption{\corr{RoboClaw: verificação formal dos dois defeitos}}
\label{tab:roboclaw-resultados}
\footnotesize
\begin{tabular}{|p{7.6cm}|p{2.6cm}|p{2.6cm}|}
\hline
\textbf{Defeito} & \textbf{Sem correção} & \textbf{Com \texttt{-DCORRECAO}} \\
\hline
\corr{Bug 1 --- byte de velocidade $>127$ em \texttt{sendUnsigned7Bit()}} & \corr{\textsc{failed}} & \corr{\textsc{successful}} \\
\hline
\corr{Bug 2 --- erro de \texttt{write()} não detectado em \texttt{writeCommandWithPayload()}} & \corr{\textsc{failed}} & \corr{\textsc{successful}} \\
\hline
\end{tabular}

\scriptsize
\begin{flushleft}
\corr{Mesmo padrão nas duas ferramentas (ESBMC e CBMC), na máquina da autora, em
28/10/2026. Registros e \textit{hashes} em
\texttt{testes\_hardware/roboclaw/evidencia\_maquina/}.}
\end{flushleft}
\end{table}

\corr{O \textbf{Bug~1} decorre do caminho de terminação de voo: o
\texttt{module.yaml} admite valores de \textit{failsafe} até 257 e o
\textit{mixer} os repassa ao driver sem limitar, enquanto o manual da RoboClaw
(Rev.~6.0, pp.~48--49) restringe o byte de dados dos comandos 0, 1, 4 e 5 à faixa
$0$--$127$. Um valor de \textit{failsafe} igual a 257 produz o byte 130. O
\textit{replay} com UART emulada, sobre a lógica real do driver, envia de fato o
byte 130 de forma determinística. O \textbf{Bug~2} vem de armazenar o retorno de
\texttt{write()} num \texttt{size\_t}: quando \texttt{write()} devolve $-1$, o valor
vira um inteiro sem sinal grande, a checagem \texttt{bytes\_written < 2} não
detecta o erro e a função retorna \texttt{OK}; o driver então espera uma resposta
que não virá, e a falha só aparece como \textit{timeout} em \texttt{select()}.
A correção dos dois está em \texttt{correcao\_roboclaw.patch}.}

\corr{O driver foi reescrito no \texttt{main} do PX4 em 28/10/2026
(PR~\#29473), depois do commit analisado. Pela leitura do código novo, o limite de
\textit{failsafe} passou a 254, o valor é restringido antes da conversão e o
retorno da escrita é guardado em \texttt{ssize\_t}; essa versão não foi
verificada nem ensaiada neste trabalho. O código antigo continua no ramo
\texttt{release/1.18}, em que \texttt{Roboclaw.cpp}, \texttt{Roboclaw.hpp} e
\texttt{module.yaml} são idênticos aos do commit analisado (conferido em
12/11/2026, commit \texttt{8303efda43cb}). Por isso os dois defeitos foram
reportados com esse ramo como alvo (\textit{Issue}~\#29785).}

\corr{\paragraph{Validação em hardware do Bug~1.} O Bug~1 foi observado também
em bancada, numa Pixhawk~6X~RT com firmware compilado do mesmo commit
(\texttt{px4\_fmu-v6xrt\_default}). A porta TELEM2 foi ligada a um ESP32 que
faz o papel da RoboClaw: responde às consultas do driver e registra cada byte
recebido. Com \texttt{RBCLW\_FAIL1} igual a 257 e a terminação de voo acionada
pelo console (\texttt{commander termination}), a placa passou a transmitir o
pacote \texttt{80 00 82 8A 90}: endereço \texttt{0x80}, comando~0, byte de
dados 130 e CRC válido. O pacote se repetiu 24\,136 vezes nas duas capturas
guardadas, sem erro de CRC. Dois ensaios complementares, na mesma montagem,
isolam a causa (Tabela~\ref{tab:roboclaw-hardware}): com o parâmetro em 254, o
maior valor que ainda resulta em byte válido, o mesmo firmware envia 127; e,
com a correção aplicada e o parâmetro de volta em 257, o byte também fica em
127, em 5\,267 pacotes e nenhum acima da faixa.}

\begin{table}[htb]
\centering
\caption[RoboClaw: byte de dados observado na Pixhawk~6X~RT]{\corr{RoboClaw: byte de dados observado na serial da Pixhawk~6X~RT durante a terminação}}
\label{tab:roboclaw-hardware}
\begin{tabular}{llll}
\hline
\corr{\textbf{Firmware}} & \corr{\textbf{\texttt{RBCLW\_FAIL1}}} & \corr{\textbf{Pacote}} & \corr{\textbf{Byte de dados}} \\
\hline
\corr{Original} & \corr{257} & \corr{\texttt{80 00 82 8A 90}} & \corr{130 (fora de 0--127)} \\
\corr{Original} & \corr{254} & \corr{\texttt{80 00 7F B4 22}} & \corr{127} \\
\corr{Corrigido} & \corr{257} & \corr{\texttt{80 00 7F B4 22}} & \corr{127} \\
\hline
\end{tabular}
\end{table}

\corr{Esse resultado tem quatro limites, que delimitam o que ele prova.
Primeiro, a RoboClaw é emulada: o ensaio mostra que a placa \emph{transmite} o
byte inválido, mas não o que o controlador real faz ao recebê-lo, o que o
manual não especifica. Segundo, o firmware não é o de fábrica: o driver não
integra a configuração padrão dessa placa e foi incluído com
\texttt{CONFIG\_DRIVERS\_ROBOCLAW=y}; além disso, foi preciso inicializar dois
membros da classe (\texttt{\_uart\_fd\_set} e \texttt{\_uart\_fd\_timeout}).
Sem isso, em duas partidas registradas, o driver não emitiu mensagem alguma e a
fila \texttt{wq:hp\_default} ficou quase parada (\texttt{battery\_status} a
$0{,}1$~Hz, contra $99{,}9$~Hz depois da alteração; \texttt{roboclaw} a
$0{,}0$~Hz, contra $3{,}4$~Hz). Pelo código, \texttt{Run()} chama a atualização
do \textit{mixer} antes de inicializar a UART, de modo que o primeiro envio
alcança \texttt{select()} com esses membros ainda não inicializados; a causa não
foi confirmada com depurador, e os parâmetros do driver estavam sendo ajustados
entre essas partidas. Registra-se, por isso, como observação de bancada, e não
como terceiro defeito confirmado. Essa alteração não toca o caminho de saída. Terceiro, a placa registrou \textit{timeout} de ACK
durante os ensaios, de modo que o driver foi observado no caminho de erro de
resposta; a medida, que é o byte lido diretamente do fio, não depende disso.
Quarto, só o Bug~1 foi exercitado em hardware; a evidência do Bug~2 continua
sendo a formal e o \textit{replay} com UART emulada. As capturas, os
\textit{diffs} dos firmwares e o código do emulador estão em
\texttt{testes\_hardware/roboclaw/hardware\_pixhawk6xrt/}.}

\subsection{\corr{Frente C --- Driver de bateria batmon}}
\label{subsec:2c-batmon}

\corr{No driver \texttt{smart\_battery/batmon}, o número de células lido do
registrador SMBus da bateria define \texttt{\_cell\_count} até
\texttt{MAX\_CELL\_COUNT}${}=16$, e o laço de publicação escreve
\texttt{new\_report.voltage\_cell\_v[i]} para \texttt{i < \_cell\_count}. Mas
\texttt{voltage\_cell\_v} é um \texttt{float32[14]}
(\texttt{msg/versioned/BatteryStatus.msg}). Uma bateria legítima de 15S ou 16S,
dentro dos registradores que o próprio driver define, já basta para escrever fora
do vetor. Com 16 células, o valor extra cai em \texttt{full\_charge\_capacity\_wh}
(o \textit{replay} publica 3,715, a tensão da 16.ª célula), e \texttt{cell\_count}${}=16$
é publicado num vetor de 14 posições. O driver está nas configurações padrão de
várias placas (fmu-v3 a v6u, cubeorange, entre outras).}

\corr{A Tabela~\ref{tab:batmon-resultados} reúne os resultados. Diferentemente das
outras frentes, a conferência foi feita contra o \texttt{main} de 13/10/2026
(commit \texttt{c9b9e4e25e}), e não \texttt{8be2590c15} --- ressalva que se
mantém explícita.}

\begin{table}[!htb]
\centering
\caption{\corr{batmon: resultados (commit \texttt{c9b9e4e25e})}}
\label{tab:batmon-resultados}
\footnotesize
\begin{tabular}{|p{6.8cm}|p{3.2cm}|p{2.6cm}|}
\hline
\textbf{Ferramenta / verificação} & \textbf{Código atual} & \textbf{\texttt{-DCORRECAO}} \\
\hline
\corr{ESBMC~8.1.0 (\texttt{--unwind 20 --overflow-check})} & \corr{\textsc{failed}} & \corr{\textsc{successful}} \\
\hline
\corr{ESBMC~8.5.0 (\texttt{--multi-property})} & \corr{\textsc{failed} (2 de 13: \texttt{i < 14} e \textit{array bounds})} & \corr{\textsc{successful}} \\
\hline
\corr{CBMC~5.95.1 (\texttt{--unwind 20}, \textit{bounds} e \textit{pointer check})} & \corr{1 de 15 falha} & \corr{0 de 15} \\
\hline
\corr{UBSan (\texttt{./replay 16})} & \corr{\texttt{index 14 out of bounds for type 'float [14]'}} & \corr{---} \\
\hline
\corr{\textit{replay} sem sanitizador} & \corr{\texttt{full\_charge\_capacity\_wh=3.7150}} & \corr{---} \\
\hline
\end{tabular}
\end{table}

\corr{O mesmo defeito foi corrigido no driver UAVCAN/DroneCAN
(\texttt{cbat\_sub\_cb}) pelo PR~\#28606, em 14/09/2026, relatado por um terceiro;
o \texttt{batmon} ficou de fora dessa correção, e as buscas por ``batmon'',
``cell\_count'' e ``voltage\_cell\_v'' nas \textit{issues} do PX4 não encontraram
registro. O código foi reconferido em 12/11/2026 contra o \texttt{main} de
11/11/2026 (commit \texttt{8e94784e1dc0}), sem alteração no driver nem na
mensagem, e o \textit{replay} foi reexecutado com o mesmo resultado. O achado foi
então reportado (\textit{Issue}~\#29787). Não houve ensaio em hardware para este
driver.}

\subsection{\corr{Protocolo de conferência contra o código de produção}}
\label{subsec:2c-protocolo}

\corr{As três frentes seguem o mesmo protocolo de evidência, que operacionaliza a
etapa~(v) (Seção~\ref{subsec:etapa-replay}). As entradas do código verificado são
declaradas não determinísticas e restritas por pré-condições
(\texttt{\_\_VERIFIER\_assume}, aceita pelo ESBMC e pelo CBMC), com a origem de
cada pré-condição anotada no cabeçalho do \textit{harness}: um limite de
parâmetro declarado no \texttt{module.yaml}, o modo como o próprio PX4 chama a
função, ou uma constante do código. No RoboClaw, por exemplo, a saída recebida do
\textit{mixer} foi restrita à faixa $[0;\,257]$ porque o \texttt{module.yaml}
declara \texttt{failsafe: \{min: 0, max: 257\}}; uma pré-condição mais estreita
teria ocultado o Bug~1.}

\corr{Cada candidato alcançável passa por \textit{replay}: a função \emph{real}
do PX4 é compilada com um compilador comum (g++) e executada com os valores do
contraexemplo, sob ASan e UBSan. O critério é binário --- se o código real
reproduz o efeito previsto (o byte 130 no RoboClaw, o acesso a
\texttt{voltage\_cell\_v[14]} no batmon, a escrita em \texttt{\_port[128]} no
MicroStrain), o candidato é defeito reproduzido; se não reproduz, investiga-se a
guarda de produção que o \textit{harness} omitiu, como no \texttt{bst}. Quando há
correção, o mesmo \textit{harness} é reverificado com o \textit{patch}
(\texttt{-DCORRECAO}) e deve passar; para os três candidatos corrigidos
\textit{upstream}, a reverificação usa a correção do próprio PX4
(\texttt{-DCORRECAO\_UPSTREAM}).}

\corr{A validação cruzada com o CBMC foi estendida a todos os \textit{harnesses},
com as opções mapeadas para os equivalentes mais próximos
(Tabela~\ref{tab:equivalencia-opcoes}). Divergências entre as ferramentas são
registradas como resultado da comparação: elas vêm de modelagem de biblioteca
(por exemplo, \texttt{isfinite} tratado pelo CBMC como retorno arbitrário, que
exigiu um prelúdio de modelagem) ou de diferença de versão na contagem de
propriedades, e não de desacordo sobre a existência das violações.}

\begin{table}[!htb]
\centering
\caption{\corr{Equivalência de opções entre ESBMC e CBMC~5.95}}
\label{tab:equivalencia-opcoes}
\footnotesize
\begin{tabular}{|p{6.5cm}|p{7.0cm}|}
\hline
\textbf{ESBMC} & \textbf{CBMC} \\
\hline
\corr{checagens padrão (limites, ponteiros, divisão por zero)} & \corr{\texttt{--bounds-check --pointer-check --div-by-zero-check}} \\
\hline
\corr{\texttt{--overflow-check}} & \corr{\texttt{--signed-overflow-check --float-overflow-check}} \\
\hline
\corr{com \textit{unwinding assertions} (padrão) / \texttt{--no-unwinding-assertions}} & \corr{\texttt{--unwinding-assertions} / padrão} \\
\hline
\corr{\textit{solver} Boolector ou Z3} & \corr{MiniSat (padrão)} \\
\hline
\end{tabular}
\end{table}

\corr{A segunda campanha rendeu ainda um relato fora do PX4. Ao verificar o
\texttt{logger}, o modelo de \texttt{snprintf} do ESBMC devolvia um comprimento
incorreto quando o argumento de \texttt{\%s} era montado em tempo de execução,
o que produzia alarme falso num programa correto. O caso foi reduzido a um
programa mínimo e reportado ao projeto ESBMC (\textit{Issue}~\#8524, de
12/10/2026); os mantenedores confirmaram o problema também no ramo principal e o
corrigiram (PR~\#8563).}

\section{Resumo}
\label{sec:resumo-metodologia}

Este capítulo apresentou a metodologia de verificação formal aplicada ao PX4 Autopilot com o ESBMC. O procedimento foi organizado em oito etapas e aplicado em duas campanhas, que abrangeram 22 módulos distintos. Foram descritos os critérios de seleção dos módulos, a abstração das dependências de \textit{hardware} e do NuttX, a definição das propriedades de segurança e a configuração do verificador. Também foram apresentados os procedimentos de análise dos contraexemplos, conferência com o código de produção, validação cruzada com o CBMC e verificação de correções. A aplicação ao GPS uBlox ilustrou o fluxo da primeira campanha; a segunda ampliou a análise para dez módulos complementares e incorporou a execução do código original sob sanitizadores. Os resultados quantitativos, bem como a classificação e a discussão dos achados, são apresentados no Capítulo~\ref{cap:experimentos}.

\chapter{Experimentos e Resultados}
\label{cap:experimentos}
\label{cap:resultados}

\lc{Sem baseline: executar um analisador estático ou os testes existentes nestes módulos.}

Este capítulo apresenta a avaliação empírica da metodologia descrita no
Capítulo~\ref{cap:metodologia}, aplicada a 12 módulos do PX4 Autopilot
\artigo{(ramo \textit{main}, \textit{commit}~\texttt{8be2590c15})}.
\corr{As Seções~\ref{sec:config-experimental} a~\ref{sec:severidade} tratam
dessa primeira campanha; os resultados da segunda, com dez módulos
complementares e a validação em bancada, estão na
Seção~\ref{sec:resultados-segunda-campanha}.}

\revisao{O resultado principal da campanha é a evidência limitada em
nível de modelo: dez dos doze módulos, cobrindo as principais classes de
entrada não confiável, verificam-se livres de violações de segurança de memória
e de \textit{overflow}} \corr{das classes habilitadas, nos caminhos alcançados
sob $k = 10$\artigo{; a reexecução com as asserções de desenrolamento ligadas
mostrou que nenhum laço desses dez \textit{harnesses} excede o limite
(Seção~\ref{sec:config-experimental})} --- qualificação cujo alcance exato é discutido na
Seção~\ref{sec:ameacas-validade}.} \revisao{O resultado secundário é
a triagem disciplinada de três contraexemplos, que levou a três conclusões
distintas: um falso positivo no contador de FIFO do BMI088, refutado \artigo{na
conferência contra o código de produção}; uma observação de temperatura no BMI088 recusada pelos
mantenedores; e um \textit{underflow} de inteiro sem sinal no driver GPS uBlox}
\artigo{,
proposto upstream como endurecimento defensivo (PR~\#27576) e posteriormente
mostrado inalcançável no driver de produção}\revisao{. Essa distinção entre os três
desfechos é ela própria um resultado: contraexemplos de BMC expõem estados
alcançáveis no modelo que o código de produção pode nunca alcançar, de modo que
a triagem e a revisão pela comunidade são companhias necessárias às provas.}

\section{Configuração Experimental}
\label{sec:config-experimental}

A verificação foi executada em um servidor Intel Core i7-13700K com 64~GB de
RAM sob Ubuntu 24.04, sobre as fontes do PX4
\artigo{(ramo \textit{main}, \textit{commit}~\texttt{8be2590c15}, de 04 de
junho de 2026, 2.935 \textit{commits} após a v1.16.0 e posterior ao lançamento
da v1.17.0)}. Os tempos
reportados são a \textbf{mediana de três execuções}.

O artefato completo --- harnesses, logs de execução e instruções de reprodução
--- está arquivado no Zenodo sob o DOI~\artigo{\texttt{10.5281/zenodo.23282236}
(versão \texttt{v1.2-sbesc}; DOI de conceito \texttt{10.5281/zenodo.21401958})} e
disponível em \url{https://github.com/marcela2881/dissertacao-px4-esbmc}.

\artigo{Todos os doze módulos foram verificados com ESBMC v8.1.0 utilizando
Boolector~3.2.4~\cite{Brummayer2009} como solver; os 36 logs da campanha
(três execuções por módulo) registram \texttt{ESBMC version 8.1.0} e
\texttt{Solving with solver Boolector 3.2.4}. O Boolector é também a
configuração sob a qual o reporte upstream do GPS (Issue~\#26866) foi
registrado; em execuções exploratórias, o Z3~\cite{deMoura2008} não produziu
veredito no harness do GPS dentro do \textit{timeout}
(Seção~\ref{sec:ameacas-validade}).} A invocação padrão foi:

\begin{lstlisting}[
  caption={Comando padrão de execução do ESBMC.},
  label={lst:comando-esbmc},
  numbers=left,
  frame=single,
  language=bash]
esbmc modulo.cpp --unwind 10 --no-unwinding-assertions \
  --overflow-check --default-solver boolector --timeout 120s
\end{lstlisting}

Três módulos exigiram ajuste. O MAVLink utilizou \texttt{--unwind~300}, pois
seus laços de \textit{parsing} são profundos. O Navigator/Mission utilizou
\texttt{--unwind~15}. O VTOL utilizou \texttt{--no-library}, que desabilita a
biblioteca C abstrata do ESBMC: \texttt{sinf} e \texttt{sqrtf} passam a
retornar valores não determinísticos, uma sobreaproximação sob a qual nenhuma
violação das demais verificações é perdida; em contrapartida, as
\artigo{verificações internas da biblioteca deixam de ser feitas}, de modo que o
veredito cobre \artigo{correspondentemente} menos propriedades. A aritmética de ponto
flutuante usa a codificação SMT padrão do ESBMC
(\texttt{--floatbv})~\cite{Gadelha2020}.

Como \texttt{--no-unwinding-assertions} está ativo, laços que excedem $k = 10$
são truncados silenciosamente --- consequência discutida na
Seção~\ref{sec:ameacas-validade} \artigo{e verificada a seguir}.

%% LC: \textbf{[IMPORTANTE]} RECOMENDAÇÃO TRANSVERSAL Nº 3 --- A QUE MAIS FORTALECE O RESULTADO PRINCIPAL, E CUSTA CERCA DE TRÊS MINUTOS DE MÁQUINA.
%% LC: O resultado principal declarado é ``dez dos doze módulos verificam-se livres de violações até o limite de desenrolamento''. Mas a campanha roda com --no-unwinding-assertions. A combinação tem um efeito que a dissertação reconhece com todas as letras nas Ameaças à Validade: laços que excedem k são truncados SILENCIOSAMENTE, ``de modo que um veredito SUCCESSFUL pode ocultar falsos negativos''. Ou seja: o resultado principal é enunciado sob a única configuração que o impede de ser conclusivo, e o próprio texto o admite.
%% LC: Basta re-executar os DEZ módulos sem violação COM asserções de desenrolamento (removendo --no-unwinding-assertions). Só há dois desfechos, e ambos melhoram o trabalho:
%% LC: \hangindent=1em\hangafter=1 - nenhum módulo dispara a asserção -> está PROVADO que nenhum laço excede k, e o veredito passa a ser conclusivo dentro do recorte;
%% LC: \hangindent=1em\hangafter=1 - algum módulo dispara -> descobre-se onde o k é insuficiente e eleva-se o limite naquele módulo. Como a campanha inteira leva ~160 s, o custo é desprezível.
%% LC: Alternativa, se algum laço não fechar: --k-induction nos módulos afetados.

\artigo{\textbf{Reexecução com as asserções de desenrolamento ligadas.} Para
retirar do resultado principal a dependência do truncamento silencioso, os doze
\textit{harnesses} foram reexecutados em 12/10/2026 com as mesmas opções da
Listagem~\ref{lst:comando-esbmc} e os mesmos ajustes por módulo, mas
\textbf{sem} \texttt{--no-unwinding-assertions}, com o ESBMC~8.1.0 e o
Boolector~3.2.4 recompilados e validados contra os logs de julho. A
Tabela~\ref{tab:unwinding} resume os desfechos.}

\begin{table}[H]
\centering
\caption{\artigo{Reexecução com as asserções de desenrolamento ligadas}}
\label{tab:unwinding}
\footnotesize
\begin{tabular}{|p{6.3cm}|p{3.0cm}|p{4.2cm}|}
\hline
\textbf{Módulos} & \textbf{Veredito} & \textbf{Propriedade violada} \\
\hline
\artigo{Barômetro, Battery Monitor, EKF2 Covariance, Expo RC, MAVLink Parser
($k = 300$), Magnetômetro, Navigator/Mission ($k = 15$), PWM Output, RC Update
e VTOL} & \artigo{\textit{sem violação}} & \artigo{---} \\
\hline
\artigo{IMU BMI088} & \artigo{\textit{violação}} & \artigo{a mesma da campanha
(candidato do FIFO)} \\
\hline
\artigo{GPS uBlox} & \artigo{\textit{violação}} & \artigo{asserção de
desenrolamento do laço do \texttt{memcpy}} \\
\hline
\end{tabular}
\end{table}

\artigo{Nos dez módulos sem violação, nenhuma asserção de desenrolamento
falha: nenhum laço desses \textit{harnesses} excede o limite usado, e o veredito
deixa de depender do truncamento silencioso. No IMU, a primeira violação
continua sendo a do candidato do FIFO. No GPS, a primeira falha passa a ser a
asserção de desenrolamento do laço do \texttt{memcpy}, que copia mais de dez
bytes: a violação da Tabela~\ref{tab:resultados} continua válida no modelo,
porque o BMC é \textit{sound}, mas a análise do GPS com $k = 10$ não cobre esse
laço por inteiro. Como controle, o MAVLink com $k = 10$ e as asserções ligadas
falha apenas na asserção de desenrolamento do \texttt{memcpy}, e com $k = 300$
passa. Os logs estão em \texttt{verificacao\_unwinding\_2026-10/} no
repositório de artefatos.}

\lc{Laço do GPS excede k = 10; reexecutar com k cobrindo 210 bytes; reportar resultado.}

%% LC: ``nove dos onze logs'' --- são doze módulos. Onze logs de quê? Se o BMI088 e o GPS geram um log cada e os dez restantes também, seriam doze. Se algum módulo compartilha log, dizer qual. Como está, o leitor tropeça num número que não fecha com nenhum outro do capítulo.

\revisao{Quanto à proveniência das versões:} \artigo{os 36 logs de
\texttt{medianas\_logs/} (três execuções de cada um dos doze módulos)
registram explicitamente \texttt{ESBMC version 8.1.0}, conforme a nota de
rodapé do artigo (``read from all twelve run logs'').} \revisao{Os contraexemplos do BMI088 e do GPS foram
originalmente reportados upstream (Issues~\#26865 e~\#26866) sob o ESBMC
v7.10.0; os tempos aqui reportados foram remedidos sob a configuração da
campanha em v8.1.0.}

%% [BANCA - Prof. Barreto] "Tres valores diferentes de --unwind. Nao deveria justificar
%% melhor? A disparidade e muito grande, de 10, que e o padrao, para 300."
\subsection{\revisao{Escolha do Limite de Desenrolamento}}
\label{subsec:escolha-k}

\revisao{A campanha utilizou três valores distintos de $k$, e a disparidade entre eles
--- sobretudo entre o padrão $k = 10$ e o $k = 300$ do MAVLink --- merece justificativa
explícita, por ser o parâmetro que mais afeta o alcance das conclusões.}

\revisao{O limite de desenrolamento governa um compromisso direto. Aumentá-lo amplia o
conjunto de execuções cobertas e, portanto, a força do veredito; em contrapartida, cada
iteração adicional replica as instruções do laço na fórmula SMT, de modo que o custo cresce
rapidamente e pode tornar a verificação intratável dentro do \textit{timeout}. Não existe,
assim, um valor universalmente correto: o valor adequado é o menor que cobre o comportamento
de interesse do módulo.}

\revisao{O critério adotado foi derivar $k$ da estrutura do laço dominante de cada
módulo, e não arbitrá-lo:}

\begin{itemize}
\item \revisao{\textbf{$k = 10$ (padrão, nove módulos).} Nesses módulos, os laços
percorrem estruturas de dimensão pequena e fixa --- canais de rádio, saídas de PWM, elementos
de matriz de covariância --- todos com menos de dez iterações no pior caso. O valor cobre o
laço inteiro, e aumentá-lo não acrescentaria execuções novas.} \artigo{A reexecução com as
asserções de desenrolamento ligadas confirma isso: nenhuma delas falha nesses
\textit{harnesses} (Tabela~\ref{tab:unwinding}).}

\item \revisao{\textbf{$k = 15$ (Navigator/Mission).} O laço itera sobre
\textit{waypoints} da missão carregada\artigo{, e com $k = 15$ nenhuma asserção de
desenrolamento falha no \textit{harness} (Tabela~\ref{tab:unwinding}).}}

\item \revisao{\textbf{$k = 300$ (MAVLink Parser).} Aqui o valor não é arbitrário nem
uma tentativa de aumentar cobertura por precaução: o \texttt{memcpy} \artigo{do
\textit{harness}, que copia para o buffer de transmissão,} itera
proporcionalmente ao tamanho do pacote, e o buffer \artigo{é de 280 bytes, o tamanho
máximo de pacote}. Qualquer $k$
inferior a 280 deixaria de cobrir pacotes de tamanho máximo\artigo{. A reexecução com
as asserções ligadas confirma isso: com $k = 10$, a única falha é a asserção de
desenrolamento do \texttt{memcpy} --- um resultado inconclusivo por profundidade
insuficiente, e não um defeito ---, e com $k = 300$ o \textit{harness} passa
(Seção~\ref{sec:config-experimental}).}
O valor 300 é o menor múltiplo redondo acima do
tamanho do buffer.}
\end{itemize}

\revisao{Segue-se disso que a disparidade entre 10 e 300 não reflete inconsistência de
critério, mas heterogeneidade dos módulos: o MAVLink é o único cujo laço dominante percorre um
buffer de centenas de bytes. Duas consequências devem ser registradas. Primeiro, os totais de
VCC por módulo misturam limites diferentes e, portanto, não são diretamente comparáveis entre
si. Segundo, o procedimento é reativo por construção --- parte-se do padrão e eleva-se $k$
quando o verificador sinaliza insuficiência --- de modo que um módulo cujo laço exceda o limite
sem disparar asserção de desenrolamento poderia passar despercebido, limitação agravada pelo
uso de \texttt{--no-unwinding-assertions} e discutida na
Seção~\ref{sec:ameacas-validade}.} \artigo{A reexecução com as asserções ligadas
(Tabela~\ref{tab:unwinding}) afasta esse risco nos dez módulos sem violação.}

\section{Questões Experimentais}
\label{sec:objetivos-verificacao}

A avaliação responde a três questões experimentais (QE). \corr{Para cada uma
indicam-se a métrica observada e o critério adotado para considerá-la
respondida.}

\begin{itemize}
    \item \textbf{QE1:} Qual é a latência de verificação por módulo em hardware
    de prateleira? \corr{\textit{Métrica:} tempo do procedimento de decisão,
    mediana de três execuções, com mínimo e máximo.
    \textit{Critério:} todo módulo conclui dentro do \textit{timeout} de
    $120$\,s.}
    \item \textbf{QE2:} O BMC expõe violações alcançáveis no modelo em código
    que os testes unitários existentes do PX4 não referenciam?
    \corr{\textit{Métrica:} número de contraexemplos, desfecho de cada um após
    a triagem e existência de teste unitário que referencie a função afetada.
    \textit{Critério:} ao menos um contraexemplo que sobreviva à triagem como
    violação alcançável no modelo, em função sem teste unitário direto.}
    \item \textbf{QE3:} Quais configurações por módulo são necessárias para
    tornar cada módulo tratável? \corr{\textit{Métrica:} número de módulos que
    exigem opções diferentes da invocação padrão, e quais opções.
    \textit{Critério:} questão descritiva, sem limiar de aceitação.}
\end{itemize}

\corr{Os critérios foram fixados depois da campanha, ao redigir este capítulo;
servem para tornar explícito o que cada resposta afirma, e não como hipóteses
registradas previamente.}

\section{Módulos Verificados}
\label{sec:modulos-verificados}

Os módulos foram priorizados pelo impacto na segurança ao longo do caminho de
sensoriamento até atuação e pela viabilidade de isolamento por \textit{stubs}.
Foram verificados 12 módulos por meio de \textit{harnesses} construídos com assistência de IA generativa e revisados pela autora,
totalizando aproximadamente $2.800$ linhas, medidas com \texttt{wc -l} sobre os
arquivos \texttt{test\_*.cpp} do repositório de artefatos.

\revisao{Em vez de reportar percentual de cobertura de linhas --- que
seria enganoso, já que cada harness abstrai a maior parte do seu módulo e
explora apenas caminhos alcançáveis sob $k = 10$ --- a cobertura é
caracterizada por \textbf{classe de entrada}, conforme a última coluna da
Tabela~\ref{tab:modulos}.}

\begin{table}[H]
\centering
\caption[Módulos verificados na primeira campanha]{Módulos verificados: tamanho do alvo, tamanho do harness, localização
no código-fonte do PX4 e classe de entrada exercitada}
\label{tab:modulos}
\footnotesize
\begin{tabular}{|l|l|r|r|l|}
\hline
\textbf{Módulo} & \textbf{Caminho (\texttt{src/})} & \textbf{Alvo} &
\textbf{Harness} & \textbf{Classe de entrada} \\
\hline
GPS uBlox         & \texttt{drivers/gps}          & \artigo{1712} & 294 & Sensor serial \\
IMU BMI088        & \texttt{drivers/imu/.../bmi088} & \artigo{700} & 319 & Sensor serial \\
MAVLink           & \texttt{modules/mavlink}      & \artigo{3624} &  82 & Telemetria \\
Expo RC           & \texttt{lib/mathlib/math}     &  \artigo{303} & 129 & Rádio \\
Battery Monitor   & \texttt{lib/battery}          &  \artigo{449} & 242 & Sensor de energia \\
PWM Output        & \texttt{drivers/pwm\_out}     &  337 & 208 & Atuação \\
RC Update         & \texttt{lib/rc}               & \artigo{769} & 253 & Rádio \\
Navigator         & \texttt{modules/navigator}    & \artigo{497} & 291 & Lógica de missão \\
EKF2 Covariance   & \texttt{modules/ekf2/EKF}     &  385 & 296 & Estimação \\
Magnetômetro      & \texttt{modules/sensors}      &  792 & 239 & \corr{Agregação de sensores} \\
Barômetro         & \texttt{modules/sensors}      &  \artigo{246} & 278 & \corr{Agregação de sensores} \\
VTOL              & \texttt{modules/vtol\_att\_control} & \artigo{591} & 146 & Atuação \\
\hline
\end{tabular}

\footnotesize
Alvo = LoC \artigo{do(s) arquivo(s)-fonte do PX4 que o harness extrai (\textit{commit} fixado)}; Harness = LoC do harness (ambos por \texttt{wc -l}).
\end{table}

\lc{Acrescentar coluna: linhas do PX4 efetivamente extraídas por harness, separadas do código de modelo.}
\corr{Os arquivos-alvo somam $10.405$ linhas, cerca de $1{,}5\,\%$ das $710.259$
linhas de C/C++ em \texttt{src/} no \textit{commit} analisado ($2.710$ arquivos,
contados com \texttt{wc -l}). Os alvos do Magnetômetro e do Barômetro estão em
\texttt{modules/sensors}: são módulos que agregam e selecionam leituras já
publicadas pelos drivers, e não drivers de barramento serial, razão pela qual
aparecem como ``agregação de sensores''.}

\section{Resultados Gerais}
\label{sec:resultados-gerais}

A Tabela~\ref{tab:resultados} sumariza os desfechos. Dez dos doze módulos
verificam-se sem violação até o limite\artigo{; os dois restantes (IMU BMI088 e
GPS uBlox) param na primeira violação encontrada. Como no artigo derivado, as
contagens de condições de verificação (VCCs) por \textit{harness} não são
tabeladas: elas misturam limites de desenrolamento diferentes (10, 15 e 300) e,
nos dois módulos com violação, cobrem só o caminho até a primeira violação. Os
valores constam dos logs do repositório de artefatos.}

\begin{table}[H]
\centering
\caption{Resultados da verificação para os 12 módulos do PX4}
\label{tab:resultados}
\footnotesize
\setlength{\tabcolsep}{4pt}
\begin{tabular}{|l|l|c|r|l|r|r|}
\hline
\textbf{Módulo} & \textbf{Desfecho} &
\textbf{VAM}$^\P$ & \textbf{Unw.} & \textbf{Solver} & \textbf{Tempo}$^\|$ & \corr{\textbf{Decisão}$^{\star}$ (mín--máx)} \\
\hline
IMU BMI088        & \corr{\textit{contraex.\ (triado)}} & 0$^\ddagger$ &  10 & \artigo{Bool.} & \artigo{0,1\,s}$^{\#}$ & \corr{0,122 (0,121--0,122)} \\
GPS uBlox         & \textit{violação}     & 1 &  10 & \artigo{Bool.} & 64,9\,s$^{\#}$ & \corr{64,857 (64,765--64,874)} \\
MAVLink Parser    & \textit{sem violação} & 0 & 300 & \artigo{Bool.} & 49,2\,s & \corr{46,970 (46,916--47,083)} \\
Expo RC           & \textit{sem violação} & 0 &  10 & \artigo{Bool.} &  4,9\,s & \corr{4,852 (4,840--4,866)} \\
Battery Monitor   & \textit{sem violação} & 0 &  10 & \artigo{Bool.} &  6,7\,s & \corr{6,727 (6,703--6,731)} \\
PWM Output        & \textit{sem violação} & 0 &  10 & \artigo{Bool.} &  0,2\,s & \corr{0,156 (0,156--0,156)} \\
RC Update         & \textit{sem violação} & 0 &  10 & \artigo{Bool.} &  9,5\,s & \corr{9,457 (9,384--9,534)} \\
Navigator/Mission & \textit{sem violação} & 0 &  15 & \artigo{Bool.} &  0,1\,s & \corr{0,095 (0,095--0,095)} \\
EKF2 Covariance   & \textit{sem violação} & 0 &  10 & \artigo{Bool.} & 12,2\,s & \corr{12,200 (12,181--12,228)} \\
Magnetômetro      & \textit{sem violação} & 0 &  10 & \artigo{Bool.} &  6,6\,s & \corr{6,621 (6,616--6,629)} \\
Barômetro         & \textit{sem violação} & 0 &  10 & \artigo{Bool.} &  0,4\,s & \corr{0,360 (0,351--0,363)} \\
VTOL$^\S$         & \textit{sem violação} & 0 &  10 & \artigo{Bool.} &  4,0\,s & \corr{3,942 (3,911--3,943)} \\
\hline
\textbf{Total} & \corr{\textbf{10 sem violação, 1 violação,}} & \textbf{1} &
--- & --- & \artigo{\textbf{$\approx$94\,s}$^{\dagger}$} & \corr{\textbf{156,359}} \\
               & \corr{\textbf{1 contraex.\ refutado/recusado}} & & & & & \\
\hline
\end{tabular}

\footnotesize
\begin{flushleft}
$^\P$VAM = violações alcançáveis no modelo (\textbf{não} defeitos
confirmados).\\
\artigo{Bool.\ = Boolector~3.2.4, com ESBMC v8.1.0, em todos os módulos.}\\
$^\ddagger$\corr{Dois contraexemplos, ambos triados como não-defeito:
FIFO, falso positivo; temperatura, recusada upstream.} \artigo{O tempo da
tabela é o do candidato FIFO (\texttt{test\_choice = 2}), em que as três
execuções param (0,122; 0,122; 0,121\,s); o contraexemplo de temperatura vem de
uma execução separada (0,272\,s), Seção~\ref{sec:caso-imu}.}\\
$^{\dagger}$\artigo{93,8\,s nos módulos que concluem (dez sem violação e o
IMU); 158,7\,s incluindo o GPS (64,9\,s).}\\
$^\S$VTOL com \texttt{--no-library}: \artigo{verificações internas da
biblioteca não são feitas.}\\
$^\|$Mediana de três execuções.\\
\corr{$^{\star}$Tempo do procedimento de decisão (\textit{Runtime decision procedure}), em segundos: mediana, mínimo e máximo das três execuções. É a única medida registrada nos logs dos doze módulos.}\\
$^{\#}$\artigo{Módulos com violação param na primeira violação; tempo do
procedimento de decisão (\textit{Runtime decision procedure}); nos demais,
\textit{BMC program time}.}
\end{flushleft}
\end{table}

O tempo máximo foi de $64,9$\,s (GPS) e o mínimo de $0,1$\,s
\artigo{(Navigator/Mission e IMU BMI088)}.

\corr{A coluna \textit{Tempo} reproduz a Tabela~II do artigo e mistura duas
medidas, conforme a nota~$^{\#}$: os logs dos dois módulos com violação não
registram o tempo total do programa. Para permitir comparação uniforme, a
última coluna traz o tempo do procedimento de decisão dos doze módulos, com
mínimo e máximo; a dispersão entre execuções é inferior a $0{,}2$\,s em todos
os casos. O tempo de relógio de ponta a ponta não foi registrado nos logs da
campanha. A Figura~\ref{fig:visao-resultados} resume desfechos e tempos.}

\begin{figure}[!htb]
\centering
\begin{tikzpicture}
  \draw[palEdge] (0,0.3) -- (0,-4.97);
  \draw[palEdge, densely dotted] (0.000,0.3) -- (0.000,-4.97);
  \node[font=\scriptsize, anchor=north] at (0.000,-4.97) {0,01};
  \draw[palEdge, densely dotted] (2.250,0.3) -- (2.250,-4.97);
  \node[font=\scriptsize, anchor=north] at (2.250,-4.97) {0,1};
  \draw[palEdge, densely dotted] (4.500,0.3) -- (4.500,-4.97);
  \node[font=\scriptsize, anchor=north] at (4.500,-4.97) {1};
  \draw[palEdge, densely dotted] (6.750,0.3) -- (6.750,-4.97);
  \node[font=\scriptsize, anchor=north] at (6.750,-4.97) {10};
  \draw[palEdge, densely dotted] (9.000,0.3) -- (9.000,-4.97);
  \node[font=\scriptsize, anchor=north] at (9.000,-4.97) {100};
  \node[anchor=east, font=\scriptsize] at (-0.1,-0.00) {GPS uBlox};
  \draw[fill=palStrong, draw=palEdgeDark, thick] (0,-0.14) rectangle (8.577,0.14);
  \node[anchor=west, font=\scriptsize] at (8.627,-0.00) {64,857\,s};
  \node[anchor=east, font=\scriptsize] at (-0.1,-0.42) {MAVLink Parser};
  \draw[fill=palFill, draw=palEdge] (0,-0.56) rectangle (8.262,-0.28);
  \node[anchor=west, font=\scriptsize] at (8.312,-0.42) {46,970\,s};
  \node[anchor=east, font=\scriptsize] at (-0.1,-0.84) {EKF2 Covariance};
  \draw[fill=palFill, draw=palEdge] (0,-0.98) rectangle (6.944,-0.70);
  \node[anchor=west, font=\scriptsize] at (6.994,-0.84) {12,200\,s};
  \node[anchor=east, font=\scriptsize] at (-0.1,-1.26) {RC Update};
  \draw[fill=palFill, draw=palEdge] (0,-1.40) rectangle (6.695,-1.12);
  \node[anchor=west, font=\scriptsize] at (6.745,-1.26) {9,457\,s};
  \node[anchor=east, font=\scriptsize] at (-0.1,-1.68) {Battery Monitor};
  \draw[fill=palFill, draw=palEdge] (0,-1.82) rectangle (6.363,-1.54);
  \node[anchor=west, font=\scriptsize] at (6.413,-1.68) {6,727\,s};
  \node[anchor=east, font=\scriptsize] at (-0.1,-2.10) {Magnetômetro};
  \draw[fill=palFill, draw=palEdge] (0,-2.24) rectangle (6.347,-1.96);
  \node[anchor=west, font=\scriptsize] at (6.397,-2.10) {6,621\,s};
  \node[anchor=east, font=\scriptsize] at (-0.1,-2.52) {Expo RC};
  \draw[fill=palFill, draw=palEdge] (0,-2.66) rectangle (6.043,-2.38);
  \node[anchor=west, font=\scriptsize] at (6.093,-2.52) {4,852\,s};
  \node[anchor=east, font=\scriptsize] at (-0.1,-2.94) {VTOL};
  \draw[fill=palFill, draw=palEdge] (0,-3.08) rectangle (5.840,-2.80);
  \node[anchor=west, font=\scriptsize] at (5.890,-2.94) {3,942\,s};
  \node[anchor=east, font=\scriptsize] at (-0.1,-3.36) {Barômetro};
  \draw[fill=palFill, draw=palEdge] (0,-3.50) rectangle (3.502,-3.22);
  \node[anchor=west, font=\scriptsize] at (3.552,-3.36) {0,360\,s};
  \node[anchor=east, font=\scriptsize] at (-0.1,-3.78) {PWM Output};
  \draw[fill=palFill, draw=palEdge] (0,-3.92) rectangle (2.685,-3.64);
  \node[anchor=west, font=\scriptsize] at (2.735,-3.78) {0,156\,s};
  \node[anchor=east, font=\scriptsize] at (-0.1,-4.20) {IMU BMI088};
  \draw[fill=palPaper, draw=palEdgeDark, dashed] (0,-4.34) rectangle (2.444,-4.06);
  \node[anchor=west, font=\scriptsize] at (2.494,-4.20) {0,122\,s};
  \node[anchor=east, font=\scriptsize] at (-0.1,-4.62) {Navigator/Mission};
  \draw[fill=palFill, draw=palEdge] (0,-4.76) rectangle (2.200,-4.48);
  \node[anchor=west, font=\scriptsize] at (2.250,-4.62) {0,095\,s};
  \draw[fill=palFill, draw=palEdge] (0,-6.37) rectangle (0.35,-6.17); \node[anchor=west, font=\scriptsize] at (0.4,-6.27) {sem violação};
  \draw[fill=palStrong, draw=palEdgeDark, thick] (2.6,-6.37) rectangle (2.95,-6.17); \node[anchor=west, font=\scriptsize] at (3.0,-6.27) {violação alcançável no modelo};
  \draw[fill=palPaper, draw=palEdgeDark, dashed] (7.2,-6.37) rectangle (7.55,-6.17); \node[anchor=west, font=\scriptsize] at (7.6,-6.27) {contraexemplo triado};
  \node[font=\scriptsize, anchor=north] at (4.50,-5.27) {tempo do procedimento de decisão (s, escala logarítmica)};
\end{tikzpicture}
\caption[Desfecho e tempo de decisão por módulo (primeira campanha)]{\corr{Desfecho e tempo do procedimento de decisão por módulo na primeira campanha (mediana de três execuções). Elaborado pela autora.}}
\label{fig:visao-resultados}
\end{figure}

\lc{Um main multiplexa vários testes; a primeira falha mascara as demais; executar cada teste separadamente.}

\begin{tcolorbox}[colback=blue!5!white, colframe=blue!60!black,
title=\textbf{Resposta à QE1 --- Latência de verificação}]
Todo módulo cronometrado verifica-se em menos de noventa segundos em hardware
de prateleira. Como cada um é verificado isoladamente por meio de um harness
pequeno, isso demonstra \textbf{viabilidade por módulo}, e não escalabilidade
de sistema completo.
\end{tcolorbox}

\revisao{Os módulos verificados não constituem uma amostra uniforme:
situam-se no caminho crítico de sensoriamento até atuação, onde bytes não
confiáveis de I2C, SPI e UART entram no \textit{flight stack} pela primeira
vez, cobrindo as principais classes de entrada não confiável --- barramentos
seriais de sensores, enlace de rádio e telemetria, atuação, estimação e lógica
de missão.} \corr{O que ficou de fora, porém, não é código periférico. No
\textit{commit} analisado, as $710.259$ linhas de C/C++ de \texttt{src/}
distribuem-se entre \texttt{modules} ($45{,}8\,\%$), \texttt{drivers}
($39{,}4\,\%$), \texttt{lib} ($11{,}4\,\%$) e \texttt{systemcmds}
($2{,}9\,\%$), e os maiores blocos não verificados incluem código crítico para
o voo: \texttt{drivers/uavcan} ($58.267$ linhas), \texttt{drivers/imu}
($46.127$), \texttt{modules/ekf2} ($33.516$, do qual só a rotina de covariância
foi verificada), \texttt{modules/commander} ($25.377$) e
\texttt{modules/navigator} ($17.613$). A campanha cobre, portanto, uma fração
pequena e escolhida do \textit{firmware}.} \revisao{As exceções críticas para voo (\textit{commander},
\textit{failsafe} e lógica de \textit{geofence}) são trabalho futuro
prioritário, e os barramentos CAN/UAVCAN, a ponte RTPS/uXRCE-DDS e os drivers
de fluxo óptico e telêmetro permanecem não verificados.}

\section{Validação Cruzada com o CBMC}
\label{sec:cbmc}

\revisao{Para avaliar a dependência dos resultados em relação à
ferramenta, o CBMC~5.95.1~\cite{Kroening2014} foi executado sobre os}
\corr{harnesses adaptados aos intrínsecos do CBMC
(\texttt{\_\_CPROVER\_assume} e \texttt{\_\_CPROVER\_assert} no lugar de
\texttt{\_\_ESBMC\_assume} e \texttt{\_\_ESBMC\_assert}), preservando
pré-condições e asserções,} \revisao{com opções equivalentes
(\texttt{--unwind 10}, verificações de \textit{overflow} e de limites, sem
asserções de desenrolamento) e o mesmo \textit{timeout} de $120$\,s.}
\corr{A adaptação é parte do experimento e está registrada no repositório de
artefatos}\artigo{, que contém o \textit{harness} do CBMC para o GPS
(\texttt{cbmc/}) e os logs do CBMC para o IMU e o GPS.}

\revisao{Para o IMU BMI088, as duas ferramentas concordaram. Para o GPS
uBlox, sob a pré-condição \texttt{len}~$\geq$~$190$
(\texttt{GPS\_DUMP\_DATA\_SIZE}~$-$~$10$, semeando o campo de comprimento
próximo à fronteira do buffer), ambas sinalizam a violação}
\artigo{, mas o CBMC apenas com limites reduzidos (\texttt{--unwind} 2, 3 e 5:
\textit{memcpy destination region writeable}); com \texttt{--unwind 10} ele não
terminou dentro de um limite de 30 minutos. O harness original do CBMC, com a
pré-condição \texttt{len}~$<$~$200$, retornou \textsc{verification successful},
porque a pré-condição exclui os estados violadores. O harness do ESBMC usa a
janela $\texttt{GPS\_DUMP\_DATA\_SIZE} - 5$. As duas ferramentas reportam a
mesma propriedade --- acesso fora dos limites dentro do \texttt{memcpy}
(\textit{dereference failure} no ESBMC); a cadeia CWE-191 $\rightarrow$
CWE-119/787 é uma classificação deste trabalho, e não saída das ferramentas.}

\corr{No IMU, o CBMC~5.95.1 reporta 2 de 6 propriedades violadas: as duas
asserções do \textit{harness} (faixa de temperatura e
\texttt{count <= FIFO\_SIZE}) falham, e as quatro verificações de
\textit{overflow} aritmético passam; o tempo do procedimento de decisão é
inferior a $0{,}01$\,s. São as mesmas duas propriedades violadas no ESBMC. O log
(\texttt{cbmc/resultado\_cbmc\_imu.txt}) não registra a linha de comando usada.
No GPS, o comando registrado é o da Listagem~\ref{lst:comando-cbmc}:}

\begin{lstlisting}[
  caption={\corr{Comando do CBMC para o \textit{harness} do GPS ($u \in \{2, 3, 5, 10\}$).}},
  label={lst:comando-cbmc},
  numbers=left,
  frame=single,
  language=bash]
cbmc test_gps_cbmc_fair.cpp --unwind $u --no-unwinding-assertions \
  --bounds-check --pointer-check \
  --signed-overflow-check --unsigned-overflow-check
\end{lstlisting}

\corr{As janelas de entrada dos dois \textit{harnesses} do GPS não são
idênticas: o do CBMC assume \texttt{len}~$\geq 190$ e o do ESBMC,
\texttt{len}~$\geq 195$. Os dois foram escritos em momentos diferentes e a
diferença não foi uniformizada. Ela não altera a comparação, porque os estados
que violam a propriedade (\texttt{len}~$> 200$) pertencem às duas janelas; uma
reexecução com a mesma janela nas duas ferramentas tornaria a comparação mais
limpa e fica registrada como pendência. As versões de compilador registradas
nos artefatos são as do \textit{replay} da segunda campanha (g++~13.3.0); o
ESBMC e o CBMC usam seus próprios \textit{front-ends}.}

\lc{Validação cruzada cobre dois módulos; executar CBMC também nos dez harnesses sem violação.}

\revisao{Cabe uma ressalva importante: isso \textbf{não} constitui
confirmação independente de alcançabilidade. Com a mesma pré-condição no
harness, a concordância corrobora apenas a codificação do problema, não que o
estado surja no driver de produção.} \corr{A ressalva é ainda mais forte quando
se considera que as duas execuções partem de codificações distintas do mesmo
problema: a concordância entre elas atesta que o defeito sobrevive à tradução
entre os dois conjuntos de intrínsecos, e nada além disso.}

\section{Estudo de Caso: Driver IMU BMI088}
\label{sec:caso-imu}

\corr{O driver do BMI088 (\artigo{$700$}~LoC no arquivo verificado, conforme a
Tabela~\ref{tab:modulos})} produziu dois candidatos.
\revisao{A triagem levou a conclusões diferentes para cada um --- e
nenhuma delas é a de defeito confirmado.}

\subsection{Temperatura fora da faixa operacional (contestado, recusado upstream)}
\label{subsec:imu-temperatura}

O driver não valida a temperatura antes da aritmética de compensação térmica; o
datasheet do BMI088~\cite{Bosch2021} restringe a operação à faixa de
$[-40,\,+85]\,^\circ$C (Equação~\ref{eq:prop_temp}). O ESBMC produziu
\texttt{temp\_msb = 63} e \texttt{temp\_lsb = 95}
(\texttt{Temp\_uint11 = 506}), resultando em $86{,}25\,^\circ$C.
\corr{A conversão especificada pelo fabricante é
$T = \texttt{Temp\_int11} \times 0{,}125 + 23$, de modo que
$506 \times 0{,}125 + 23 = 86{,}25\,^\circ$C; o traço literal consta do
Apêndice~\ref{apend:contraexemplos}.}

\corr{Se a observação fosse aceita como defeito, a classe em que se enquadraria
é a CWE-1284 (\textit{Improper Validation of Specified Quantity in Input}),
subclasse de CWE-20 (\textit{Improper Input Validation}): o driver consome uma
leitura fora da faixa do \textit{datasheet} sem verificação. Como a observação
foi recusada, o rótulo indica apenas a classe do cenário investigado, e não uma
vulnerabilidade confirmada.}

\revisao{Se isso constitui defeito é ponto contestado. Os mantenedores
observam que a temperatura não alimenta o caminho inercial, e que descartar
quadros da IMU diante de uma excursão transitória arrisca perda de controle em
plataformas com uma única IMU; tais exceções seriam melhor tratadas pelas
camadas de estimação a jusante, como a fusão de sensores do EKF2, do que por
paradas em nível de driver. Trata-se de um compromisso entre rigor no driver e
tolerância a falhas em nível de sistema.}

\revisao{Um \textit{patch} upstream que limitava a temperatura fora de
faixa (PX4~PR~\#27575) foi \textbf{rejeitado}: os mantenedores objetaram à
abordagem de limitar e avisar, e fecharam tanto o PR quanto a
Issue~\#26865\footnote{\url{https://github.com/PX4/PX4-Autopilot/issues/26865}}.
A observação é mantida nesta dissertação apenas como nota de projeto, \textbf{sem
consequência de segurança confirmada e sem remediação aceita}, ilustrando a
etapa de triagem da metodologia em vez de um defeito.} \corr{A
Figura~\ref{fig:ce-imu-temp} reproduz a saída do verificador.}

\begin{figure}[H]
    \centering
    \definecolor{ceFundo}{HTML}{F7F8FA}\definecolor{ceCinza}{HTML}{5A6070}%
    \definecolor{ceAmarelo}{HTML}{0B5CAD}\definecolor{ceVermelho}{HTML}{B3261E}%
    \begin{tcolorbox}[colback=ceFundo, colframe=black!45, boxrule=0.6pt, arc=2pt,
      left=8pt, right=8pt, top=6pt, bottom=6pt, width=0.95\linewidth,
      fontupper=\ttfamily\scriptsize\linespread{1.45}\selectfont]
    \textcolor{ceVermelho}{\textbf{[Counterexample]}}\\
    \textcolor{ceCinza}{State 2 -{}- test\_bmi088\_imu.cpp:120 -{}- test\_temperature\_calculation}\\
    \hspace*{1em}\fcolorbox{ceAmarelo}{blue!6}{\textcolor{ceAmarelo}{\textbf{temp\_msb = 63 (00111111)}}}\\
    \textcolor{ceCinza}{State 3 -{}- test\_bmi088\_imu.cpp:121 -{}- test\_temperature\_calculation}\\
    \hspace*{1em}\fcolorbox{ceAmarelo}{blue!6}{\textcolor{ceAmarelo}{\textbf{temp\_lsb = 95 (01011111)}}}\\
    \textcolor{ceCinza}{State 7 -{}- test\_bmi088\_imu.cpp:42 -{}- updateTemperature}\\
    \hspace*{1em}\fcolorbox{ceAmarelo}{blue!6}{\textcolor{ceAmarelo}{\textbf{Temp\_uint11 = 506 (00000001 11111010)}}}\\
    \textcolor{ceCinza}{State 11 -{}- test\_bmi088\_imu.cpp:51 -{}- updateTemperature}\\
    \hspace*{1em}\fcolorbox{ceAmarelo}{blue!6}{\textcolor{ceAmarelo}{\textbf{temperature = 8.625000e+1f (...00000000000)}}}\\
    \textcolor{ceVermelho}{\textbf{Violated property:}}\\
    \textcolor{ceVermelho}{\hspace*{1em}file test\_bmi088\_imu.cpp line 126 function test\_temperature\_calculation}\\
    \textcolor{ceVermelho}{\hspace*{1em}assertion temperature >= -40.0f \&\& temperature <= 85.0f}\\
    \textcolor{ceVermelho}{\hspace*{1em}temperature >= -40.000000f \&\& temperature <= 85.000000f}\\
    \textcolor{ceVermelho}{\textbf{VERIFICATION FAILED}}\\[-2pt]
    \textcolor{ceCinza}{\rule{\linewidth}{0.3pt}}\\
    \textcolor{ceAmarelo}{8.625000e+1f = 86,25\,$^{\circ}$C > 85\,$^{\circ}$C (limite operacional do datasheet BMI088)}\\
    \textcolor{ceCinza}{ESBMC 8.1.0 + Boolector 3.2.4 -{}- log: resultado\_temperatura\_8.1.log}\\
    \textcolor{ceCinza}{PX4 main, commit 8be2590c15 -{}- Issue \#26865 (fechada) / PR \#27575 (fechado)}
    \end{tcolorbox}
    \caption[Contraexemplo de temperatura do BMI088 (saída do ESBMC)]{\corr{Saída do ESBMC para o contraexemplo de temperatura do BMI088, montada a partir do log (ESBMC~v8.1.0, Boolector~3.2.4).}}
    \label{fig:ce-imu-temp}
\end{figure}

%% LC: \textbf{[EM ABERTO]} Conferir dois dados internos desta figura contra o texto: ela cita o datasheet como ``BST-BMI088-DS001 §5.3.7'', ao passo que o .bib registra ``BST-BMI088-DS000-19, rev. 1.9'' e a Tab. \ref{tab:proveniencia} repete essa segunda forma. São dois números de documento diferentes para a mesma fonte. Verificar qual é o correto no PDF da Bosch e uniformizar figura, .bib e tabela.
%% LC: \textbf{[AJUSTE]} Renomear figuras/imagematual.jpeg (usada na Fig. \ref{fig:ce-gps}). Sugestão: ce\_gps\_underflow.jpeg. Aproveitar para remover de figuras/ os arquivos que o LEIA-ME já identificou como contraditórios com o texto atual: bug1\_temperatura.jpeg (mostra 512 -> 87,0 $^\circ$C), bug\_gps.png (input\_len=1, len=216) e bug2\_fifo.jpeg (count=1032). Enquanto estiverem no .zip entregue, são um risco: basta alguém abrir a pasta.

\subsection{Contador do FIFO (falso positivo)}
\label{subsec:imu-fifo}

\revisao{O ESBMC produziu} \artigo{\texttt{count = 9088}} \revisao{bytes, mais de oito
vezes o FIFO de $1.024$ bytes, o que sugeria CWE-190 e CWE-787. A}
\artigo{conferência} \revisao{desse contraexemplo} \artigo{contra o código} \revisao{do driver de produção \textbf{não
instrumentado} revelou um \textbf{falso positivo}.}

\revisao{O driver liberado protege a leitura em três níveis: o contador
de bytes é comparado ao tamanho do FIFO e reiniciado sempre que
\texttt{fifo\_byte\_counter}~$\geq$~\texttt{FIFO::SIZE}; a contagem de amostras
é limitada por \texttt{FIFO\_MAX\_SAMPLES}; e o comprimento da transferência é
restringido por
\texttt{math::min(samples*sizeof(FIFO::DATA)+4,\,FIFO::SIZE)}. O valor \artigo{$9088$}
$\geq 1.024$ dispara a reinicialização da primeira guarda, de modo que nenhum
acesso fora dos limites é alcançável.}

\revisao{O candidato surgiu porque o harness omitia essas guardas, e foi
capturado pela etapa obrigatória de \artigo{conferência contra o código de produção}. Este caso é o exemplo mais
didático da tese central deste trabalho:} \artigo{o candidato chegou a ser
proposto \textit{upstream}, como uma limitação incluída no PR~\#27575, antes da
conferência; foi a leitura do driver de produção que mostrou tratar-se de um
contraexemplo formalmente válido para um defeito inexistente.}

\artigo{A refutação foi feita por leitura do código-fonte do driver, e não por
\textit{replay}: para repeti-la, basta conferir as três guardas acima em
\texttt{BMI088\_Accelerometer.cpp}. O artefato depositado não contém programas
de \textit{replay} nem registros de sanitizadores.}

\lc{Refutação depende de leitura; verificar formalmente as três guardas, como feito para o GPS.}

\section{Estudo de Caso: Driver GPS uBlox}
\label{sec:caso-gps}

O \artigo{harness do} driver uBlox (\artigo{arquivo-fonte de $1.712$}~LoC) foi verificado em $64,9$\,s. A verificação de
espaço disponível \artigo{subtrai o nível de preenchimento da capacidade}
(Listagem~\ref{lst:gps-vulneravel}):

\begin{lstlisting}[
  caption={\artigo{Verificação de espaço disponível no \textit{harness} do GPS uBlox}.},
  label={lst:gps-vulneravel},
  numbers=left,
  frame=single,
  language=C]
if (write_len > GPS_DUMP_DATA_SIZE - dump_data->len)
\end{lstlisting}

\artigo{No harness, \texttt{GPS\_DUMP\_DATA\_SIZE} é a constante inteira $200$
(\texttt{\#define}), e o comprimento \texttt{uint8\_t} é promovido a
\texttt{int}; quando \texttt{dump\_data->len} excede a capacidade, a diferença
é negativa ($-7$ para \texttt{len = 207}) e, na comparação com
\texttt{write\_len} (\texttt{size\_t}), converte-se em um valor sem sinal
enorme, tornando a guarda sempre falsa. No driver de produção a expressão
correspondente é \texttt{sizeof(dump\_data->data) - dump\_data->len}
(\texttt{gps.cpp}, linha 741), com buffer de 79 bytes
(\texttt{msg/GpsDump.msg}); ali a subtração já é feita em \texttt{size\_t}, sem
sinal, e daria a volta módulo $2^{32}$ no alvo de voo STM32 de 32 bits ---
\textit{se} \texttt{len} pudesse exceder 79, o que a
Seção~\ref{subsec:gps-inducao} mostra não ocorrer.}

O ESBMC produziu \artigo{\texttt{input\_len = 10} e \texttt{dump\_data->len = 207}
(\texttt{test\_choice = 3}, função \texttt{test\_gps\_real\_bit\_operation};
idêntico nas três execuções)}; o
\texttt{memcpy} subsequente escreve fora do buffer. Classificação:
%% [BANCA - Prof. Yuri] "Isso ai nao e underflow, isso e overflow."
%% DECISAO CONFIRMADA PELO LC: mantida a terminologia da CWE, com nota de
%% rodape reconhecendo a distincao. O CWE-191 do MITRE denomina-se literalmente
%% "Integer Underflow (Wrap or Wraparound)"; trocar o termo desalinharia a
%% dissertacao da classificacao adotada e do artigo submetido.
%% (Caixa \pergunta removida: questao respondida.)
\corr{no \textit{harness}, em que a subtração é feita com sinal e o resultado
negativo é depois convertido para um tipo sem sinal,} \textbf{CWE-195}
\corr{(\textit{Signed to Unsigned Conversion Error})} $\rightarrow$
\textbf{CWE-119} (\artigo{\textit{Improper Restriction of Operations within the
Bounds of a Memory Buffer}}); \corr{na expressão do driver de produção, em que os
dois operandos já são sem sinal, a classe correspondente seria} \textbf{CWE-191}
(\textit{Integer Underflow \artigo{(Wrap or Wraparound)}}).\footnote{\revisao{Registra-se
uma ambiguidade terminológica. No vocabulário de análise numérica,
\textit{underflow} designa o resultado pequeno demais para ser representado, ao
passo que o fenômeno aqui descrito --- uma subtração sem sinal que dá a volta
para um valor muito grande --- corresponderia a \textit{overflow} por
\textit{wraparound}. Adota-se, ainda assim, a denominação da \textit{Common
Weakness Enumeration}, cuja entrada CWE-191 intitula-se \textit{Integer Underflow
(Wrap or Wraparound)} e cobre exatamente esse caso, de modo a preservar a
correspondência entre o texto, a classificação empregada e os relatórios
submetidos à comunidade.}} \corr{A Figura~\ref{fig:ce-gps} reproduz a saída do
verificador para esse contraexemplo.} \corr{O \texttt{dump\_buffer} local do \textit{harness}
corresponde ao buffer \texttt{dump\_data} passado ao driver, e o traço literal
consta do Apêndice~\ref{apend:contraexemplos}.}

\begin{figure}[H]
    \centering
    \definecolor{ceFundo}{HTML}{F7F8FA}\definecolor{ceCinza}{HTML}{5A6070}%
    \definecolor{ceAmarelo}{HTML}{0B5CAD}\definecolor{ceVermelho}{HTML}{B3261E}%
    \begin{tcolorbox}[colback=ceFundo, colframe=black!45, boxrule=0.6pt, arc=2pt,
      left=8pt, right=8pt, top=6pt, bottom=6pt, width=0.95\linewidth,
      fontupper=\ttfamily\scriptsize\linespread{1.45}\selectfont]
    \textcolor{ceVermelho}{\textbf{[Counterexample]}}\\
    \textcolor{ceCinza}{State 1 -{}- test\_gps\_driver.cpp:224 -{}- main}\\
    \hspace*{1em}\fcolorbox{ceAmarelo}{blue!6}{\textcolor{ceAmarelo}{\textbf{test\_choice = 3 (...00000000011)}}}\\
    \textcolor{ceCinza}{State 3 -{}- test\_gps\_driver.cpp:179 -{}- test\_gps\_real\_bit\_operation}\\
    \hspace*{1em}\fcolorbox{ceAmarelo}{blue!6}{\textcolor{ceAmarelo}{\textbf{input\_len = 10 (...00000001010)}}}\\
    \textcolor{ceCinza}{State 8 -{}- test\_gps\_driver.cpp:185 -{}- test\_gps\_real\_bit\_operation}\\
    \hspace*{1em}\fcolorbox{ceAmarelo}{blue!6}{\textcolor{ceAmarelo}{\textbf{dump\_buffer.len = \{ .data=\{...\}, .len=207, ... \}}}}\\
    \textcolor{ceVermelho}{\textbf{Violated property:}}\\
    \textcolor{ceVermelho}{\hspace*{1em}file string.c line 280 function \_\_memcpy\_impl}\\
    \textcolor{ceVermelho}{\hspace*{1em}dereference failure: Access to object out of bounds}\\
    \textcolor{ceVermelho}{\textbf{VERIFICATION FAILED}}\\[-2pt]
    \textcolor{ceCinza}{\rule{\linewidth}{0.3pt}}\\
    \textcolor{ceAmarelo}{harness: 200 - 207 = -7 (int) -> size\_t enorme na compara\c{c}\~ao; guarda sempre falsa}\\
    \textcolor{ceCinza}{ESBMC 8.1.0 + Boolector 3.2.4 -{}- log: medianas\_logs/GPS\_run1.log (igual em run2 e run3)}\\
    \textcolor{ceCinza}{PX4 main, commit 8be2590c15 -{}- Issue \#26866 / PR \#27576}
    \end{tcolorbox}
    \caption[Contraexemplo do GPS uBlox (saída do ESBMC)]{\corr{Saída do ESBMC para o contraexemplo do GPS uBlox, montada a partir do log (ESBMC~v8.1.0, Boolector~3.2.4).}}
    \label{fig:ce-gps}
\end{figure}

Em hardware Pixhawk sem MMU, uma escrita fora dos limites pode corromper
memória adjacente; sob o modelo de memória plana do NuttX ela não
necessariamente levanta uma falha limpa. \revisao{O impacto preciso,
porém, depende do arranjo em tempo de ligação}\artigo{, que não foi analisado
aqui, e da alcançabilidade no alvo, tratada a seguir}\revisao{.}

\artigo{A análise do código-fonte do driver de produção resolve essa questão
negativamente: o laço de escrita limita \texttt{write\_len} ao espaço livre
restante, de modo que o nível de preenchimento nunca excede o buffer antes de o
registro ser publicado e o campo zerado. A violação é, portanto, alcançável no
modelo, mas não no driver de produção. Ela é mantida aqui como a ilustração
mais nítida da tese central da metodologia: um contraexemplo formalmente
válido exige triagem contra o código de produção antes de poder ser chamado de
defeito.} \artigo{Essa conclusão, obtida pelo mantenedor e por leitura do
código, foi depois verificada formalmente (Seção~\ref{subsec:gps-inducao}).}

\begin{tcolorbox}[colback=blue!5!white, colframe=blue!60!black,
title=\textbf{Resposta à QE2 --- Violações em código sem teste unitário direto}]
\corr{Parcial. O BMC expôs, no \textit{harness} do GPS, uma violação de
segurança de memória alcançável no modelo, em uma função
(\texttt{dumpGpsData}) que nenhum teste unitário do PX4 referencia. O estado
violador, porém, é inalcançável no driver de produção
(Seção~\ref{subsec:gps-inducao}); nenhum teste, unitário ou de integração,
poderia exercitá-lo. O que a campanha demonstra para a QE2 é, portanto, mais
restrito: a função não tem cobertura unitária direta, e a análise formal foi o
que estabeleceu tanto a violação no modelo quanto a sua inalcançabilidade na
produção. Na primeira campanha, nenhum defeito do driver foi exposto.}
\end{tcolorbox}

\corr{A ausência de teste unitário direto foi conferida em 09/07/2026 por busca
textual na árvore do PX4 (registro em \texttt{evidencia\_unit\_test\_gps.txt},
no repositório de artefatos). As únicas ocorrências de \texttt{dumpGpsData}
estão em \texttt{src/drivers/gps/gps.cpp} (declaração na linha~303, chamadas nas
linhas~428, 435, 457 e~660, definição na linha~727); nenhum arquivo de teste
inclui \texttt{gps.cpp} nem instancia a classe \texttt{GPS}; e o único teste do
diretório, \texttt{devices/gps-parser-test.cpp}, não menciona a função. O
registro não guarda os comandos exatos nem o \textit{commit}; os números de
linha coincidem com os do \textit{commit} \texttt{8be2590c15}. A busca evidencia
ausência de cobertura unitária direta, e não de exercício indireto por SITL ou
por testes de integração.}

\subsection{Verificação do \textit{patch}}
\label{subsec:patch-gps}

\artigo{O artigo registra uma limitação da verificação do \textit{patch} então
disponível: o CBMC~5.95.1 reportava $0$ de $7$ propriedades violadas no
harness com a guarda do PR~\#27576, mas o harness \textbf{sem} a guarda
reportava o mesmo $0$ de $7$, porque ambos assumiam
\texttt{dump\_buffer.len}~$<$~\texttt{GPS\_DUMP\_DATA\_SIZE} --- pré-condição
que exclui justamente os estados que a guarda trata. A comparação, portanto,
não isolava o efeito da guarda, e não havia log do ESBMC para o harness
corrigido.}

\artigo{Para fechar essa lacuna, a verificação foi refeita em 11/10/2026 com
\textbf{as mesmas pré-condições com e sem o \textit{patch}}, em duas formas da
guarda: (a) a do PR~\#27576 (subtração em \texttt{int}; se não há espaço,
zera \texttt{len} e interrompe o laço); e (b) a forma pedida pelo mantenedor na
revisão do PR (no topo do laço, se \texttt{len} excede a capacidade, zera
\texttt{len} e prossegue). A Tabela~\ref{tab:patch-gps} resume os
desfechos.}

\begin{table}[H]
\centering
\caption{\artigo{Verificação do \textit{patch} do GPS sob as mesmas pré-condições}}
\label{tab:patch-gps}
\footnotesize
\begin{tabular}{|p{6.2cm}|c|c|c|}
\hline
\textbf{Ferramenta e harness} & \textbf{Sem \textit{patch}} & \textbf{PR \#27576} & \textbf{Forma do mantenedor} \\
\hline
\artigo{ESBMC~8.1.0 + Boolector~3.2.4; harness da Tabela~\ref{tab:resultados}, mesmas \textit{flags}} &
\artigo{\textit{violação}} & \artigo{sem violação} & \artigo{sem violação} \\
\artigo{CBMC~5.95.1; janela $\texttt{len} \geq 190$; \texttt{--unwind} 2, 3 e 5} &
\artigo{\textit{violação}} & \artigo{sem violação} & \artigo{sem violação} \\
\hline
\end{tabular}

\footnotesize
\begin{flushleft}
\artigo{ESBMC: sem \textit{patch}, 4.463 VCCs e o mesmo contraexemplo da
Tabela~\ref{tab:resultados} (\texttt{input\_len = 10},
\texttt{dump\_buffer.len = 207}); PR~\#27576, 3.453 VCCs; forma do mantenedor,
4.760 VCCs. CBMC: sem \textit{patch}, 1 de 73 propriedades violadas
(\textit{memcpy destination region writeable}) nos três limites; com as
guardas, 0 de 73 e 0 de 79. As cópias do harness com \textit{patch} diferem do
original apenas no laço de \texttt{dumpGpsData()}. O ESBMC~8.1.0 usado foi
recompilado da \textit{tag} \texttt{v8.1} com Boolector~3.2.4 e validado contra
os logs de julho (mesmas contagens de VCC e mesmos contraexemplos). Execuções
em máquina de nuvem: só os vereditos são reportados; os tempos não são
comparáveis aos da Tabela~\ref{tab:resultados}.}
\end{flushleft}
\end{table}

\artigo{Sob as mesmas pré-condições, as duas formas da guarda eliminam a
violação que o harness sem \textit{patch} exibe, nas duas ferramentas. O
escopo continua limitado: são verificações até o limite de desenrolamento, sem
asserções de desenrolamento, como na Tabela~\ref{tab:resultados}; e, como a
próxima seção mostra, o estado que a guarda trata não ocorre no driver de
produção, de modo que o \textit{patch} só se justifica como endurecimento
defensivo --- exatamente o enquadramento dado pelo mantenedor.}

\subsection{\artigo{Inalcançabilidade no driver de produção: prova indutiva}}
\label{subsec:gps-inducao}

\artigo{O argumento do mantenedor --- que \texttt{len} nunca excede a
capacidade no programa real --- foi verificado formalmente sobre uma cópia
literal do laço de \texttt{GPS::dumpGpsData()} (\texttt{gps.cpp}, linhas
727--761 do \textit{commit} \texttt{8be2590c15}), com os tipos reais:
\texttt{len} em \texttt{size\_t}, \texttt{sizeof(dump\_data->data)}~$= 79$ e
\texttt{dump\_data->len} em \texttt{uint8\_t}. A prova é por indução sobre o
invariante $I$: \texttt{dump\_data->len}~$< 79$.}

\begin{itemize}
\item \artigo{\textbf{Base.} \texttt{GPS::initializeCommunicationDump()} zera os
dois buffers com \texttt{memset} (linhas 717--718), logo $I$ vale antes da
primeira chamada.}
\item \artigo{\textbf{Passo.} A partir de um estado arbitrário que satisfaz $I$
e de um comprimento de entrada arbitrário ($1 \leq \texttt{len} \leq 4096$), uma
iteração do laço não acessa fora dos limites no \texttt{memcpy} e preserva $I$.
Verificado com o CBMC~5.95.1 (0 de 124 propriedades violadas) e com o
ESBMC~8.1.0 + Boolector~3.2.4 (447 VCCs, \textsc{verification successful}, com
asserções de desenrolamento \textbf{ligadas}).}
\item \artigo{\textbf{Premissa.} \texttt{dumpGpsData()} é a única função que
escreve \texttt{dump\_data->len}: no \textit{commit} verificado, as escritas
estão nas linhas 747--757, e os ponteiros para os buffers são membros privados
da classe.}
\end{itemize}

\artigo{Como controle negativo, o mesmo corpo sem a hipótese $I$ é refutado
pelas duas ferramentas (acesso fora dos limites no \texttt{memcpy}), o que mostra
que o harness detectaria o defeito se o invariante não valesse. Diferentemente
dos resultados da Tabela~\ref{tab:resultados}, este não depende de limite de
desenrolamento: o harness não tem laço próprio, e o único laço (o modelo de
\texttt{memcpy}, de no máximo 79 iterações) é coberto com asserção de
desenrolamento ativa. Logo, em toda execução do driver real o \texttt{memcpy}
permanece dentro do buffer, e o estado violador do harness
(\texttt{len = 207}) é inalcançável na produção. O resultado transforma a
triagem por leitura de código em evidência formal e ilustra a etapa~(v) da
metodologia aplicada ao próprio contraexemplo.}

\corr{A premissa de escritor único foi conferida por busca dos dois ponteiros
em \texttt{gps.cpp} no \textit{commit} verificado. \texttt{\_dump\_to\_device} e
\texttt{\_dump\_from\_device} aparecem apenas na declaração (linhas 224--225),
na liberação (410--411), na alocação e no teste de nulo (709--712), no
\texttt{memset} (717--718) e na seleção do buffer dentro de
\texttt{dumpGpsData()} (729). Não são passados a outra função nem expostos por
método de acesso, e a publicação recebe o conteúdo por referência para cópia;
não há, portanto, outro nome pelo qual \texttt{len} possa ser escrito
(\textit{aliasing}).}

\corr{A conclusão depende de três hipóteses, que a prova não verifica: (i)~as
chamadas a \texttt{dumpGpsData()} sobre um mesmo buffer não se intercalam, isto
é, não há escrita concorrente; (ii)~a memória do buffer não é corrompida por um
defeito em outra parte do \textit{firmware}; e (iii)~compilador e
\textit{hardware} executam o código conforme a semântica da linguagem. Sob
qualquer violação dessas hipóteses o invariante pode deixar de valer, e a guarda
proposta no PR~\#27576 passaria a ter efeito.}

\section{Discussão}
\label{sec:discussao-resultados}

\subsection{Configuração, \textit{soundness} e tempo}

\begin{tcolorbox}[colback=blue!5!white, colframe=blue!60!black,
title=\textbf{Resposta à QE3 --- Configurações por módulo}]
\corr{Nove dos doze módulos verificam-se sob a invocação padrão, sem qualquer
ajuste. Apenas três exigiram configuração específica: MAVLink com
\texttt{--unwind~300} (laços profundos de \textit{parsing}), Navigator com
\texttt{--unwind~15} e VTOL com \texttt{--no-library} (desabilita a biblioteca C
abstrata). Todas estão documentadas no repositório de artefatos. O custo de
configuração é, portanto, baixo: a maioria dos módulos não requer
sintonia alguma.}
\end{tcolorbox}

%% LC (resolvido 12/11/2026): o esforço humano não foi registrado; a afirmação
%% de baixo custo foi restrita ao tempo de máquina e à configuração.
\corr{Essa resposta refere-se apenas à configuração do verificador. O esforço de
escrever cada \textit{harness}, triar os contraexemplos e conferir o código de
produção não foi registrado durante a pesquisa, e por isso não é reportado em
horas por módulo; é esse esforço, e não o tempo de máquina, que domina o custo
de adoção (Seção~\ref{sec:conclusao-limitacoes}).}

%% [BANCA - Prof. Barreto] "Ha uma tendencia do pessoal de metodos formais em achar que
%% o teste e pior que a verificacao formal. Cada um tem o seu lugar."
\revisao{Convém evitar uma assimetria comum na literatura de métodos formais: apontar a
cobertura limitada dos testes sem reconhecer que a verificação formal está sujeita à mesma
limitação, apenas por outra via. A cobertura do teste é limitada pelos casos que alguém
escreveu; a do BMC é limitada por três fatores igualmente humanos --- as \emph{propriedades}
que se escolheu especificar, a \emph{profundidade} até a qual se decidiu desenrolar, e o
\emph{recorte} de código submetido ao verificador. Uma falha de lógica de negócio para a qual
nenhuma propriedade foi formulada passa despercebida por um verificador tão silenciosamente
quanto por uma suíte de testes que não a exercita.}

\revisao{A diferença entre as duas abordagens não é, portanto, de completude, mas de
\emph{natureza da garantia}: dentro do recorte especificado, o BMC cobre exaustivamente todas
as entradas possíveis até o limite, enquanto o teste cobre exatamente as entradas escolhidas.
Cada uma alcança o que a outra não alcança, e é nesse sentido preciso que são complementares
--- e não em uma suposta superioridade de uma sobre a outra.}

O BMC é \textit{sound} \artigo{e completo até o limite $k$}: violações reportadas são
reais para o modelo codificado, mas a ausência de violações até a profundidade
$k$ não é prova~\cite{Biere1999}.
Dentro desse escopo, BMC e teste unitário são
complementares~\cite{Godefroid2005}.

\corr{Os dois módulos mais caros da campanha, à parte o GPS, são o MAVLink
($49{,}2$\,s) e o EKF2 ($12{,}2$\,s). No caso do MAVLink, o custo acompanha o
limite de desenrolamento elevado: com $k = 300$, cada iteração replica as
instruções do laço na fórmula SMT, conforme discutido na
Seção~\ref{subsec:escolha-k}. Não se atribui aqui o custo ao \textit{bit-blasting}
de IEEE-754~\cite{Gadelha2020}, porque os dados da Tabela~\ref{tab:resultados}
não sustentam essa explicação: o Expo RC é aritmética de ponto flutuante pura e
verifica-se em $4{,}9$\,s, o Barômetro em $0{,}4$\,s e o VTOL, que emprega
\texttt{sinf} e \texttt{sqrtf}, em $4{,}0$\,s. Isolar a contribuição do ponto
flutuante exigiria execução controlada adicional, registrada como trabalho
futuro.}

\subsection{\revisao{Generalidade do padrão encontrado}}

\revisao{Dois outros candidatos com subtração próxima a \texttt{memcpy}
foram inspecionados (o \textit{statustext} do MAVLink e o RC Update) e não
apresentam \textit{underflow} desprotegido. O padrão do GPS, portanto, não é
disseminado nos módulos inspecionados.}

\lc{Método de inspeção não informado; dizer se ESBMC ou leitura, e listar locais no código.}

\subsection{\revisao{Observação entre autopilotos e integração contínua}}
\label{subsec:ardupilot-ci}

\revisao{A mesma observação de ausência de validação de faixa \artigo{(CWE-1284,
\textit{Improper Validation of Specified Quantity in Input}, subclasse de
CWE-20)} aparece no driver equivalente do ArduPilot~\cite{ArduPilot}: a execução do
CBMC~5.95.1 sobre \texttt{AP\_InertialSensor\_BMI088.cpp} (linhas $383$--$385$)
a reproduz, registrada na Issue~\#33272\footnote{\url{https://github.com/ArduPilot/ardupilot/issues/33272}}.
Como os mantenedores do PX4 recusaram essa mesma observação, ela é reportada
apenas como verificação de consistência entre autopilotos, e \textbf{não} como
achado independente.} \artigo{Um mantenedor do ArduPilot encerrou a
\textit{issue} citando a faixa de medição de $-104$ a $150\,^{\circ}$C do
\textit{datasheet}; os dois projetos, portanto, recusaram a observação.}
\revisao{Vale notar que o ArduPilot já protege o contador do FIFO
(\texttt{fifo\_length}~$\leq 8 \times 7 = 56$ bytes,} \artigo{linhas 314--315 do
ramo \textit{master} consultado em 12/10/2026}\revisao{).}

\revisao{Um fluxo de trabalho do GitHub Actions em um \textit{fork} do
PX4\footnote{\url{https://github.com/marcela2881/PX4-Autopilot}} executa os
harnesses do IMU e do GPS a cada \textit{push} e \textit{pull request},
integrando verificação baseada em harness a um \textit{pipeline} de CI padrão.
Os tempos de relógio da CI não foram medidos sob configuração controlada.}

%% [BANCA - Prof. Yuri] "A metodologia promete avaliar impacto e consequencias da falha
%% no contexto de voo, mas eu nao senti que isso esta sendo feito nos exemplos."
%% [BANCA - Prof. Reginaldo] "Nao e trivial fazer essa analise de criticidade. Para essa
%% temperatura, sendo simples, da ate para comentar."
%% DECISAO CONFIRMADA PELO LC: a analise e feita onde e tratavel (temperatura),
%% com o criterio do Prof. Reginaldo, e declara-se por que nao se estende aos
%% demais. O LC respondeu "vale tentar algo mais ambicioso" quanto ao SITL, que
%% fica registrado como trabalho futuro por nao caber no cronograma ate a defesa.
%% (Caixa \pergunta removida: questao respondida.)
\section{\revisao{Análise de Severidade e Plausibilidade Operacional}}
\label{sec:severidade}

\revisao{Um contraexemplo demonstra que existe uma execução que viola a propriedade,
mas nada diz sobre duas questões que interessam a quem opera o veículo: \emph{quão grave} é a
violação e \emph{quão plausível} é que o estado que a produz ocorra em voo. O verificador não
ordena achados por severidade nem avalia verossimilhança operacional --- essa etapa é humana e
constitui parte da triagem. Esta seção a executa onde ela é tratável e explicita por que não
se estende a todos os casos.}

\subsection{\revisao{Temperatura do BMI088}}

\revisao{Para esta propriedade a análise de severidade é decidível a partir do próprio
código, e reduz-se a uma pergunta: existe, a jusante do driver, algum mecanismo que sinalize a
leitura anômala?}

\revisao{Se a pilha de software emite alarme ao ultrapassar os $85\,^\circ$C
especificados, então o fato de a variável admitir $86{,}25\,^\circ$C é apenas margem de
representação, sem consequência operacional: o sistema já sabe que a leitura é inválida. Se,
ao contrário, nenhuma camada sinaliza a excursão, a situação é qualitativamente distinta ---
o veículo pode estar operando fora da faixa certificada do componente sem que isso seja
comunicado ao operador ou à lógica de \textit{failsafe}. A responsabilidade em jogo não é a de
impedir o sobreaquecimento, que é fenômeno físico, mas a de \emph{medir e informar} a
temperatura, que é responsabilidade de software.}

\revisao{No PX4} \artigo{(\textit{commit} \texttt{8be2590c15})}\revisao{, a temperatura do BMI088 não alimenta o caminho inercial e não
existe alarme dedicado a excursões dessa grandeza no nível do driver. A severidade, portanto,
não é a de perda de controle, mas a de \textbf{perda de observabilidade}: uma leitura fora de
faixa é consumida silenciosamente. Essa é precisamente a razão pela qual a observação foi
mantida como nota de projeto após a recusa upstream --- é um achado de qualidade de
instrumentação, não de segurança de voo, e essa distinção só se torna visível quando a análise
de severidade é feita.}

\corr{Os consumidores da leitura foram levantados por busca textual no
\textit{commit} analisado. O driver publica a temperatura depois de testar
apenas se o valor é finito (\texttt{PX4\_ISFINITE}); não há comparação com a
faixa do \textit{datasheet}. A jusante, o campo é lido em quatro lugares:
\texttt{sensors/vehicle\_imu}, que calcula a média e a publica em
\texttt{vehicle\_imu\_status}; os fluxos \texttt{SCALED\_IMU} do MAVLink, que a
enviam à estação de solo; \texttt{drivers/heater}, que a usa como realimentação
do aquecedor da IMU; e \texttt{modules/temperature\_compensation}, que, quando
habilitado por parâmetro (\texttt{TC\_A\_ENABLE}), calcula a partir dela a
correção de \textit{offset} do acelerômetro. Nenhum desses pontos compara o
valor com a faixa de $-40$ a $+85\,^{\circ}$C. Dois deles qualificam a afirmação
de que a temperatura não alimenta o caminho inercial: com a compensação térmica
habilitada ela influencia a correção aplicada à medida, e no controle do
aquecedor ela fecha uma malha. O efeito de uma leitura fora de faixa nesses dois
caminhos não foi analisado nesta dissertação.}

\subsection{\revisao{Plausibilidade operacional da violação do GPS}}

\revisao{Para o GPS a análise inicialmente ficou incompleta. A
violação exige que \texttt{dump\_data->len} exceda 200 no momento da chamada. Como o campo é
acumulado entre chamadas sucessivas, e como o recurso de \textit{dump} é ativado por
parâmetro, estabelecer plausibilidade exigiria modelar a sequência de chamadas a montante ---
isto é, demonstrar que os módulos abstraídos pelos \textit{stubs} podem efetivamente conduzir
o driver a esse estado.}

\artigo{Essa modelagem foi feita posteriormente, não por simulação da sequência de
chamadas a montante, mas por análise direta do laço de escrita do próprio driver: ele
limita \texttt{write\_len} ao espaço livre restante antes de cada incremento de
\texttt{dump\_data->len}, de modo que o campo nunca excede a capacidade do buffer
antes de o registro ser publicado e zerado. A pré-condição que o \textit{harness} assume
--- \texttt{dump\_data->len} já acima de 200 ao entrar na verificação --- não tem
como surgir no driver de produção, porque é o próprio laço de escrita que mantém
esse invariante.}

\revisao{A severidade, caso o estado fosse atingido, seria alta: em hardware sem
MMU uma escrita fora dos limites pode corromper estado adjacente de voo.}
\artigo{A plausibilidade, no entanto, está resolvida: o estado que dispara a
violação é inalcançável no driver de produção. É essa combinação --- consequência
potencialmente severa, mas restrita ao modelo do \textit{harness} --- que melhor
descreve o achado, e que motivou a mudança de enquadramento junto ao mantenedor,
de vulnerabilidade para endurecimento defensivo (PR~\#27576).}

\artigo{A confirmação por simulação \textit{software-in-the-loop}, cogitada
como caminho para estabelecer plausibilidade, deixou de ser necessária para
este achado específico: a análise do código-fonte a resolveu diretamente. SITL
permanece de interesse para casos futuros em que a leitura do código não for
conclusiva por si só, e fica registrada nesse sentido mais geral entre as
direções de trabalho futuro.}

\subsection{\revisao{Por que a análise não se estende a todos os módulos}}

\revisao{Dos doze módulos, dez não produziram violação, de modo que não há severidade a
avaliar. Entre os três contraexemplos, o candidato do FIFO foi refutado \artigo{na conferência contra o código de produção} e
sua severidade é, por construção, nula. Restam os dois casos analisados acima. A análise de
criticidade em sentido pleno --- estabelecer a relação causal entre a violação e o desfecho da
missão --- exigiria modelo da dinâmica do veículo e da propagação de erro entre módulos,
elemento que está fora do escopo desta dissertação e que se registra como direção de trabalho
futuro (Seção~\ref{sec:trabalhos-futuros-hw}).}

\section{\corr{Resultados da Segunda Campanha}}
\label{sec:resultados-segunda-campanha}

\corr{A segunda campanha aplicou o método a dez módulos complementares, com
verificação sobre o código original do PX4, validação cruzada com o CBMC em
todos os \textit{harnesses} e \textit{replay} sob sanitizadores. O protocolo,
as tabelas por frente e os registros estão na
Seção~\ref{sec:segunda-campanha}; esta seção reúne os desfechos
(Tabela~\ref{tab:desfechos-segunda}) e a validação em bancada.}

\begin{table}[!htb]
\centering
\caption{\corr{Desfechos da segunda campanha, por módulo}}
\label{tab:desfechos-segunda}
\footnotesize
\begin{tabular}{|p{5.6cm}|c|p{7.4cm}|}
\hline
\textbf{Desfecho} & \textbf{Módulos} & \textbf{Quais} \\
\hline
\corr{Defeito reproduzido no código real e corrigido \textit{upstream} por terceiros} & \corr{3} & \raggedright\corr{\texttt{externalChecks}, \texttt{ft7\_technologies}, \texttt{IridiumSBD}} \tabularnewline
\hline
\corr{Defeito reproduzido e reportado neste trabalho} & \corr{2} & \corr{RoboClaw (\textit{Issue}~\#29785), batmon (\textit{Issue}~\#29787)} \\
\hline
\corr{Defeito reproduzido, relatado por terceiro e não corrigido} & \corr{1} & \corr{\texttt{sumd}} \\
\hline
\corr{Defeito reproduzido, não reportado (impacto desprezível)} & \corr{1} & \corr{\texttt{MicroStrain}} \\
\hline
\corr{Falso positivo de modelagem, barrado na conferência} & \corr{1} & \corr{\texttt{bst}} \\
\hline
\corr{Sem defeito de memória} & \corr{2} & \corr{\texttt{logger}, \texttt{Common}} \\
\hline
\end{tabular}
\end{table}

\corr{Em sete dos dez módulos o efeito previsto pelo verificador foi reproduzido
executando o código real do PX4. Em três deles, o próprio projeto corrigiu o
mesmo ponto de forma independente, o que corrobora os achados sem que fosse
preciso reportá-los. O único falso positivo, no \texttt{bst}, foi barrado pela
conferência antes de qualquer relato público. O contraste com a primeira
campanha é direto: lá, três contraexemplos não resultaram em defeito confirmado
algum; aqui, nenhum candidato chegou a ser reportado sem ter sido antes
reproduzido no código real. Os dois relatos abertos aguardavam triagem na data
de fechamento deste texto. A campanha rendeu também uma correção fora do PX4: um
defeito no modelo de \texttt{snprintf} do ESBMC, reportado e corrigido pelos
mantenedores da ferramenta.}

\subsection{\corr{Validação em bancada}}
\label{subsec:validacao-bancada}

\corr{Dois ensaios foram feitos numa Pixhawk~6X~RT com \textit{firmware}
compilado do \textit{commit} analisado. Eles respondem a perguntas opostas: um
mostra que um estado apontado pelo verificador \emph{não} ocorre na placa, e o
outro, que um defeito apontado \emph{ocorre}.}

\corr{\textbf{GPS (primeira campanha).} Com o registro da comunicação do GPS
habilitado (\texttt{GPS\_DUMP\_COMM}~$= 1$), foram coletadas 175 mensagens
\texttt{gps\_dump} em dois registros gravados pela placa em bancada, analisadas com
\texttt{pyulog}. Em todas, o campo \texttt{len} vale 207, isto é, o bit~7
(que marca dado enviado ao dispositivo) somado a 79~bytes, a capacidade do
buffer: nenhuma mensagem excede a capacidade. O número coincide por acaso com o
\texttt{len}~$= 207$ do contraexemplo, mas tem outro significado: no
\textit{harness}, 207 é um nível de preenchimento acima de um buffer de
200~bytes; na placa, é um buffer de 79~bytes exatamente cheio. O ensaio é
consistente com a prova da Seção~\ref{subsec:gps-inducao}. Seu limite é que só
o sentido da placa para o receptor foi exercitado, sem receptor GPS conectado.}

\corr{\textbf{RoboClaw (segunda campanha).} Com um ESP32 no papel do
controlador de motores, o valor de \textit{failsafe} em 257 e a terminação de
voo acionada, a placa transmitiu o byte de dados 130, fora da faixa de 0 a 127
aceita pelo protocolo; com o valor em 254, transmitiu 127; e, com a correção
aplicada, voltou a transmitir 127 mesmo com 257. Os pacotes, as contagens e os
limites do ensaio (controlador emulado, \textit{firmware} com o driver
habilitado e dois membros inicializados, apenas o primeiro dos dois defeitos)
estão na Seção~\ref{subsec:2c-roboclaw}.}

\corr{Os dois ensaios cobrem, cada um, um único achado. Os demais resultados
desta dissertação apoiam-se na verificação formal e, na segunda campanha, no
\textit{replay} do código real em computador.}

\section{Ameaças à Validade}
\label{sec:ameacas-validade}

\subsection{Validade interna}

Os tempos reportados são medianas de três execuções e devem ser lidos apenas
como ordens de grandeza. \corr{Na fase exploratória, um resultado do MAVLink com \texttt{--unwind~10}
foi descartado como artefato de configuração: era desenrolamento insuficiente,
resolvido com \texttt{--unwind~300} e confirmado pela reexecução com as
asserções ligadas (Seção~\ref{subsec:escolha-k}). Um segundo candidato inicial,
no EKF2, foi descartado na mesma fase, mas não há registro que permita
descrevê-lo; por isso ele não é contado nem discutido neste trabalho.}
\revisao{O \textit{underflow} do GPS é descarregado apenas pelo
Boolector dentro do \textit{timeout}, de modo que o achado é sensível ao
solver; a validação cruzada com CBMC mitiga artefatos específicos de
ferramenta.}
\artigo{Com o Z3, o ESBMC não produziu veredito dentro do \textit{timeout} no
GPS, nem no MAVLink e no EKF2 com \texttt{--unwind 10}, nas re-execuções com a
versão 8.1 (\texttt{validacao\_esbmc\_8.1/} no repositório); e a validação
cruzada com o CBMC é parcial no GPS, pois só termina com limites reduzidos.}

\corr{O limite de desenrolamento também é ameaça à validade interna.} A verificação foi limitada a $k = 10$ e, com
\texttt{--no-unwinding-assertions} habilitado, laços que excedem o limite são
truncados silenciosamente --- de modo que um veredito \textsc{successful} pode
ocultar falsos negativos mesmo abaixo da profundidade operacional pretendida.
\artigo{A reexecução com as asserções ligadas mostrou que isso não ocorre nos dez
módulos sem violação (Tabela~\ref{tab:unwinding}); no GPS, porém, o laço do
\texttt{memcpy} excede $k = 10$, e a análise desse módulo não o cobre por
inteiro. Permanece, em todos os casos, o limite do recorte: o veredito vale
para o \textit{harness}, e não para o módulo inteiro.}
Os totais de VCC por módulo misturam limites de desenrolamento (10, 15 e 300) e
não são diretamente comparáveis.

\subsection{Validade externa}

Selecionar módulos por viabilidade de \textit{stubbing} sub-representa módulos
com acoplamento pesado a hardware ou ao escalonador, de modo que a amostra não
é estatisticamente representativa do PX4; a mitigação é cobrir as principais
classes de entrada não confiável. A análise modular não captura bugs de
integração decorrentes de interações concorrentes, e os barramentos CAN/UAVCAN,
a ponte RTPS/uXRCE-DDS e os drivers de fluxo óptico e telêmetro permanecem não
verificados.

\subsection{Validade de construto}

\corr{A questão de construto é se as propriedades verificadas medem o que se
pretende afirmar. Três ressalvas se aplicam. Primeiro, as asserções foram
escritas pela autora a partir de \textit{datasheets} e da documentação, e uma
delas, a faixa de temperatura do BMI088, foi contestada pelos mantenedores como
requisito de software: violar a asserção não equivale a violar um requisito do
projeto. Segundo, a asserção final do \textit{harness} do GPS é tautológica, e a
propriedade efetivamente verificada ali é a checagem automática de limites do
\texttt{memcpy} (Capítulo~\ref{cap:metodologia}). Terceiro, toda propriedade é
verificada sobre o \textit{harness}, que transcreve e adapta o código: o
veredito mede a segurança do modelo, e só a conferência contra o código de
produção permite estendê-lo ao driver.}
\corr{A comparação entre ESBMC com Boolector e CBMC com MiniSAT varia
ferramenta e \textit{backend} conjuntamente, de modo que a propriedade distinta
que cada uma reporta sobre o defeito do GPS reflete diferença de
\textit{toolchain}, e não de SMT versus SAT isoladamente.}

\subsection{\corr{Validade de conclusão}}

\corr{Cada módulo foi executado três vezes e reporta-se a mediana, sem medida de
dispersão. Isso sustenta afirmações de ordem de grandeza --- ``menos de noventa
segundos'' --- mas não comparações finas entre módulos cujos tempos diferem por
poucos segundos. Não há, além disso, \textit{baseline} de comparação: nenhuma
outra técnica foi executada sobre os mesmos doze módulos para contraste, à
exceção do CBMC em dois casos. Segue-se que a afirmação de viabilidade é
absoluta, e não relativa: demonstra-se que a verificação é barata em termos
próprios, não que seja mais ou menos barata que alternativas.}

\subsection{\corr{Construção dos \textit{harnesses} com apoio de ferramenta de IA}}

\corr{Conforme registrado na Declaração de Uso de IA, os \textit{harnesses}
foram inicialmente montados com apoio de assistente de IA, a partir das regras
de instrumentação do ESBMC e do código-fonte de cada módulo, e em seguida
revisados manualmente pela autora. Trata-se de ameaça legítima à validade e
merece registro explícito: um \textit{harness} gerado dessa forma pode
reproduzir padrões plausíveis porém incorretos --- em particular, \textbf{omitir
guardas presentes no código de produção}, que foi exatamente o mecanismo do
falso positivo do contador de FIFO discutido na Seção~\ref{subsec:imu-fifo}.}

\corr{A mitigação praticada tem duas camadas. A primeira é a revisão manual de
cada \textit{harness} contra o código-fonte do módulo antes da execução. A
segunda, e mais importante, é a etapa obrigatória de \artigo{conferência} contra o
driver não instrumentado \artigo{(por \textit{replay} ou leitura do código)}, que capturou o único caso em que essa ameaça se
materializou. O episódio do FIFO é, nesse sentido, evidência de que o método
dispõe de mecanismo de refutação embutido, e não apenas de detecção.}

\subsection{\revisao{Correção dos \textit{stubs}}}

\revisao{Os \textit{stubs} que retornam \texttt{nondet\_T()}
sobreaproximam apenas as \textbf{entradas} de hardware, não a lógica do driver
ao redor: nenhuma violação é perdida por conta da abstração de entrada, mas
quando o harness omite guardas de produção --- como mostrou o candidato do FIFO
--- o modelo admite estados que o driver real rejeita. A legitimidade de um
achado repousa, portanto, sobre a etapa obrigatória de \artigo{conferência} contra o
driver não instrumentado.}

\revisao{Inversamente, os \textit{stubs} \textbf{sub}aproximam o
comportamento interno do hardware (máquinas de estado I2C e SPI, escalonamento
do NuttX); defeitos disparados por esse estado interno não seriam detectados, o
que torna a validação com hardware no laço um complemento necessário.}

\subsection{\corr{Defeitos no verificador ou no solver}}

\corr{Um veredito depende da correção do próprio verificador, do seu modelo das
bibliotecas e do solver. A ameaça não é hipotética: durante a segunda campanha,
o modelo de \texttt{snprintf} do ESBMC produziu um alarme falso, depois
confirmado e corrigido pelos mantenedores da ferramenta
(Seção~\ref{sec:segunda-campanha}). A mitigação praticada é parcial. Na
primeira campanha, o CBMC, com outro \textit{front-end} e outro solver, foi
executado apenas sobre os dois módulos com contraexemplo; os dez módulos sem
violação foram verificados com uma única ferramenta e um único solver, de modo
que um defeito que suprimisse uma violação neles passaria despercebido. Na
segunda campanha, a validação cruzada com o CBMC cobre todos os
\textit{harnesses}, e os candidatos alcançáveis foram ainda reproduzidos fora
dos verificadores, por \textit{replay} sob sanitizadores.}

\section{\corr{Lições Aprendidas}}
\label{sec:licoes}

\corr{Quatro lições atravessam as duas campanhas.}

\corr{\textbf{Fidelidade do \textit{harness}.} Os dois falsos positivos deste
trabalho, o contador de FIFO do BMI088 e o \texttt{bst}, têm a mesma origem: o
\textit{harness} omitia guardas que o código de produção contém. Transcrever
trechos é rápido, mas o veredito vale para a cópia. A segunda campanha passou a
verificar o código original e a exigir \textit{replay} da função real, e foi
essa exigência que barrou o \texttt{bst}.}

\corr{\textbf{Reporte prematuro.} Na primeira campanha, as \textit{issues} do
GPS e do FIFO foram abertas antes da conferência contra o código de produção, e
a conferência depois enfraqueceu ou refutou os dois candidatos. A regra adotada
em seguida foi reportar somente depois da conferência, do \textit{replay}, da
consulta à especificação com página citada, da busca por relatos duplicados e
da verificação do estado atual do código. O caso do RoboClaw acrescentou um
item: conferir também em qual ramo o código analisado ainda existe, pois o
driver havia sido reescrito no ramo principal entre a análise e o relato.}

\corr{\textbf{Projeto das hipóteses de entrada.} A pré-condição decide o que o
verificador pode encontrar, nos dois sentidos. No GPS, admitir um estado inicial
acima da capacidade produziu uma violação que o driver real não alcança. No
RoboClaw, restringir a saída à faixa declarada no \texttt{module.yaml}, de 0 a
257, foi o que expôs o defeito; uma faixa mais estreita, de aparência razoável,
o teria ocultado. Cada pré-condição precisa de origem documentada no código ou
na especificação.}

\corr{\textbf{Limites da ferramenta.} Três comportamentos do verificador
exigiram atenção: o truncamento silencioso de laços com
\texttt{--no-unwinding-assertions}, a sensibilidade do contraexemplo do GPS ao
solver e um defeito no modelo de uma função de biblioteca. Nenhum invalidou os
resultados, mas todos reforçam que um veredito isolado, de uma única ferramenta
e configuração, é ponto de partida para a triagem, e não conclusão.}


\chapter{Conclusão}
\label{cap:conclusao}

%% LC (resolvido 12/11/2026): resposta explícita à questão de pesquisa e ligação
%% de cada objetivo à sua evidência.
A questão que orientou esta pesquisa foi se é viável aplicar \textit{Bounded
Model Checking} a um \textit{firmware} de controle de voo de porte industrial,
com tempo de verificação por módulo compatível com o uso prático e produzindo
evidências auditáveis sobre propriedades de memória e de aritmética. A resposta
é afirmativa, com um alcance que precisa ser dito junto: a viabilidade foi
demonstrada por módulo, sobre \textit{harnesses} que isolam cada unidade, e não
para o sistema completo. Vinte e dois módulos do PX4 Autopilot foram
verificados em duas campanhas; cada verificação da primeira concluiu em menos
de noventa segundos, e todo veredito pode ser refeito a partir dos artefatos
publicados.

Os três objetivos específicos foram atendidos da seguinte forma. O objetivo~O1,
propor a metodologia, corresponde às oito etapas do
Capítulo~\ref{cap:metodologia}. O objetivo~O2, aplicá-la e classificar os
candidatos, corresponde às duas campanhas do Capítulo~\ref{cap:experimentos}:
três contraexemplos triados na primeira e dez módulos com desfecho individual
na segunda (Tabelas~\ref{tab:resultados} e~\ref{tab:desfechos-segunda}). O
objetivo~O3, avaliar a viabilidade, corresponde aos tempos e às configurações
por módulo da Tabela~\ref{tab:resultados} e às respostas às questões
experimentais.

\section{Síntese dos Resultados}
\label{sec:conclusao-sintese}

%% LC (resolvido 12/11/2026): síntese condensada em um parágrafo por achado.
\textbf{Dez módulos sem violação.} Na primeira campanha, os
\textit{harnesses} de dez dos doze módulos foram verificados sem violação das
propriedades de segurança de memória e de \textit{overflow} habilitadas. A
reexecução com as asserções de desenrolamento ligadas mostrou que nenhum laço
desses \textit{harnesses} excede o limite usado, de modo que o resultado não
depende de truncamento silencioso.

\textbf{Temperatura do BMI088.} O verificador produziu uma leitura de
$86{,}25\,^{\circ}$C, acima do limite operacional de $85\,^{\circ}$C do
\textit{datasheet}. Os mantenedores recusaram a observação como defeito, e ela
permanece como especificação disputada.

\textbf{Contador do FIFO do BMI088.} O contraexemplo, com
\texttt{count}~$= 9088$ para um FIFO de $1.024$~bytes, era um falso positivo: o
driver de produção contém três guardas que o \textit{harness} omitia.

\textbf{GPS uBlox.} O verificador encontrou, no \textit{harness}, uma escrita
fora dos limites em \texttt{dumpGpsData()}. Uma prova indutiva sobre o laço real
do driver mostrou que o estado violador é inalcançável na produção, e um ensaio
em bancada foi consistente com essa conclusão. A guarda proposta foi enquadrada
pelo mantenedor como endurecimento defensivo.

\textbf{Segunda campanha.} Em sete de dez módulos complementares, o efeito
previsto foi reproduzido executando o código real do PX4 sob sanitizadores.
Três desses defeitos foram corrigidos pelo próprio projeto, de forma
independente; dois foram reportados neste trabalho (RoboClaw e batmon) e
aguardavam triagem no fechamento do texto; um já havia sido relatado por
terceiro; e um não foi reportado, pelo impacto desprezível. Um falso positivo
foi barrado pela conferência antes de qualquer relato.

\textbf{Validação em bancada.} Numa Pixhawk~6X~RT, o driver do RoboClaw
transmitiu o byte de dados 130, fora da faixa de 0 a 127 do protocolo, e deixou
de fazê-lo com a correção aplicada. O controlador de motores foi emulado.

\textbf{Correção na ferramenta.} A campanha expôs um defeito no modelo de
\texttt{snprintf} do ESBMC, confirmado e corrigido pelos mantenedores.

\section{Contribuições}
\label{sec:conclusao-contribuicoes}

%% LC (resolvido 12/11/2026): lista alinhada às cinco contribuições da
%% Introdução, com a numeração (i) a (viii) das etapas.
As contribuições são as cinco enunciadas na Introdução
(Seção~\ref{sec:contribuicoes}).

\begin{enumerate}
\item \textbf{Metodologia em oito etapas.} Seleção de módulos~(i),
instrumentação com \textit{stubs}~(ii), especificação de propriedades~(iii),
execução do verificador~(iv), análise de contraexemplos com conferência contra
o código de produção~(v), reporte~(vi), validação cruzada e verificação de
\textit{patch}~(vii), e reprodução em outro autopiloto com integração
contínua~(viii).
\item \textbf{Evidência experimental delimitada pelo modelo.} Dez módulos sem
violação até o limite de desenrolamento, com tempos por módulo medidos em três
execuções.
\item \textbf{Triagem documentada.} Três contraexemplos com três desfechos
distintos, e a prova formal de que o estado violador do GPS não ocorre no
driver de produção.
\item \textbf{Segunda campanha.} Dez módulos verificados sobre o código
original, com defeitos reproduzidos por \textit{replay}, um deles observado em
bancada, dois relatos ao PX4 e uma correção aceita no ESBMC.
\item \textbf{Artefato público.} \textit{Harnesses}, registros, programas de
\textit{replay}, capturas de bancada e roteiros de reprodução.
\end{enumerate}

A ausência de teste unitário direto para \texttt{dumpGpsData()} foi conferida
por busca na árvore do PX4 (Seção~\ref{sec:caso-gps}). Ela evidencia falta de
cobertura unitária direta, e não de exercício indireto por simulação ou por
testes de integração.

\section{Limitações do Estudo}
\label{sec:conclusao-limitacoes}

\textbf{Cobertura.} Os arquivos-alvo da primeira campanha somam cerca de
$1{,}5\,\%$ das $710$~mil linhas de C/C++ de \texttt{src/}. Ficaram de fora os
mecanismos de \textit{failsafe} do Commander, a lógica de \textit{geofence}, os
barramentos CAN/UAVCAN, a ponte RTPS/uXRCE-DDS e os drivers de fluxo óptico e
telêmetro, as mesmas exclusões da Seção~\ref{sec:ameacas-validade}.

\textbf{Limite de desenrolamento.} As garantias valem até o limite $k$
adotado: 10 como padrão, 15 no Navigator/Mission e 300 no MAVLink Parser
(Seção~\ref{subsec:escolha-k}). A campanha original usou
\texttt{--no-unwinding-assertions}; a reexecução com as asserções ligadas
afastou o risco de truncamento nos dez módulos sem violação, mas, no GPS, o
laço do \texttt{memcpy} excede $k = 10$ e não é coberto por inteiro.

\textbf{Modelo.} Os \textit{stubs} não representam temporização, atrasos de
barramento nem falhas intermitentes, e a análise modular não captura defeitos
que surgem da interação entre tarefas concorrentes.

%% LC (resolvido 12/11/2026): limitação sobre defeitos confirmados.
\textbf{Eficácia em encontrar defeitos.} A primeira campanha não produziu
defeito confirmado algum: de três contraexemplos, um era falso positivo, um foi
recusado e um é inalcançável na produção. A segunda campanha reproduziu
defeitos no código real, e três deles coincidem com correções feitas pelo
próprio PX4; ainda assim, nenhum relato deste trabalho havia sido aceito pelos
mantenedores no fechamento do texto. A eficácia do método em encontrar defeitos
que o projeto reconheça como tais está, portanto, demonstrada de forma
indireta.

%% LC (resolvido 12/11/2026): limitação sobre esforço humano.
\textbf{Esforço humano.} O tempo de máquina é pequeno, mas o custo do método é
dominado pelo trabalho de escrever \textit{harnesses}, triar contraexemplos e
conferir o código de produção. Esse esforço não foi medido, e nenhuma afirmação
sobre custo total de adoção se sustenta nos dados deste trabalho.

\textbf{Validação em bancada.} Os ensaios cobrem dois achados. No do RoboClaw,
o controlador de motores foi emulado, o \textit{firmware} precisou de duas
alterações declaradas e só um dos dois defeitos foi exercitado; no do GPS, não
havia receptor conectado.

%% LC (resolvido 12/11/2026): aplicabilidade ao SAGRES.
\textbf{Relação com o projeto SAGRES.} O projeto motivou a pesquisa, mas não
foi levantada a configuração de aeronave que ele utiliza, nem quais dos módulos
verificados fazem parte dela. A aplicabilidade dos resultados às suas missões é
uma hipótese, e não um resultado.

\section{Trabalhos Futuros}
\label{sec:conclusao-futuros}

\subsection{Expansão da cobertura}

%% LC (resolvido 12/11/2026): Navigator retirado da lista de não analisados;
%% LPE definido; citações mantidas só onde sustentam a frase.
A metodologia pode ser estendida aos subsistemas não analisados: os mecanismos
de \textit{failsafe} do Commander, a lógica de \textit{geofence}, os barramentos
CAN/UAVCAN e o estimador de posição local do PX4 (\textit{Local Position
Estimator}, LPE). Nos módulos já verificados, cabe ampliar os cenários; no
MAVLink Parser, por exemplo, sequências de mensagens malformadas do tipo
descrito por Ficco et al.~\cite{Ficco2022}.

\subsection{Validação em bancada e em simulação}
\label{sec:trabalhos-futuros-hw}

%% LC (resolvido 12/11/2026): seção numerada, para que as referências apontem
%% para ela.
A validação em bancada pode ser estendida aos demais achados da segunda
campanha. O sensor de vento FT7 e o batmon são os candidatos imediatos: o
primeiro usa a mesma interface serial do ensaio do RoboClaw, e o segundo
exigiria emular uma bateria SMBus. Para o RoboClaw, falta observar a reação de
um controlador real ao byte fora da faixa, que o manual não especifica. Para o
GPS, falta o sentido do receptor para a placa.

Os ensaios feitos injetam a condição de disparo diretamente, por parâmetro e
por comando de console. Eles mostram que o defeito ocorre na placa quando a
condição se apresenta, mas não que a cadeia de módulos a montante a produz em
operação. Responder a essa pergunta exige exercitar a pilha completa, e a
simulação \textit{software-in-the-loop} é o ambiente adequado, sem risco ao
equipamento.

\subsection{Integração em \textit{pipeline} de desenvolvimento}

Um fluxo do GitHub Actions já executa os \textit{harnesses} do IMU e do GPS a
cada \textit{push} (Seção~\ref{subsec:ardupilot-ci}). Os passos seguintes são a
verificação incremental, restrita aos módulos modificados, a medição do tempo
de relógio sob configuração controlada e a proposição do fluxo à integração
contínua oficial do projeto.

\subsection{Verificação composicional e concorrente}

A verificação módulo a módulo não detecta defeitos de composição. Dois recursos
do ESBMC atacam essa limitação. Os contratos de função permitem substituir um
\textit{stub} não determinístico por uma especificação verificada do que a
função exige e garante. A verificação de programas concorrentes, não explorada
aqui, permitiria tratar condições de corrida entre as tarefas que o PX4 executa
sob o escalonador do NuttX; exigiria modelar o escalonamento e a semântica de
publicação do uORB.

\subsection{Propriedades além do software}

As propriedades verificadas referem-se a estados do programa. O que uma
violação significa para o veículo, isto é, em que condição de voo ocorreria e
com que consequência, exige acoplar à verificação um modelo abstrato da
dinâmica do sistema físico. Essa foi uma pergunta da banca de qualificação que
este trabalho não responde.

\section{Considerações Finais}
\label{sec:conclusao-finais}

As três questões experimentais tiveram as seguintes respostas. Quanto à
\textbf{QE1}, todo módulo da primeira campanha foi verificado em menos de
noventa segundos em \textit{hardware} de prateleira, o que demonstra
viabilidade por módulo. Quanto à \textbf{QE2}, a resposta é parcial: o BMC
expôs, no \textit{harness} do GPS, uma violação em função sem teste unitário
direto, mas o estado é inalcançável na produção, e nenhum defeito do driver foi
exposto na primeira campanha. Quanto à \textbf{QE3}, nove dos doze módulos
foram verificados com a invocação padrão, e três exigiram opções específicas.

%% LC (resolvido 12/11/2026): alegações de custo e de casos de canto qualificadas.
Os resultados indicam que o BMC é aplicável a módulos de \textit{firmware} de
VANTs em produção com tempo de máquina pequeno. Não permitem afirmar que o
custo total seja compatível com um ciclo de desenvolvimento, porque o esforço
humano e o tempo em integração contínua não foram medidos. Também não se
sustenta, para a primeira campanha, a afirmação de que a abordagem alcança
casos extremos que os testes não cobrem: o único caso extremo encontrado ali
não ocorre no driver real.

%% LC (resolvido 12/11/2026): argumento final apoiado na segunda campanha.
O valor da abordagem aparece com mais clareza na segunda campanha, e por uma
razão metodológica. Quando a verificação foi feita sobre o código original, e
cada candidato só foi aceito depois de reproduzido fora do verificador, sete de
dez módulos exibiram defeitos que o código real de fato tem, três deles
corrigidos pelo próprio projeto sem conhecimento deste trabalho. A primeira
campanha mostrou o que acontece sem essa disciplina: contraexemplos formalmente
válidos para defeitos que não existem.

O que este trabalho entrega é, assim, uma metodologia de verificação formal
para \textit{firmware} de sistemas críticos, aplicada ao PX4 como estudo de
caso. Os achados são produto da aplicação do método, e não a sua medida:
aplicado a outro \textit{firmware}, ele poderia não revelar violação alguma e
ainda assim produzir evidência da mesma natureza, limitada, reproduzível e
auditável. A lição que atravessa as duas campanhas é que um contraexemplo é o
início da triagem, e que a conferência contra o código de produção é a etapa
que separa um achado de um artefato do modelo.



\bibliographystyle{cls_bst/bst_ufam}
\bibliography{referencias}

\appendix
\chapter[Produção Científica Derivada desta Dissertação]{\corr{Produção Científica Derivada desta Dissertação}}
\label{apend:producoes}

\artigo{O artigo científico derivado desta pesquisa foi aceito para publicação:}
\begin{itemize}
  \item \textbf{Título:} Formal Verification of Critical Modules in the PX4
        Autopilot using Bounded Model Checking
  \item \textbf{Autores:} Marcela Alves Rodrigues, Lucas Carvalho Cordeiro,
        José Reginaldo Hughes de Carvalho
  \item \textbf{Veículo:} XVI Simpósio Brasileiro de Engenharia de Sistemas
        Computacionais (SBESC~2026), trilha de artigo completo, Salvador/BA,
        24 a 27 de novembro de 2026
  \item \textbf{Status:} \artigo{Aceito (submissão \#33507); versão final
        enviada em outubro de 2026}
  \item \textbf{Artefato:} \textit{harnesses}, logs de execução e instruções de
        reprodução arquivados no Zenodo sob o DOI~\artigo{\texttt{10.5281/zenodo.23282236}
        (versão \texttt{v1.2-sbesc}; DOI de conceito \texttt{10.5281/zenodo.21401958})} e
        disponíveis em \url{https://github.com/marcela2881/dissertacao-px4-esbmc}
\end{itemize}

\chapter[Harness Completo --- GPS uBlox]{\corr{Harness Completo --- GPS uBlox}}
\label{apend:harness-gps}

\corr{Este apêndice reproduz na íntegra o \textit{harness} de
verificação do módulo GPS uBlox, \texttt{test\_gps\_driver.cpp}, utilizado
no estudo descrito no Capítulo~\ref{cap:metodologia}. Esse caso foi analisado
em maior profundidade nesta dissertação: dele provém o contraexemplo apresentado
na Seção~\ref{sec:contraexemplo-gps}.} \artigo{O mesmo caso também foi usado
na avaliação do \textit{patch} e na análise de alcançabilidade discutidas no
Capítulo~\ref{cap:experimentos}.}

\corr{A reprodução integral permite ao leitor examinar a implementação
utilizada na verificação, as pré-condições de cada teste e as diferenças entre
a adaptação incluída no \textit{harness} e o código do driver de produção. Essa
distinção é necessária para interpretar corretamente o contraexemplo: o
\textit{harness} preserva partes da lógica analisada, mas não reproduz
integralmente as estruturas nem todas as condições do driver.}

\corr{O arquivo define cinco funções de teste, das quais uma é selecionada
em cada execução por meio de uma variável não determinística em
\texttt{main()}. Cada teste estabelece pré-condições próprias para o tamanho
da entrada e para o estado inicial do buffer. Assim, os valores concretos do
contraexemplo dependem do teste selecionado e das escolhas feitas pelo solver.}

\artigo{O arquivo tem 294 linhas (Tabela~\ref{tab:modulos}) e é reproduzido sem
edição a partir do repositório de artefatos (Listagem~\ref{lst:harness-gps}). O contraexemplo da
Tabela~\ref{tab:resultados} vem do teste selecionado por
\texttt{test\_choice = 3}, \texttt{test\_gps\_real\_bit\_operation()}
(linhas 178--197).}

%% LC (resolvido 12/11/2026): as duas observacoes sobre o cabecalho e os
%% comentarios do harness viraram o paragrafo abaixo. O arquivo e reproduzido
%% literalmente, por isso os comentarios nao foram editados; a divergencia e
%% declarada no texto. Conferido contra gps.cpp e GpsDump.msg no commit 8be2590c15.
\corr{Como o arquivo é reproduzido sem edição, seus comentários internos
foram mantidos, embora contenham descrições que não correspondem integralmente
à versão do driver analisada. Para evitar ambiguidade, registram-se três
diferenças verificadas em relação ao commit \texttt{8be2590c15}. (i)~O
cabeçalho do \textit{harness} cita a versão~1.16 e a linha~643 de
\texttt{gps.cpp}; nesse commit, \texttt{GPS::dumpGpsData()} começa na
linha~727. (ii)~Os comandos indicados no rodapé do arquivo
(\texttt{--unwind 8} e \texttt{--unwind 5}) não são os usados na campanha
reportada nesta dissertação, que empregou
\texttt{--unwind 10 --no-unwinding-assertions --overflow-check}
(Seção~\ref{sec:config-experimental}). (iii)~A função apresentada é uma
adaptação, e não uma extração literal: o \textit{harness} usa um buffer de
200~bytes, enquanto a mensagem \texttt{GpsDump.msg} declara
\texttt{uint8[79] data}; além disso, o buffer e o modo ativo são passados como
parâmetros, e a estrutura modelada não inclui o campo
\texttt{device\_id} da mensagem. Portanto, a correspondência com o driver é
lógica e parcial. Em particular, a diferença de capacidade do buffer limita o
alcance das conclusões que podem ser extraídas do contraexemplo; essa questão é
retomada na Seção~\ref{subsec:gps-inducao}.}

\begin{lstlisting}[
  caption={\artigo{\textit{Harness} \texttt{test\_gps\_driver.cpp} (literal)}},
  label={lst:harness-gps},
  language=C++,
  numbers=left,
  frame=single,
  basicstyle=\ttfamily\scriptsize,
  breaklines=true,
  extendedchars=true,
  literate={á}{{\'a}}1 {à}{{`a}}1 {â}{{\^a}}1 {ã}{{\~a}}1 {é}{{\'e}}1 {ê}{{\^e}}1
    {í}{{\'i}}1 {ó}{{\'o}}1 {ô}{{\^o}}1 {õ}{{\~o}}1 {ú}{{\'u}}1 {ç}{{\c c}}1
    {Á}{{\'A}}1 {À}{{`A}}1 {Â}{{\^A}}1 {Ã}{{\~A}}1 {É}{{\'E}}1 {Ê}{{\^E}}1
    {Í}{{\'I}}1 {Ó}{{\'O}}1 {Ô}{{\^O}}1 {Õ}{{\~O}}1 {Ú}{{\'U}}1 {Ç}{{\c C}}1]
/**
 * @file teste_gps_driver_real_esbmc.cpp
 * @author Dissertação Mestrado - Verificação Formal PX4 v1.16
 * 
 * OBJETIVO: Teste de código REAL extraído do PX4 GPS Driver
 * FUNÇÃO TESTADA: GPS::dumpGpsData() - linha ~643 de src/drivers/gps/gps.cpp
 * MÉTODO: Bounded Model Checking com ESBMC
 * ESTRUTURAS: Baseadas no USO real no código gps.cpp
 */

#include <assert.h>
#include <cstring>
#include <cstdint>

// ================== FUNÇÕES ESBMC ==================
extern int nondet_int();
extern uint8_t nondet_uint8();
extern size_t nondet_size_t();
extern bool nondet_bool();
extern void __ESBMC_assume(int condition);

// ================== ESTRUTURAS REAIS EXTRAÍDAS DO PX4 ==================
/**
 * ESTRUTURA REAL baseada no USO no gps.cpp:
 * - dump_data->data + dump_data->len
 * - dump_data->len += write_len  
 * - sizeof(dump_data->data)
 * - dump_data->len |= 1 << 7
 * - dump_data->timestamp = hrt_absolute_time()
 * - dump_data->instance = (uint8_t)_instance
 */

// Tamanho realista que funciona com uint8_t len
#define GPS_DUMP_DATA_SIZE 200

struct gps_dump_s {
    uint8_t data[GPS_DUMP_DATA_SIZE];    // Buffer principal
    uint8_t len;                         // Índice atual (usado com bitwise ops)
    uint8_t instance;                    // Instância GPS
    uint64_t timestamp;                  // hrt_absolute_time()
};

enum class gps_dump_comm_mode_t : int32_t {
    Disabled = 0,
    Full = 1,
    RTCM = 2
};

// ================== FUNÇÃO REAL EXTRAÍDA DO PX4 ==================

/**
 * FUNÇÃO ORIGINAL dumpGpsData() extraída EXATAMENTE do gps.cpp linha ~643
 * Lógica preservada: while loop, memcpy, aritmética de ponteiros, bit operations
 */
void dumpGpsData(uint8_t *data, size_t len, gps_dump_comm_mode_t mode, bool msg_to_gps_device, 
                 gps_dump_s *dump_data, gps_dump_comm_mode_t active_mode)
{
    // Verificação de modo (do código original)
    if (active_mode != mode || !dump_data) {
        return;
    }

    dump_data->instance = 0; // Simular: dump_data->instance = (uint8_t)_instance;

    // LOOP CRÍTICO REAL DO PX4
    while (len > 0) {
        size_t write_len = len;

        // CÁLCULO CRÍTICO: potencial underflow se dump_data->len > GPS_DUMP_DATA_SIZE
        if (write_len > GPS_DUMP_DATA_SIZE - dump_data->len) {
            write_len = GPS_DUMP_DATA_SIZE - dump_data->len;
        }

        // OPERAÇÃO CRÍTICA: memcpy com aritmética de ponteiros (do código real)
        memcpy(dump_data->data + dump_data->len, data, write_len);
        
        // ATUALIZAÇÕES (exatamente como no gps.cpp)
        data += write_len;
        dump_data->len += write_len;
        len -= write_len;

        // LÓGICA DE PUBLICAÇÃO (do código original)
        if (dump_data->len >= GPS_DUMP_DATA_SIZE) {
            // BIT OPERATION do código real
            if (msg_to_gps_device) {
                dump_data->len |= 1 << 7;  // Set bit 7
            }

            // Simular: dump_data->timestamp = hrt_absolute_time();
            dump_data->timestamp = 12345;
            
            // Simular: _dump_communication_pub.publish(*dump_data);
            // Reset para próxima iteração (do código real)
            dump_data->len = 0;
        }
    }
}

// ================== TESTES DE VERIFICAÇÃO FORMAL ==================

/**
 * TESTE 1: Verificar buffer bounds na operação memcpy real
 * PROPRIEDADE: memcpy nunca deve escrever além do buffer data[]
 */
void test_gps_real_buffer_bounds() {
    size_t input_len = nondet_size_t();
    bool msg_to_device = nondet_bool();
    
    // Entrada realista mas potencialmente perigosa
    __ESBMC_assume(input_len > 0 && input_len <= 300);
    
    uint8_t input_data[300];
    gps_dump_s dump_buffer;
    dump_buffer.len = nondet_uint8();
    __ESBMC_assume(dump_buffer.len < GPS_DUMP_DATA_SIZE);
    
    uint8_t initial_len = dump_buffer.len;
    
    // Chamar função REAL do PX4
    dumpGpsData(input_data, input_len, gps_dump_comm_mode_t::Full, msg_to_device, 
                &dump_buffer, gps_dump_comm_mode_t::Full);
    
    // PROPRIEDADE CRÍTICA: len nunca deve exceder tamanho do buffer
    // (considerando que pode ser resetado para 0)
    assert(dump_buffer.len <= GPS_DUMP_DATA_SIZE);
}

/**
 * TESTE 2: Verificar underflow no cálculo GPS_DUMP_DATA_SIZE - dump_data->len
 * PROPRIEDADE: Subtração deve ser segura mesmo se len estiver corrompido
 */
void test_gps_real_underflow_protection() {
    size_t input_len = nondet_size_t();
    
    __ESBMC_assume(input_len > 0 && input_len <= 50);
    
    uint8_t input_data[50];
    gps_dump_s dump_buffer;
    
    // CENÁRIO CRÍTICO: len próximo ou maior que GPS_DUMP_DATA_SIZE
    dump_buffer.len = nondet_uint8();
    __ESBMC_assume(dump_buffer.len >= GPS_DUMP_DATA_SIZE - 10);
    
    // Chamar função REAL
    dumpGpsData(input_data, input_len, gps_dump_comm_mode_t::Full, false, 
                &dump_buffer, gps_dump_comm_mode_t::Full);
    
    // PROPRIEDADE: Resultado deve ser coerente
    assert(dump_buffer.len <= GPS_DUMP_DATA_SIZE);
}

/**
 * TESTE 3: Verificar terminação do loop while(len > 0) real
 * PROPRIEDADE: Loop deve sempre terminar, mesmo com inputs extremos
 */
void test_gps_real_loop_termination() {
    size_t input_len = nondet_size_t();
    
    __ESBMC_assume(input_len > 0 && input_len <= 100);
    
    uint8_t input_data[100];
    gps_dump_s dump_buffer;
    dump_buffer.len = nondet_uint8();
    __ESBMC_assume(dump_buffer.len < GPS_DUMP_DATA_SIZE);
    
    // Chamar função REAL
    dumpGpsData(input_data, input_len, gps_dump_comm_mode_t::Full, false, 
                &dump_buffer, gps_dump_comm_mode_t::Full);
    
    // PROPRIEDADE: Se chegou aqui, loop terminou
    assert(true);
}

/**
 * TESTE 4: Verificar bit operation real do código
 * PROPRIEDADE: dump_data->len |= 1 << 7 deve ser segura
 */
void test_gps_real_bit_operation() {
    size_t input_len = nondet_size_t();
    
    __ESBMC_assume(input_len > 0 && input_len <= GPS_DUMP_DATA_SIZE + 10);
    
    uint8_t input_data[GPS_DUMP_DATA_SIZE + 10];
    gps_dump_s dump_buffer;
    dump_buffer.len = nondet_uint8();
    __ESBMC_assume(dump_buffer.len >= GPS_DUMP_DATA_SIZE - 5);
    
    uint8_t initial_len = dump_buffer.len;
    
    // Chamar função REAL que pode executar bit operation
    dumpGpsData(input_data, input_len, gps_dump_comm_mode_t::Full, true, 
                &dump_buffer, gps_dump_comm_mode_t::Full);
    
    // PROPRIEDADE: Operação bit deve ser válida OU len deve ser resetado
    assert(dump_buffer.len == 0 || dump_buffer.len <= GPS_DUMP_DATA_SIZE || 
           (dump_buffer.len & 0x7F) <= GPS_DUMP_DATA_SIZE);
}

/**
 * TESTE 5: Verificar edge case quando buffer está completamente cheio
 * PROPRIEDADE: Função deve lidar corretamente com buffer no limite
 */
void test_gps_real_full_buffer_edge_case() {
    size_t input_len = nondet_size_t();
    
    __ESBMC_assume(input_len > 0 && input_len <= 20);
    
    uint8_t input_data[20];
    gps_dump_s dump_buffer;
    
    // Buffer exatamente cheio
    dump_buffer.len = GPS_DUMP_DATA_SIZE;
    
    // Chamar função REAL
    dumpGpsData(input_data, input_len, gps_dump_comm_mode_t::Full, false, 
                &dump_buffer, gps_dump_comm_mode_t::Full);
    
    // PROPRIEDADE: Deve lidar graciosamente com buffer cheio
    assert(dump_buffer.len <= GPS_DUMP_DATA_SIZE);
}

// ================== MAIN PARA ESBMC ==================
int main() {
    int test_choice = nondet_int();
    __ESBMC_assume(test_choice >= 0 && test_choice < 5);
    
    switch(test_choice) {
        case 0:
            test_gps_real_buffer_bounds();
            break;
        case 1:
            test_gps_real_underflow_protection();
            break;
        case 2:
            test_gps_real_loop_termination();
            break;
        case 3:
            test_gps_real_bit_operation();
            break;
        case 4:
            test_gps_real_full_buffer_edge_case();
            break;
    }
    
    return 0;
}

/* 
 * ================================================================
 * DOCUMENTAÇÃO
 * ================================================================
 * 
 * METODOLOGIA DEMONSTRADA:
 * 
 * 1. EXTRAÇÃO DE CÓDIGO 100% REAL:
 *    - Função dumpGpsData() copiada exatamente de src/drivers/gps/gps.cpp
 *    - Estrutura gps_dump_s baseada no USO real observado no código
 *    - Lógica preservada: while loop, memcpy, bit operations, reset
 * 
 * 2. ESTRUTURA DERIVADA DO CÓDIGO REAL:
 *    - data[]: deduzido de sizeof(dump_data->data) 
 *    - len: deduzido de dump_data->len += write_len e operações bit
 *    - instance: observado em dump_data->instance = (uint8_t)_instance
 *    - timestamp: observado em dump_data->timestamp = hrt_absolute_time()
 * 
 * 3. PROPRIEDADES BASEADAS EM VULNERABILIDADES REAIS:
 *    - Buffer overflow em memcpy + aritmética de ponteiros
 *    - Integer underflow em GPS_DUMP_DATA_SIZE - dump_data->len
 *    - Loop infinito se write_len não decrementar len corretamente
 *    - Bit operation safety em len |= 1 << 7
 * 
 * 4. TÉCNICA DE VERIFICAÇÃO:
 *    - Bounded Model Checking com ESBMC
 *    - Entrada simbólica para tamanhos e estados do buffer
 *    - Assumptions para cenários críticos específicos
 *    - Assertions para propriedades de memory safety
 * 
 * COMANDOS DE EXECUÇÃO:
 * esbmc teste_gps_driver_real_esbmc.cpp --unwind 8 --timeout 300s
 * esbmc teste_gps_driver_real_esbmc.cpp --overflow-check --unwind 5
 * 
 * VULNERABILIDADES ALVO:
 * - Buffer overflow se dump_data->len corrompido
 * - Integer underflow na subtração GPS_DUMP_DATA_SIZE - len
 * - Aritmética de ponteiros insegura em memcpy
 * - Loop infinito se write_len == 0
 * 
 * DIFERENCIAL DESTA VERSÃO:
 * - Estruturas baseadas em USO real, não especulação
 * - Função extraída palavra por palavra do gps.cpp
 * - Tamanhos realistas compatíveis com uint8_t
 * - Testes focados em padrões observados no código real
 * 
 * ================================================================
 */
\end{lstlisting}

\corr{A leitura do código também exige duas ressalvas. Primeiro, o teste
\texttt{test\_gps\_real\_loop\_termination()} termina com
\texttt{assert(true)}; essa asserção é tautológica e, isoladamente, não
constitui uma propriedade formal de terminação. Segundo, as asserções dos
outros testes verificam condições ao final da execução, mas não substituem a
verificação direta de todas as operações de memória realizadas durante o
caminho. O contraexemplo reproduzido a seguir é relevante porque o ESBMC acusa
especificamente uma desreferência fora dos limites no \texttt{memcpy}. Essas
limitações são consideradas na interpretação dos resultados e impedem que o
\textit{harness} seja descrito como uma cópia integral ou uma validação completa
do driver de produção.}

\chapter{Contraexemplos Gerados pelo ESBMC}
\label{apend:contraexemplos}

%% Prof. Lucas (qualificacao): "Poderias mostrar o contraexemplo gerado pelo ESBMC
%% no apendice." -- ATENDIDO: apendice adicionado.

Este apêndice apresenta contraexemplos gerados pelo ESBMC nos módulos IMU
BMI088 e GPS uBlox. Conforme discutido no Capítulo~\ref{cap:experimentos},
apenas o contraexemplo do GPS uBlox permaneceu como violação alcançável no
\textit{harness} correspondente. Isso não significa, por si só, que o estado
seja alcançável no driver de produção: a análise de alcançabilidade descrita na
Seção~\ref{subsec:gps-inducao} concluiu que ele não é alcançável nesse código.
Os traços do BMI088 são apresentados para documentar a etapa de triagem; um foi
classificado como falso positivo e o outro como especificação disputada. \corr{Os contraexemplos da segunda campanha (RoboClaw,
batmon e candidatos adicionais) não são reproduzidos aqui; os registros
completos estão no repositório de artefatos, nas pastas indicadas na
Seção~\ref{sec:segunda-campanha}.}

\corr{Cada traço inclui os valores concretos atribuídos pelo solver
SMT, o rastro de execução até o estado de falha e a propriedade violada reportada
pelo verificador.} \artigo{Os traços abaixo são trechos literais dos logs do
ESBMC~v8.1.0 com Boolector~3.2.4 depositados no repositório de artefatos (os
mesmos que fundamentam a Tabela~\ref{tab:resultados} e as Listagens do artigo);
a única edição é a elisão de vetores de bytes longos, indicada no próprio
texto. Os contraexemplos foram originalmente reportados} \corr{à comunidade PX4 nas
Issues~\#26865 (IMU BMI088)\footnote{\url{https://github.com/PX4/PX4-Autopilot/issues/26865}}
e~\#26866 (GPS uBlox)\footnote{\url{https://github.com/PX4/PX4-Autopilot/issues/26866}}
\artigo{a partir de execuções anteriores com o ESBMC~v7.10.0; os tempos e os
valores aqui reportados vêm das execuções com o ESBMC~v8.1.0, conforme a nota
de rodapé do artigo.}}

\corr{Cabe uma observação metodológica sobre a leitura destes traços.
As entradas dos \textit{harnesses} são não-determinísticas, de modo que o solver
é livre para escolher qualquer testemunha que satisfaça a condição de violação.
Execuções distintas do mesmo \textit{harness}, sob versões ou solvers diferentes,
produzem valores concretos diferentes para a mesma propriedade violada, sem que
isso constitua divergência de resultado.} \artigo{Os valores reproduzidos
abaixo são os das execuções com o ESBMC~v8.1.0, idênticos nas três execuções
de cada módulo.}

\section{Contraexemplo 1 --- Temperatura fora da faixa operacional (BMI088)}
\label{sec:contraexemplo-bmi088-temp}

O traço da Listagem~\ref{lst:contraex_bmi088_temp} viola a propriedade
$\phi_{\text{temp}}$ (Equação~\ref{eq:prop_temp}): a temperatura calculada é
$86{,}25\,^{\circ}$C, acima do limite superior de
$85\,^{\circ}$C indicado para a faixa operacional no \textit{datasheet} do
BMI088~\cite{Bosch2021}. \corr{Pela conversão indicada pelo fabricante,
$T = \texttt{Temp\_int11} \times 0{,}125 + 23$, o valor
\texttt{Temp\_uint11 = 506} resulta em
$506 \times 0{,}125 + 23 = 86{,}25\,^{\circ}$C. O resultado demonstra que
a entrada considerada pelo \textit{harness} ultrapassa o limite definido; sua
interpretação como defeito do driver depende da validade dessa entrada no
contexto de operação do sensor.}

\begin{lstlisting}[
  caption={\corr{Contraexemplo ESBMC para a violação da faixa de temperatura no BMI088} \artigo{(trecho literal do log, ESBMC~v8.1.0)}},
  label={lst:contraex_bmi088_temp},
  numbers=left,
  frame=single,
  language=bash,
  basicstyle=\ttfamily\scriptsize,
  breaklines=true]
ESBMC version 8.1.0 64-bit x86_64 linux
Solving with solver Boolector 3.2.4
[Counterexample]

State 2 file resultados/IMU_BMI088/test_bmi088_imu.cpp line 120 column 5 function test_temperature_calculation thread 0
----------------------------------------------------
  temp_msb = 63 (00111111)

State 3 file resultados/IMU_BMI088/test_bmi088_imu.cpp line 121 column 5 function test_temperature_calculation thread 0
----------------------------------------------------
  temp_lsb = 95 (01011111)

State 7 file resultados/IMU_BMI088/test_bmi088_imu.cpp line 42 column 5 function updateTemperature thread 0
----------------------------------------------------
  Temp_uint11 = 506 (00000001 11111010)

State 8 file resultados/IMU_BMI088/test_bmi088_imu.cpp line 48 column 9 function updateTemperature thread 0
----------------------------------------------------
  Temp_int11 = 506 (00000001 11111010)

State 11 file resultados/IMU_BMI088/test_bmi088_imu.cpp line 51 column 5 function updateTemperature thread 0
----------------------------------------------------
  temperature = 8.625000e+1f (01000010 10101100 10000000 00000000)

State 12 file resultados/IMU_BMI088/test_bmi088_imu.cpp line 123 column 5 function test_temperature_calculation thread 0
----------------------------------------------------
  temperature = 8.625000e+1f (01000010 10101100 10000000 00000000)

State 13 file resultados/IMU_BMI088/test_bmi088_imu.cpp line 126 column 5 function test_temperature_calculation thread 0
----------------------------------------------------
Violated property:
  file resultados/IMU_BMI088/test_bmi088_imu.cpp line 126 column 5 function test_temperature_calculation
  assertion temperature >= -40.0f && temperature <= 85.0f
  temperature >= -40.000000f && temperature <= 85.000000f

VERIFICATION FAILED
\end{lstlisting}

\corr{O resultado foi inicialmente associado à validação da faixa de entrada
e relacionado às categorias CWE-1284 (\textit{Improper Validation of Specified
Quantity in Input}) e CWE-20 (\textit{Improper Input Validation}). Essa
associação deve ser entendida como classificação do cenário investigado, não
como confirmação de uma vulnerabilidade. Conforme discutido na
Seção~\ref{subsec:imu-temperatura}, a observação foi reportada à comunidade
PX4, mas a abordagem de saturação proposta no PR~\#27575 não foi aceita pelos
mantenedores. Assim, o caso permanece como especificação disputada, sem defeito
de segurança confirmado nem remediação aceita.}

\section{Contraexemplo 2 --- Contador FIFO excede capacidade do buffer (BMI088)}
\label{sec:contraexemplo-bmi088-fifo}

\corr{O traço da Listagem~\ref{lst:contraex_bmi088_fifo} apresenta
\texttt{count = 9088} bytes, valor superior ao
limite de \texttt{FIFO\_SIZE = 1.024} bytes e, portanto, uma violação da
propriedade verificada no \textit{harness}. Na conferência com o driver de
produção não instrumentado (Seção~\ref{subsec:imu-fifo}), porém, o candidato foi
classificado como falso positivo: o código de produção contém três níveis de
proteção, e esse valor aciona a primeira guarda de reinicialização antes do
acesso ao buffer. O resultado do ESBMC decorre da ausência dessas guardas no
\textit{harness}; não demonstra um acesso fora dos limites alcançável no driver
de produção.}

\artigo{O traço é o das três execuções da Tabela~\ref{tab:resultados} para o
IMU BMI088 (\texttt{test\_choice = 2}): é deste candidato, e não do de
temperatura, o tempo de 0,1\,s registrado na tabela.}

\begin{lstlisting}[
  caption={\artigo{Contraexemplo ESBMC para o contador do FIFO no BMI088 (trecho literal do log, ESBMC~v8.1.0)}},
  label={lst:contraex_bmi088_fifo},
  numbers=left,
  frame=single,
  language=bash,
  basicstyle=\ttfamily\scriptsize,
  breaklines=true]
ESBMC version 8.1.0 64-bit x86_64 linux
Solving with solver Boolector 3.2.4
[Counterexample]

State 2 file resultados/IMU_BMI088/test_bmi088_imu.cpp line 245 column 5 function main thread 0
----------------------------------------------------
  test_choice = 2 (00000000 00000000 00000000 00000010)

State 4 file resultados/IMU_BMI088/test_bmi088_imu.cpp line 138 column 5 function test_fifo_count_calculation thread 0
----------------------------------------------------
  fifo_len_0 = 128 (10000000)

State 5 file resultados/IMU_BMI088/test_bmi088_imu.cpp line 139 column 5 function test_fifo_count_calculation thread 0
----------------------------------------------------
  fifo_len_1 = 227 (11100011)

State 6 file resultados/IMU_BMI088/test_bmi088_imu.cpp line 61 column 5 function fifoReadCount thread 0
----------------------------------------------------
  FIFO_LENGTH_1_MASKED = 35 (00100011)

State 8 file resultados/IMU_BMI088/test_bmi088_imu.cpp line 141 column 5 function test_fifo_count_calculation thread 0
----------------------------------------------------
  count = 9088 (00100011 10000000)

State 9 file resultados/IMU_BMI088/test_bmi088_imu.cpp line 144 column 5 function test_fifo_count_calculation thread 0
----------------------------------------------------
Violated property:
  file resultados/IMU_BMI088/test_bmi088_imu.cpp line 144 column 5 function test_fifo_count_calculation
  assertion count <= FIFO_SIZE
  (unsigned long int)count <= FIFO_SIZE

VERIFICATION FAILED
\end{lstlisting}

\section[Contraexemplo 3 --- Underflow e acesso fora dos limites (GPS uBlox)]{\corr{Contraexemplo 3 --- Underflow e acesso fora dos limites (GPS uBlox)}}
\label{sec:contraexemplo-gps}

O traço viola a propriedade $\phi_{\text{gps}}$
(Equação~\ref{eq:prop_gps}) no \textit{harness}. A verificação de espaço
modelada em \texttt{dumpGpsData()} calcula a capacidade restante por meio da
subtração entre a capacidade do buffer e o nível de preenchimento:

\begin{lstlisting}[language=C,basicstyle=\ttfamily\scriptsize,frame=none,numbers=none]
if (write_len > GPS_DUMP_DATA_SIZE - dump_data->len)
\end{lstlisting}

No estado encontrado pelo solver, \texttt{dump\_data->len = 207}, enquanto
\texttt{GPS\_DUMP\_DATA\_SIZE = 200}. A subtração é promovida para
\texttt{int} e resulta em $-7$; na comparação com \texttt{write\_len}, que é
do tipo \texttt{size\_t}, o valor negativo é convertido para um inteiro sem
sinal muito grande. Com isso, a condição da guarda não limita
\texttt{write\_len}, e o \texttt{memcpy} tenta escrever a partir da posição
207 do buffer modelado, ultrapassando seus 200 bytes. Esse comportamento é
observado no \textit{harness}; a análise de alcançabilidade no driver de
produção é tratada separadamente na Seção~\ref{subsec:gps-inducao}.

\corr{O ESBMC reporta uma única propriedade violada
(Listagem~\ref{lst:contraex_gps}): a falha de
\artigo{desreferência} no \texttt{memcpy}. O \textit{underflow} aritmético é a causa raiz,
e não uma segunda violação independente --- a ferramenta encerra na primeira
violação encontrada. A classificação CWE-191 (\textit{Integer Underflow (Wrap or Wraparound)})
$\rightarrow$~CWE-119 (\textit{Improper Restriction of Operations within the
Bounds of a Memory Buffer}) reflete essa
relação de causa e efeito.}

\begin{lstlisting}[
  caption={\corr{Contraexemplo ESBMC para o underflow com acesso fora dos limites no GPS uBlox} \artigo{(trecho literal do log, ESBMC~v8.1.0)}},
  label={lst:contraex_gps},
  numbers=left,
  frame=single,
  language=bash,
  basicstyle=\ttfamily\scriptsize,
  breaklines=true]
ESBMC version 8.1.0 64-bit x86_64 linux
Solving with solver Boolector 3.2.4
[Counterexample]

State 1 file resultados/GPS/test_gps_driver.cpp line 224 column 5 function main thread 0
----------------------------------------------------
  test_choice = 3 (00000000 00000000 00000000 00000011)

State 3 file resultados/GPS/test_gps_driver.cpp line 179 column 5 function test_gps_real_bit_operation thread 0
----------------------------------------------------
  input_len = 10 (00000000 00000000 00000000 00000000 00000000 00000000 00000000 00001010)

State 6 file resultados/GPS/test_gps_driver.cpp line 183 column 5 function test_gps_real_bit_operation thread 0
----------------------------------------------------
  input_data = { ... 210 bytes omitidos ... }

State 7 file resultados/GPS/test_gps_driver.cpp line 184 column 5 function test_gps_real_bit_operation thread 0
----------------------------------------------------
  dump_buffer = { .data={ ... 200 bytes omitidos ... }, .len=0,
    .instance=0, .anon_pad$3=281474976710655, .timestamp=0xFFFFFFFFFFFFFFFF }

State 8 file resultados/GPS/test_gps_driver.cpp line 185 column 23 function test_gps_real_bit_operation thread 0
----------------------------------------------------
  dump_buffer.len = { .data={ ... 200 bytes omitidos ... }, .len=207,
    .instance=0, .anon_pad$3=281474976710655, .timestamp=0xFFFFFFFFFFFFFFFF }

State 21 file <esbmc>/c2goto/library/string.c line 280 column 5 function __memcpy_impl thread 0
----------------------------------------------------
Violated property:
  file <esbmc>/c2goto/library/string.c line 280 column 5 function __memcpy_impl
  dereference failure: Access to object out of bounds

VERIFICATION FAILED
\end{lstlisting}

\corr{O contraexemplo foi obtido com o solver Boolector~3.2.4. Conforme
registrado na Seção~\ref{sec:ameacas-validade}, essa foi a única configuração
que encontrou a violação dentro do limite de tempo estabelecido. Esse resultado
indica sensibilidade à configuração do verificador e do solver; não significa
que a violação seja exclusiva desse solver.} \artigo{O estado
\texttt{len = 207} é alcançável no modelo do \textit{harness}, mas não no
driver de produção, conforme a análise indutiva apresentada na
Seção~\ref{subsec:gps-inducao}.} \corr{A variável local
\texttt{dump\_buffer} representa, no modelo, a estrutura usada como destino
do \texttt{memcpy}; ela não deve ser interpretada como reprodução exata da
estrutura publicada pelo driver.}

\corr{Embora uma escrita fora dos limites possa, em termos gerais, corromper
memória adjacente, este experimento não permite afirmar que tal efeito ocorra em
uma plataforma Pixhawk. O \textit{harness} utiliza uma estrutura adaptada e um
buffer de 200~bytes, e a análise de alcançabilidade não encontrou o estado
\texttt{len = 207} no driver de produção. Portanto, o resultado deve ser
interpretado como um contraexemplo do modelo analisado, e não como evidência de
corrupção de memória em execução real.}


\end{document}
